% Kinematics of particles moving in a straight line — Pure Mathematics with Mechanics Upper Sixth — Cours signé OnBuch+
% Official MINESEC syllabus. Compiler avec : tectonic PMM09-kinematics-of-particles-moving-in-a-straight-line.tex
\newcommand{\DOCMATIERE}{Pure Mathematics with Mechanics}
\newcommand{\DOCNIVEAU}{Upper Sixth}
\newcommand{\DOCMODULE}{Upper Sixth — Pure mathematics with mechanics}
\newcommand{\DOCLECON}{9}
\newcommand{\DOCTITRE}{Kinematics of particles moving in a straight line}
% OnBuch+ — préambule « cours signés » ENRICHI (compilé avec tectonic / XeLaTeX)
% Système visuel « Pop » d'OnBuch+ : fond crème (jamais blanc), bordures encre
% épaisses, ombres « dures » décalées, titres Archivo Black en capitales.
% Version MATHÉMATIQUES : graphes (pgfplots/tikz), boîtes pédagogiques
% (définition, propriété, méthode, attention, le savais-tu, exercice résolu,
% exercice, TP, bilan), page de garde et signature OnBuch+ ancrée en bas.
%
% Macros attendues dans main.tex AVANT \input{preamble} :
%   \DOCMATIERE  \DOCNIVEAU  \DOCMODULE  \DOCLECON (numéro)  \DOCTITRE
\documentclass[11pt]{article}

\usepackage{fontspec}
\usepackage[english]{babel}
\usepackage[a4paper,margin=2cm,top=3.4cm,bottom=2.8cm,headheight=1.4cm,footskip=1.3cm]{geometry}
\usepackage{amsmath,amssymb,amsfonts}
\usepackage{mathtools} % \xmapsto, \coloneqq... souvent utilisés par les modèles
\usepackage{cancel} % simplifications barrées (bilans de forces)
% \abs{z} (module, valeur absolue) et \norm{u} (norme), utilisés par les modèles
\providecommand{\abs}{}\let\abs\relax\DeclarePairedDelimiter{\abs}{\lvert}{\rvert}
\providecommand{\norm}{}\let\norm\relax\DeclarePairedDelimiter{\norm}{\lVert}{\rVert}
\usepackage[table]{xcolor} % [table] : fournit aussi \rowcolors (alternance de lignes)
\usepackage{pagecolor}
\usepackage{tikz}
\usetikzlibrary{arrows.meta,positioning,calc,shapes.geometric,shapes.misc,shapes.arrows,decorations.pathmorphing,decorations.pathreplacing,decorations.markings,patterns,shadows,mindmap,trees,fit,backgrounds,matrix,intersections,angles,quotes}
\usepackage{pgfplots}
\usepackage[version=4]{mhchem} % équations nucléaires \ce{^{235}_{92}U}
\usepackage[european,siunitx]{circuitikz} % schémas électriques
\pgfplotsset{compat=1.18}
\usepgfplotslibrary{fillbetween,groupplots} % aires entre courbes (calcul intégral)
\usepackage{enumitem}
\usepackage{fancyhdr}
% [most] charge toutes les bibliothèques tcolorbox (listings, minted, raster,
% documentation...) et gonfle inutilement la mémoire TeX sur un document long
% (~300 boîtes) : on ne charge que ce qui est réellement utilisé.
\usepackage[skins,breakable]{tcolorbox}
\usepackage{graphicx}
\usepackage{colortbl}
\usepackage{booktabs}
\usepackage{tabularx}
\usepackage{diagbox} % case d'en-tête diagonale des tableaux à double entrée
% Colonnes extensibles centrée / gauche / droite (C, L, R), souvent utilisées par les modèles
\newcolumntype{C}{>{\centering\arraybackslash}X}
\newcolumntype{L}{>{\raggedright\arraybackslash}X}
\newcolumntype{R}{>{\raggedleft\arraybackslash}X}
\usepackage{array}
\usepackage{multicol}
\usepackage{siunitx}
% parse-numbers=false : accepte aussi \SI{2,5 \times 10^{-3}}{...}
\sisetup{locale=UK,output-decimal-marker={.},group-minimum-digits=4,parse-numbers=false}
\DeclareSIUnit\ohm{\ensuremath{\symup{\Omega}}} % U+2126 absent de la police
\DeclareSIUnit\division{div} % réglages d'oscilloscope (ms/div, V/div)
\DeclareSIUnit\tr{tr} % tours (vitesse de rotation, ex. \SI{1500}{\tr\per\minute})
\usepackage{qrcode}
% \degree : commande fréquemment utilisée par les modèles pour un angle ou une
% température en dehors de \SI{}{} (non fournie par les paquets chargés).
\providecommand{\degree}{^{\circ}}
% \qed (fin de démonstration), utilisé par les modèles sans amsthm
\providecommand{\qed}{\hfill$\square$}
% \barre : double barre des valeurs interdites dans un tableau de variations (tkz-tab)
\providecommand{\barre}{\ensuremath{\|}}
% environnement proof (amsthm non chargé) utilisé par les modèles
\ifdefined\proof\else\newenvironment{proof}[1][Proof]{\par\noindent\textit{#1.}\ }{\qed\par}\fi
\usepackage{lastpage}
\usepackage{hyperref}
\hypersetup{hidelinks}

% ============================================================
% Palette « Pop » OnBuch+ — mêmes hex que la classe P (plus_ui.dart)
% ============================================================
\definecolor{poporange}{HTML}{F59321}
\definecolor{popdark}{HTML}{E07A0C}
\definecolor{popcream}{HTML}{FFF6E8}
\definecolor{popink}{HTML}{1C1714}
\definecolor{poppurple}{HTML}{7A5AE0}
\definecolor{popgreen}{HTML}{2E9E6A}
\definecolor{poppink}{HTML}{E0457B}
\definecolor{popblue}{HTML}{2F7DE1}
\definecolor{popgold}{HTML}{C98A12}
\definecolor{popmuted}{HTML}{8A7F76}
\definecolor{popline}{HTML}{EADFD2}
\definecolor{popgray}{HTML}{8A7F76}
\colorlet{popgrey}{popgray}
% Alias tolérants (le modèle nomme parfois les couleurs différemment)
\colorlet{popwhite}{white}
\colorlet{popblack}{popink}
\colorlet{popred}{poppink}
\colorlet{popyellow}{popgold}
% Teintes claires pour fonds de tableaux / zones
\colorlet{popblueL}{popblue!10!white}
\colorlet{poporangeL}{poporange!14!white}
\colorlet{popgreenL}{popgreen!12!white}
\colorlet{poppurpleL}{poppurple!12!white}
\colorlet{poppinkL}{poppink!12!white}
\colorlet{popgoldL}{popgold!14!white}

\pagecolor{popcream}

% ============================================================
% Typographie
% ============================================================
\defaultfontfeatures{Ligatures=TeX}
\setmainfont{PlusJakartaSans-Medium.ttf}[Path=../../../fonts/,BoldFont=PlusJakartaSans-Bold.ttf]
% Maths en sans-serif (Fira Math) avec chiffres et lettres droites de Plus
% Jakarta, pour que les formules se fondent dans le texte.
\usepackage[math-style=ISO,bold-style=ISO]{unicode-math}
\setmathfont{FiraMath-Regular.otf}[Path=../../../fonts/]
\setmathfont{PlusJakartaSans-Medium.ttf}[Path=../../../fonts/,range={up/num,up/{Latin,latin},\mathup}]
% (icomma retiré : décimales avec point en anglais)
% Caractères mathématiques Unicode absents des polices de texte (Plus Jakarta,
% Archivo Black) : rendus par leur équivalent mathématique, en texte comme en
% formule — sinon ils s'affichent en carré vide (ex. « Arithmétique dans ℤ »).
\usepackage{newunicodechar}
\newunicodechar{ℝ}{\ensuremath{\mathbb{R}}}
\newunicodechar{ℤ}{\ensuremath{\mathbb{Z}}}
\newunicodechar{ℂ}{\ensuremath{\mathbb{C}}}
\newunicodechar{ℕ}{\ensuremath{\mathbb{N}}}
\newunicodechar{ℚ}{\ensuremath{\mathbb{Q}}}
\newunicodechar{𝔻}{\ensuremath{\mathbb{D}}}
\newunicodechar{α}{\ensuremath{\alpha}}
\newunicodechar{β}{\ensuremath{\beta}}
\newunicodechar{γ}{\ensuremath{\gamma}}
\newunicodechar{δ}{\ensuremath{\delta}}
\newunicodechar{ε}{\ensuremath{\varepsilon}}
\newunicodechar{θ}{\ensuremath{\theta}}
\newunicodechar{λ}{\ensuremath{\lambda}}
\newunicodechar{μ}{\ensuremath{\mu}}
\newunicodechar{σ}{\ensuremath{\sigma}}
\newunicodechar{τ}{\ensuremath{\tau}}
\newunicodechar{φ}{\ensuremath{\varphi}}
\newunicodechar{ψ}{\ensuremath{\psi}}
\newunicodechar{ω}{\ensuremath{\omega}}
\newunicodechar{Σ}{\ensuremath{\Sigma}}
% majuscules produites par \MakeUppercase dans les titres (α-aminés -> Α)
\newunicodechar{Α}{\ensuremath{\alpha}}
\newunicodechar{Β}{\ensuremath{\beta}}
\newunicodechar{Γ}{\ensuremath{\Gamma}}
\newunicodechar{Δ}{\ensuremath{\Delta}}
\newunicodechar{Ε}{\ensuremath{\varepsilon}}
\newunicodechar{Θ}{\ensuremath{\Theta}}
\newunicodechar{Λ}{\ensuremath{\Lambda}}
\newunicodechar{Μ}{\ensuremath{\mu}}
\newunicodechar{Τ}{\ensuremath{\tau}}
\newunicodechar{Φ}{\ensuremath{\Phi}}
\newunicodechar{Ψ}{\ensuremath{\Psi}}
\newunicodechar{Ω}{\ensuremath{\Omega}}
\newunicodechar{↦}{\ensuremath{\mapsto}}
\newunicodechar{∈}{\ensuremath{\in}}
\newunicodechar{∉}{\ensuremath{\notin}}
\newunicodechar{∧}{\ensuremath{\wedge}}
\newunicodechar{⇔}{\ensuremath{\Leftrightarrow}}
\newunicodechar{⟺}{\ensuremath{\Longleftrightarrow}}
\newunicodechar{⇒}{\ensuremath{\Rightarrow}}
\newunicodechar{⟹}{\ensuremath{\Longrightarrow}}
\newunicodechar{∘}{\ensuremath{\circ}}
\newunicodechar{∪}{\ensuremath{\cup}}
\newunicodechar{∩}{\ensuremath{\cap}}
\newunicodechar{≡}{\ensuremath{\equiv}}
\newunicodechar{⊂}{\ensuremath{\subset}}
\newunicodechar{⊥}{\ensuremath{\perp}}
\newunicodechar{∃}{\ensuremath{\exists}}
\newunicodechar{∀}{\ensuremath{\forall}}
\newunicodechar{∛}{\ensuremath{\sqrt[3]{\,}}}
\newunicodechar{∜}{\ensuremath{\sqrt[4]{\,}}}
\newunicodechar{⃗}{\ensuremath{{}^{\to}}}
\newunicodechar{ⁿ}{\ifmmode^{n}\else\textsuperscript{\ensuremath{n}}\fi}
\newunicodechar{ˣ}{\ifmmode^{x}\else\textsuperscript{\ensuremath{x}}\fi}
\newunicodechar{ᵇ}{\ifmmode^{b}\else\textsuperscript{\ensuremath{b}}\fi}
\newunicodechar{ʳ}{\ifmmode^{r}\else\textsuperscript{\ensuremath{r}}\fi}
\newunicodechar{ᵘ}{\ifmmode^{u}\else\textsuperscript{\ensuremath{u}}\fi}
\newunicodechar{ʸ}{\ifmmode^{y}\else\textsuperscript{\ensuremath{y}}\fi}
\newunicodechar{ᵃ}{\ifmmode^{a}\else\textsuperscript{\ensuremath{a}}\fi}
\newunicodechar{ᵉ}{\ifmmode^{e}\else\textsuperscript{\ensuremath{e}}\fi}
\newunicodechar{ᵏ}{\ifmmode^{k}\else\textsuperscript{\ensuremath{k}}\fi}
\newunicodechar{ᵖ}{\ifmmode^{p}\else\textsuperscript{\ensuremath{p}}\fi}
\newunicodechar{ᴺ}{\ifmmode^{N}\else\textsuperscript{\ensuremath{N}}\fi}
\newunicodechar{ᶜ}{\ifmmode^{c}\else\textsuperscript{\ensuremath{c}}\fi}
\newunicodechar{ʰ}{\ifmmode^{h}\else\textsuperscript{\ensuremath{h}}\fi}
\newunicodechar{ᵠ}{\ifmmode^{q}\else\textsuperscript{\ensuremath{q}}\fi}
\newunicodechar{ₙ}{\ifmmode_{n}\else\textsubscript{\ensuremath{n}}\fi}
\newunicodechar{ₐ}{\ifmmode_{a}\else\textsubscript{\ensuremath{a}}\fi}
\newunicodechar{ᵢ}{\ifmmode_{i}\else\textsubscript{\ensuremath{i}}\fi}
\newunicodechar{ᵣ}{\ifmmode_{r}\else\textsubscript{\ensuremath{r}}\fi}
\newunicodechar{ₖ}{\ifmmode_{k}\else\textsubscript{\ensuremath{k}}\fi}
\newunicodechar{ᵦ}{\ifmmode_{\beta}\else\textsubscript{\ensuremath{\beta}}\fi}
\newunicodechar{ₓ}{\ifmmode_{x}\else\textsubscript{\ensuremath{x}}\fi}
\newfontfamily\popdisplay{ArchivoBlack-Regular.ttf}[Path=../../../fonts/]
\newfontfamily\poplogo{SpaceGrotesk-Bold.ttf}[Path=../../../fonts/]
\newfontfamily\popsemi{PlusJakartaSans-SemiBold.ttf}[Path=../../../fonts/]
\newfontfamily\popxbold{PlusJakartaSans-ExtraBold.ttf}[Path=../../../fonts/]
\newfontfamily\popxboldls{PlusJakartaSans-ExtraBold.ttf}[Path=../../../fonts/,LetterSpace=3.0]

\setlist[enumerate]{leftmargin=1.7em,itemsep=5pt,topsep=4pt}
\setlist[itemize]{leftmargin=1.7em,itemsep=4pt,topsep=4pt,label={\color{poporange}\textbullet}}
\setlength{\parindent}{0pt}
\setlength{\parskip}{5pt}
\renewcommand{\arraystretch}{1.25}

% pgfplots : style Pop par défaut
\pgfplotsset{
  every axis/.append style={axis line style={popink,line width=1pt},tick style={popink},
    grid style={popline,line width=0.5pt},label style={font=\small},tick label style={font=\footnotesize}},
  popcurve/.style={poporange,line width=1.8pt,no markers,smooth},
}
% Même style disponible dans un \draw TikZ simple (hors pgfplots)
\tikzset{popcurve/.style={poporange,line width=1.8pt}}
\tikzset{popgrid/.style={popline,very thin}}
\colorlet{popcurve}{poporange} % le nom du style est parfois utilisé comme couleur

% ============================================================
% Pill / squiggle
% ============================================================
\newcommand{\poppill}[3]{%
  \tikz[baseline=-0.55ex]\node[draw=popink,line width=1.2pt,rounded corners=2.6pt,%
    fill=#2,inner xsep=6.5pt,inner ysep=3pt]{%
    \popxboldls\fontsize{7}{7}\selectfont\color{#3}\MakeUppercase{#1}};%
}
\newcommand{\popsquiggle}[1]{%
  \tikz[baseline=-0.4ex]{
    \draw[#1,line width=2.6pt,line cap=round]
      plot [smooth] coordinates {(0,0.09) (0.34,0.01) (0.68,0.21) (1.02,0.01) (1.36,0.10)};
  }%
}
% Mot-clé surligné (marqueur orange sous le texte)
\newcommand{\cle}[1]{{\popxbold\color{popdark}#1}}

% ============================================================
% En-tête / pied
% ============================================================
\pagestyle{fancy}
\fancyhf{}
\renewcommand{\headrulewidth}{0pt}
\renewcommand{\footrulewidth}{0pt}
\lhead{\raisebox{-0.15ex}{\poplogo\fontsize{13}{13}\selectfont\color{popink}OnBuch\color{poporange}+}}
\rhead{\poppill{\DOCMATIERE\ \textbullet\ \DOCNIVEAU\ \textbullet\ Chapter \DOCLECON}{white}{popink}}
\lfoot{\popxbold\fontsize{7.6}{7.6}\selectfont\color{popmuted}\MakeUppercase{Signed course OnBuch+}\ \textbullet\ {\popsemi\DOCTITRE}}
\rfoot{\tikz[baseline=-0.6ex]\node[rounded corners=8.5pt,fill=popink,minimum height=17pt,minimum width=17pt,inner xsep=5pt,inner sep=0]{\popxbold\fontsize{8}{8}\selectfont\color{popcream}\thepage\,/\,\pageref*{LastPage}};}

% ============================================================
% Titres
% ============================================================
\newcommand{\chaptitre}[2]{%
  \begin{tcolorbox}[enhanced,colback=poporange,colframe=popink,boxrule=2.6pt,arc=11pt,
      drop shadow=popink,left=15pt,right=15pt,top=12pt,bottom=11pt,before skip=3pt,after skip=18pt]
    \poppill{#2}{popink}{popcream}\par\vspace{9pt}
    {\raggedright\hyphenpenalty=10000\exhyphenpenalty=10000\popdisplay\fontsize{22}{24}\selectfont\color{popcream}\MakeUppercase{#1}\par}
  \end{tcolorbox}
}
% Section numérotée : pastille orange avec le numéro + titre Archivo
\newcounter{popsec}
\newcommand{\coursec}[1]{%
  \stepcounter{popsec}\setcounter{popsub}{0}%
  \par\vspace{16pt}\noindent
  \begin{minipage}{\linewidth}
  \tikz[baseline=-0.7ex]\node[draw=popink,line width=1.6pt,fill=poporange,rounded corners=4pt,minimum size=22pt,inner sep=0]{\popdisplay\fontsize{12}{12}\selectfont\color{popcream}\thepopsec};%
  \hspace{8pt}{\popdisplay\fontsize{15}{17}\selectfont\color{popink}\MakeUppercase{#1}}\par
  \vspace{4pt}\noindent{\color{poporange}\rule{36pt}{3.6pt}}
  \end{minipage}\par\nopagebreak\vspace{8pt}
}
\newcounter{popsub}
\newcommand{\courssub}[1]{%
  \stepcounter{popsub}%
  \par\vspace{9pt}\noindent{\popxbold\fontsize{11.5}{13}\selectfont\color{popink}{\color{poporange}\thepopsec.\thepopsub}\hspace{6pt}#1}\par\nopagebreak\vspace{4pt}
}

% ============================================================
% Boîtes pédagogiques — même recette : fond blanc, bordure épaisse colorée,
% ombre dure encre, badge plein. Titre optionnel [#1].
% ============================================================
\tcbset{popbase/.style={breakable,enhanced,colback=white,boxrule=2.3pt,arc=9pt,
  shadow={3pt}{-3pt}{0pt}{popink},left=14pt,right=14pt,top=11pt,bottom=11pt,before skip=11pt,after skip=15pt,
  fontupper=\popsemi\fontsize{10.6}{14.5}\selectfont\color{popink},
  fontlower=\popsemi\fontsize{10.6}{14.5}\selectfont\color{popink}}}
% Coche dessinée (le glyphe ✓ n'existe pas dans les polices du document)
\newcommand{\popcheck}[1]{\tikz[baseline=-0.5ex]\draw[#1,line width=1.6pt,line cap=round,line join=round](-0.13,0.0)--(-0.04,-0.09)--(0.13,0.1);}
\providecommand{\coche}{\popcheck{popgreen}}
\providecommand{\resultat}[1]{\par\smallskip\noindent\fcolorbox{popgreen}{popgreenL}{\parbox{\dimexpr\linewidth-2\fboxsep-2\fboxrule\relax}{\textbf{Result.} #1}}\par\smallskip}
\newcommand{\popboxhead}[3]{\poppill{#1}{#2}{white}\ifx\relax#3\relax\else\hspace{6pt}{\popxbold\fontsize{10.6}{12}\selectfont\color{popink}#3}\fi\par\vspace{7pt}}

\newtcolorbox{aretenir}[1][]{popbase,colframe=popgreen,before upper={\popboxhead{Key points}{popgreen}{#1}}}
\newtcolorbox{exemplebox}[1][]{popbase,colframe=poporange,before upper={\popboxhead{Exemple}{poporange}{#1}}}
\newtcolorbox{definition}[1][]{popbase,colframe=popblue,before upper={\popboxhead{Definition}{popblue}{#1}}}
\newtcolorbox{propriete}[1][]{popbase,colframe=poppurple,before upper={\popboxhead{Property}{poppurple}{#1}}}
\newtcolorbox{methode}[1][]{popbase,colframe=popgold,before upper={\popboxhead{Method}{popgold}{#1}}}
\newtcolorbox{attention}[1][]{popbase,colframe=poppink,before upper={\popboxhead{Warning}{poppink}{#1}}}
\newtcolorbox{experience}[1][]{popbase,colframe=popblue,colback=popblueL,before upper={\popboxhead{Experiment}{popblue}{#1}}}
% « Le savais-tu ? » : carte violette pleine (contexte, culture, Cameroun)
\newtcolorbox{savaistu}[1][]{popbase,colback=poppurple,colframe=popink,
  fontupper=\popsemi\fontsize{10.4}{14.2}\selectfont\color{white},
  before upper={\poppill{Did you know?}{white}{poppurple}\ifx\relax#1\relax\else\hspace{6pt}{\popxbold\color{white}#1}\fi\par\vspace{7pt}}}
% Exercice résolu : énoncé en haut, \tcblower puis la solution sur fond crème
\newcounter{popexo}\newcounter{popexr}
\newtcolorbox{exoresolu}[1][]{popbase,colframe=poporange,colbacklower=popcream,
  segmentation style={popink,line width=1.2pt,dashed},
  before upper={\stepcounter{popexr}\popboxhead{Worked example \thepopexr}{poporange}{#1}},
  before lower={\poppill{Solution}{popink}{popcream}\par\vspace{6pt}}}
% Exercice (non résolu en place) — la correction est dans la partie Corrigés
\newtcolorbox{exercice}[1][]{popbase,colframe=popink,
  before upper={\stepcounter{popexo}\popboxhead{Exercise \thepopexo}{popink}{#1}}}
% Corrigé
\newtcolorbox{corrige}[1][]{popbase,colframe=popgreen,colback=popgreenL,
  before upper={\popboxhead{Solution}{popgreen}{#1}}}
% Objectifs / prérequis / bilan
\newtcolorbox{objectifs}{popbase,colframe=popink,colback=poporangeL,
  before upper={\popboxhead{Chapter objectives}{popink}{}}}
\newtcolorbox{prerequis}{popbase,colframe=popink,colback=popblueL,
  before upper={\popboxhead{Prerequisites}{popblue}{}}}
\newtcolorbox{bilan}{popbase,colframe=popink,colback=white,boxrule=2.8pt,
  before upper={\popboxhead{Fiche bilan}{poporange}{L'essentiel à connaître pour l'examen}}}
% Figure encadrée avec légende
\newtcolorbox{popfigure}[1][]{enhanced,breakable=false,colback=white,colframe=popline,boxrule=1.4pt,arc=8pt,
  left=10pt,right=10pt,top=10pt,bottom=8pt,before skip=10pt,after skip=12pt,halign=center,
  lower separated=false,halign lower=center,
  fontlower=\popsemi\fontsize{9}{11.5}\selectfont\color{popmuted},
  #1}
\newcommand{\legende}[1]{\par\vspace{6pt}{\popsemi\fontsize{9}{11.5}\selectfont\color{popmuted}#1}}

% ============================================================
% Signature OnBuch+
% ============================================================
\newcommand{\poplogoinline}[1]{{\poplogo\fontsize{#1}{#1}\selectfont\color{popink}OnBuch\color{poporange}+}}
\newcommand{\signatureonbuch}{%
  \begin{tcolorbox}[enhanced,colback=popink,colframe=popink,boxrule=2.2pt,arc=10pt,
      drop shadow=poporange,left=14pt,right=14pt,top=10pt,bottom=10pt,before skip=0pt,after skip=0pt]
    \tikz[baseline=-0.7ex]\node[circle,fill=popgreen,draw=popcream,line width=1.3pt,minimum size=22pt,inner sep=0]{\popcheck{white}};%
    \hspace{9pt}\begin{minipage}[c]{0.62\linewidth}
      {\popxbold\fontsize{12}{13}\selectfont\color{popcream}Signed\ \textbullet\ {\poplogo OnBuch\color{poporange}+}}\par\vspace{2pt}
      {\raggedright\popsemi\fontsize{8}{10.5}\selectfont\color{popline}\MakeUppercase{OnBuch+ teaching team}\\\MakeUppercase{Follows the official MINESEC syllabus}\par}
    \end{minipage}\hfill
    {\poplogo\fontsize{22}{22}\selectfont\color{popcream}OnBuch\color{poporange}+}
  \end{tcolorbox}%
}

% ============================================================
% Page de garde
% #1 titre · #2 matière · #3 niveau · #4 module · #5 n° leçon · #6 durée ·
% #7 description / objectifs (texte ou itemize)
% ============================================================
\newcommand{\pagegarde}[7]{%
  \thispagestyle{empty}
  \newgeometry{a4paper,margin=1.7cm,top=1.6cm,bottom=1.4cm}%
  \begin{tikzpicture}[remember picture,overlay]
    \fill[poporange!18] ([xshift=-3.2cm,yshift=-3.6cm]current page.north east) circle (4.6cm);
    \fill[poppurple!14] ([xshift=2.4cm,yshift=7.6cm]current page.south west) circle (3.8cm);
  \end{tikzpicture}%
  {\poplogo\fontsize{30}{30}\selectfont\color{popink}OnBuch\color{poporange}+}\hfill
  \raisebox{0.6ex}{\poppill{Signed course \textbullet\ Premium}{popink}{popcream}}\par
  \vspace{1pt}\popsquiggle{poppurple}\par
  \vspace{8pt}
  \begin{tcolorbox}[enhanced,colback=poporange,colframe=popink,boxrule=2.8pt,arc=13pt,
      drop shadow=popink,left=18pt,right=18pt,top=12pt,bottom=13pt,after skip=10pt]
    \poppill{#2 \textbullet\ #3}{popink}{popcream}\hspace{5pt}\poppill{#4}{white}{popink}\par\vspace{8pt}
    {\popdisplay\fontsize{14}{14}\selectfont\color{popink}CHAPTER #5}\par\vspace{4pt}
    {\raggedright\hyphenpenalty=10000\exhyphenpenalty=10000\popdisplay\fontsize{25}{27}\selectfont\color{popcream}\MakeUppercase{#1}\par}
    \vspace{8pt}
    {\popxbold\fontsize{9.5}{9.5}\selectfont\color{popink}\MakeUppercase{Suggested time: #6}}
  \end{tcolorbox}
  \begin{tcolorbox}[enhanced,colback=white,colframe=popink,boxrule=2.2pt,arc=10pt,
      drop shadow=popink,left=16pt,right=16pt,top=10pt,bottom=10pt,after skip=8pt]
    \poppill{In this chapter}{poppurple}{white}\par\vspace{5pt}
    \popsemi\fontsize{9.3}{12.6}\selectfont\color{popink}#7
  \end{tcolorbox}
  \noindent\begin{minipage}[t]{0.487\linewidth}
    \begin{tcolorbox}[enhanced,colback=popgreen,colframe=popink,boxrule=2.2pt,arc=10pt,
        drop shadow=popink,left=11pt,right=11pt,top=10pt,bottom=10pt,height=3.1cm,valign=center]
      \tcbox[colback=white,colframe=popink,boxrule=1.3pt,arc=5pt,boxsep=0pt,nobeforeafter,
        left=3pt,right=3pt,top=3pt,bottom=3pt,valign=center]{\qrcode[height=1.9cm]{https://play.google.com/store/apps/details?id=com.luvvix.onbuch&hl=en}}%
      \hspace{8pt}\begin{minipage}[c]{0.5\linewidth}
        {\popdisplay\fontsize{11}{12.5}\selectfont\color{white}\MakeUppercase{Download OnBuch}}\par\vspace{4pt}
        {\raggedright\popsemi\fontsize{8.4}{11}\selectfont\color{white}Free \textbullet\ Google Play\\AI tutor, past papers, results\par}
      \end{minipage}
    \end{tcolorbox}
  \end{minipage}\hfill
  \begin{minipage}[t]{0.487\linewidth}
    \begin{tcolorbox}[enhanced,colback=poppurple,colframe=popink,boxrule=2.2pt,arc=10pt,
        drop shadow=popink,left=11pt,right=11pt,top=10pt,bottom=10pt,height=3.1cm,valign=center]
      \tikz[baseline=-0.7ex]\node[circle,fill=white,draw=popink,line width=1.3pt,minimum size=1.6cm,inner sep=0]{\popdisplay\fontsize{24}{24}\selectfont\color{poppurple}+};%
      \hspace{8pt}\begin{minipage}[c]{0.6\linewidth}
        {\popdisplay\fontsize{11}{12.5}\selectfont\color{white}\MakeUppercase{Go OnBuch+}}\par\vspace{4pt}
        {\raggedright\popsemi\fontsize{8.4}{11}\selectfont\color{white}All signed courses, unlimited AI tutor, PDF sheets — OnBuch+ tab in the app\par}
      \end{minipage}
    \end{tcolorbox}
  \end{minipage}
  \vspace{8pt}
  \popcontenu
  \vfill
  \signatureonbuch
  \restoregeometry
  \newpage
}

% Rangée « Ce cours contient » (page de garde)
\newcommand{\poptile}[3]{% #1 couleur #2 gros texte #3 légende
  \begin{tcolorbox}[enhanced,colback=#1,colframe=popink,boxrule=2pt,arc=9pt,shadow={2.5pt}{-2.5pt}{0pt}{popink},
      left=6pt,right=6pt,top=9pt,bottom=9pt,halign=center,valign=center,height=1.75cm,width=\linewidth,nobeforeafter]
    {\popdisplay\fontsize{12.5}{13}\selectfont\color{white}\MakeUppercase{#2}}\par\vspace{3pt}
    {\centering\popsemi\fontsize{7.8}{9.5}\selectfont\color{white}#3\par}
  \end{tcolorbox}}
\newcommand{\popcontenu}{%
  \par\noindent{\popxboldls\fontsize{7.5}{7.5}\selectfont\color{popmuted}THIS SIGNED COURSE CONTAINS}\par\vspace{7pt}
  \noindent\begin{minipage}[t]{0.235\linewidth}\poptile{popblue}{Course}{complete, illustrated, step by step}\end{minipage}\hfill
  \begin{minipage}[t]{0.235\linewidth}\poptile{poporange}{Methods}{and worked examples}\end{minipage}\hfill
  \begin{minipage}[t]{0.235\linewidth}\poptile{poppink}{Exercises}{exam style + solutions}\end{minipage}\hfill
  \begin{minipage}[t]{0.235\linewidth}\poptile{popgreen}{Summary sheet}{the essentials to remember}\end{minipage}\par}

% Sommaire
\newtcolorbox{sommaire}{enhanced,colback=white,colframe=popink,boxrule=2.2pt,arc=10pt,
  drop shadow=popink,left=18pt,right=18pt,top=13pt,bottom=13pt,before skip=6pt,after skip=20pt,
  before upper={\poppill{Contents}{poppurple}{white}\par\vspace{9pt}\popsemi\fontsize{10.6}{14.5}\selectfont\color{popink}}}

% Signature de fin de cours — poussée tout en bas de la dernière page
\newcommand{\findecours}{%
  \par\vfill\nopagebreak
  \begin{center}\popsquiggle{poporange}\end{center}\vspace{4pt}
  \signatureonbuch
}
\AtBeginDocument{\def\llbracket{\mathopen{[\mkern-3.5mu[}}\def\rrbracket{\mathclose{]\mkern-3.5mu]}}} % intervalles d'entiers (glyphe ⟦ absent de la police)
% Glyphes absents de Fira Math : redéfinis à partir de glyphes disponibles
\AtBeginDocument{%
  \def\setminus{\mathbin{\backslash}}\let\smallsetminus\setminus
  \def\vdots{\mathord{\vbox{\baselineskip3.2pt\lineskiplimit0pt\kern4pt\hbox{.}\hbox{.}\hbox{.}}}}%
  \def\ddots{\mathinner{\mkern1mu\raise6.5pt\hbox{.}\mkern2mu\raise3.6pt\hbox{.}\mkern2mu\raise0.7pt\hbox{.}\mkern1mu}}%
  \def\triangle{\mathord{\scalebox{0.9}{$\symup{\Delta}$}}}%
  \def\nabla{\mathord{\raisebox{\depth}{\scalebox{1}[-1]{$\symup{\Delta}$}}}}%
  \def\checkmark{\popcheck{popdark}}%
  \def\star{\mathbin{\ast}}\def\bigstar{\mathord{\ast}}%
  \def\diamond{\mathbin{\scalebox{0.6}{$\Diamond$}}}%
  \def\bigcup{\mathop{\mathchoice{\raisebox{-0.3em}{\scalebox{1.7}{$\cup$}}}{\scalebox{1.25}{$\cup$}}{\cup}{\cup}}\displaylimits}%
  \def\bigcap{\mathop{\mathchoice{\raisebox{-0.3em}{\scalebox{1.7}{$\cap$}}}{\scalebox{1.25}{$\cap$}}{\cap}{\cap}}\displaylimits}%
  \def\lesseqqgtr{\mathrel{\lessgtr}}\def\bigoplus{\mathop{\oplus}\displaylimits}%
}
\providecommand{\ssi}{if and only if} % abréviation fréquente
\AtBeginDocument{\providecommand{\wideparen}{\overparen}} % arc de cercle

%% FIGFIX-BEGIN v4 — correctifs globaux des figures (docs/onbuch-generation/tools/figfix)
\usepackage{adjustbox}\usepackage{etoolbox}
% toute figure TikZ/pgfplots plus large que la ligne est réduite à la largeur disponible
\BeforeBeginEnvironment{tikzpicture}{\begin{adjustbox}{max width=\linewidth}}
\AfterEndEnvironment{tikzpicture}{\end{adjustbox}}
% légendes pgfplots lisibles par-dessus les courbes
\pgfplotsset{every axis/.append style={legend style={fill=white,fill opacity=0.92,text opacity=1,draw=black!15,inner sep=3pt}}}
% étiquettes d'arêtes (syntaxe edge["..."]) avec fond blanc
\tikzset{every edge quotes/.append style={fill=white,inner sep=1.2pt}}
%% FIGFIX-END
\begin{document}

\pagegarde{Kinematics of particles moving in a straight line}{Pure Mathematics with Mechanics}{Upper Sixth}{Upper Sixth — Pure mathematics with mechanics}{9}{10 periods}{This chapter introduces the mathematical description of motion in one dimension: displacement, velocity, acceleration and their graphical representations. Learners master the equations of uniformly accelerated motion and apply them to free fall and real-life situations, in full preparation for the GCE Advanced Level mechanics paper.
\vspace{4pt}
\begin{multicols}{2}\small
\begin{itemize}
\item By the end of this chapter I can distinguish between distance and displacement, speed and velocity, and average and instantaneous quantities.
\item By the end of this chapter I can state and use the five equations of motion under uniform acceleration (suvat equations).
\item By the end of this chapter I can select the appropriate suvat equation for a given problem and solve it systematically.
\item By the end of this chapter I can plot, read and interpret displacement–time, velocity–time and acceleration–time graphs.
\item By the end of this chapter I can deduce velocity from the gradient of an x–t graph and displacement from the area under a v–t graph.
\item By the end of this chapter I can connect the three types of motion graphs and sketch any one of them from another.
\end{itemize}\end{multicols}}

\begin{sommaire}
\begin{multicols}{2}
\begin{enumerate}
\item Section 1: Concept of Motion in One Dimension (Lesson 92)
\item Section 2: Equations of Motion under Uniform Acceleration (Lesson 93)
\item Section 3: Displacement–Time and Velocity–Time Graphs (Lessons 94, 95, 96)
\item Section 4: Acceleration–Time Graphs and Connecting the Three Graphs (Lessons 97, 98)
\item Section 5: Free Fall and Motion under Gravity (Lessons 99, 100)
\item Section 6: Real-Life Applications and Synthesis (Lessons 94 \& 101, consolidation)
\item Integration activity
\item Methods and worked examples
\item Exercises
\item Solutions to the exercises
\item Summary sheet
\end{enumerate}
\end{multicols}
\end{sommaire}

\begin{objectifs}
\begin{itemize}
\item By the end of this chapter I can distinguish between distance and displacement, speed and velocity, and average and instantaneous quantities.
\item By the end of this chapter I can state and use the five equations of motion under uniform acceleration (suvat equations).
\item By the end of this chapter I can select the appropriate suvat equation for a given problem and solve it systematically.
\item By the end of this chapter I can plot, read and interpret displacement–time, velocity–time and acceleration–time graphs.
\item By the end of this chapter I can deduce velocity from the gradient of an x–t graph and displacement from the area under a v–t graph.
\item By the end of this chapter I can connect the three types of motion graphs and sketch any one of them from another.
\item By the end of this chapter I can model free fall as motion under uniform acceleration g ≈ 9.8 m/s² and solve vertical motion problems.
\item By the end of this chapter I can apply kinematics to real-life situations such as braking distances, road safety and falling objects.
\item By the end of this chapter I can solve GCE-style problems involving two moving particles, meeting points and overtaking.
\end{itemize}
\end{objectifs}

\begin{prerequis}
\begin{itemize}
\item Algebraic manipulation: solving linear and quadratic equations, rearranging formulae (Lower Sixth algebra).
\item Gradient of a straight line and area of basic shapes (rectangle, triangle, trapezium).
\item Plotting and reading Cartesian graphs; equation of a line y = mx + c.
\item Basic SI units (metre, second) and unit conversions (km/h ↔ m/s).
\item Simple differentiation of polynomials is helpful but NOT required (flagged as enrichment only).
\end{itemize}
\end{prerequis}

\newpage

\coursec{Section 1: Concept of Motion in One Dimension (Lesson 92)}

Every morning, the \emph{bend-skin} motorbike rider in Bamenda must judge in a split second whether he can stop before the junction when the traffic light turns orange. A driver on the Yaound\'e--Douala highway who follows too closely at $100\ \mathrm{km\,h^{-1}}$ may never stop in time, while a mango falling from a tree in Buea accelerates in exactly the same way as a stone dropped from a storey building. How far does the bike travel before stopping? How long does the mango take to hit the ground? These everyday questions are all answered by one powerful tool: the \cle{kinematics} of a particle moving in a straight line.

In this first section we build the vocabulary of kinematics: what a particle is, how we describe where it is (position, displacement), how fast it moves (speed, velocity) and how its motion changes (acceleration). Everything in the rest of the chapter rests on these ideas, so we shall define them carefully.

\courssub{The Particle Model}

Look again at the bend-skin. It has a length, a width, wheels that spin, a rider who leans into corners. If we tried to describe the motion of every part of it, the problem would be impossibly complicated. But if all we want to know is \emph{how far along the road} the bike is at each instant, then its size and shape do not matter: we can pretend the whole bike, rider included, is a single point carrying the total mass.

\begin{definition}[The particle model]
A \cle{particle} is a body whose dimensions are negligible compared with the distances involved in its motion. We model it as a single point at which the whole mass of the body is concentrated.
\end{definition}

\begin{attention}[The particle model is an idealisation]
No real object is a point. The model is valid when the size of the body is small compared with the distances travelled and when we do not care about rotation or internal motion. A bend-skin over a $500\ \mathrm{m}$ stretch of road is a particle; the same bend-skin parking in a $2\ \mathrm{m}$ gap is not. Always state the modelling assumptions you use in an examination answer.
\end{attention}

\begin{savaistu}[Why physicists love simple models]
The particle model is the first of a family of idealisations you will meet in mechanics: the \emph{light inextensible string}, the \emph{smooth pulley}, the \emph{smooth plane}. Each one removes a complication (mass of the string, friction, \dots) so that the essential mathematics appears clearly. Mastering mechanics is largely learning which idealisation fits which situation.
\end{savaistu}

\courssub{Position and Displacement}

To describe motion along a straight line we turn the line into a number line: we choose an \cle{origin} $O$ (a reference point) and a \cle{positive direction}. The \cle{position} of a particle $P$ at time $t$ is then a single number $x$, its directed distance from $O$, measured in metres ($\mathrm{m}$).

\begin{definition}[Displacement]
The \cle{displacement} of a particle between two instants is the change of position
\[
s = \Delta x = x_{\text{final}} - x_{\text{initial}} .
\]
Displacement is a \textbf{vector} quantity: it has a magnitude (in metres) and a direction (given by its sign along the chosen axis). A negative displacement simply means ``in the negative direction''.
\end{definition}

The \cle{distance} travelled, by contrast, is the total length of the path actually covered. It is a \textbf{scalar}: always positive, with no direction attached.

\begin{propriete}[Distance versus displacement]
For any motion,
\[
\text{distance travelled} \;\geq\; |\text{displacement}|,
\]
with equality \textbf{if and only if} the particle moves in one fixed direction throughout (it never reverses). Whenever the particle turns back, the distance travelled strictly exceeds the magnitude of the displacement.
\end{propriete}

The reason is clear: displacement only ``sees'' the start and the finish, while distance accumulates every metre covered, including any backtracking. A round trip back to the start gives \emph{zero} displacement but a large distance.

\begin{exemplebox}[Distance and displacement on a number line]
A particle starts at $x = -2\ \mathrm{m}$, moves to $x = 5\ \mathrm{m}$, then returns to $x = 1\ \mathrm{m}$.
\begin{itemize}
\item Displacement: $\Delta x = 1 - (-2) = +3\ \mathrm{m}$.
\item Distance travelled: $|-2 \to 5| + |5 \to 1| = 7 + 4 = 11\ \mathrm{m}$.
\end{itemize}
Note that $11 > |3|$, exactly as the property predicts, because the particle reversed at $x = 5$.
\end{exemplebox}

\begin{popfigure}
\begin{center}
\begin{tikzpicture}[scale=0.95]
\draw[->, thick, popmuted] (-3.2,0) -- (6.2,0) node[below left] {$x\ (\mathrm{m})$};
\foreach \xx in {-3,-2,-1,0,1,2,3,4,5,6}
  \draw (\xx,0.08) -- (\xx,-0.08) node[below] {\small $\xx$};
\fill[popblue] (-2,0) circle (2.2pt) node[above=2pt, popblue] {\small start};
\fill[poporange] (5,0) circle (2.2pt) node[above=2pt, poporange] {\small turn};
\fill[popgreen] (1,0) circle (2.2pt) node[above=2pt, popgreen] {\small end};
\draw[->, very thick, popblue] (-2,0.55) -- (5,0.55) node[midway, above, popblue] {\small $7$ m};
\draw[->, very thick, poporange] (5,1.05) -- (1,1.05) node[midway, above, poporange] {\small $4$ m};
\draw[->, very thick, popgreen] (-2,-0.55) -- (1,-0.55) node[midway, below, popgreen] {\small displacement $= +3$ m};
\end{tikzpicture}
\end{center}
\legende{A particle goes from $x=-2$ m to $x=5$ m, then back to $x=1$ m. Distance travelled $= 11$ m; displacement $= +3$ m.}
\end{popfigure}

\courssub{Average Speed and Average Velocity}

Knowing where a particle is at two instants, the natural next question is: \emph{how fast did it get there?}

\begin{definition}[Average speed and average velocity]
\[
\text{average speed} = \frac{\text{distance travelled}}{\text{time taken}}, \qquad
\text{average velocity} = \frac{\text{displacement}}{\text{time taken}} = \frac{\Delta x}{\Delta t}.
\]
Average speed is a \textbf{scalar}; average velocity is a \textbf{vector} (its sign gives the direction). The SI unit of both is the metre per second, $\mathrm{m\,s^{-1}}$; in everyday life we also use kilometres per hour, $\mathrm{km\,h^{-1}}$.
\end{definition}

\begin{propriete}[Conversion of units]
\[
1\ \mathrm{m\,s^{-1}} = 3.6\ \mathrm{km\,h^{-1}}, \qquad 1\ \mathrm{km\,h^{-1}} = \frac{1}{3.6}\ \mathrm{m\,s^{-1}}.
\]
\emph{Derivation (not examinable but instructive):} $1\ \mathrm{m\,s^{-1}}$ means $1$ m in $1$ s, i.e. $3600$ m in $3600$ s, i.e. $3.6$ km in $1$ hour.
\end{propriete}

\begin{attention}[Mixing units]
The most frequent source of wrong answers in kinematics is mixing $\mathrm{km\,h^{-1}}$ with $\mathrm{m\,s^{-1}}$, or hours with seconds. \textbf{Before any calculation, convert everything to SI units} (metres, seconds, $\mathrm{m\,s^{-1}}$). For example, $90\ \mathrm{km/h} = 90/3.6 = 25\ \mathrm{m/s}$.
\end{attention}

\begin{exoresolu}[Average speed and average velocity]
A trader walks from her stall, taken as origin, $120$ m east to the market in $100$ s, then walks $40$ m back west in $40$ s to meet a customer. Take east as positive. Find (a) the total distance, (b) the displacement, (c) her average speed, (d) her average velocity for the whole journey.
\tcblower
\textbf{Solution.}
\begin{itemize}
\item (a) Distance $= 120 + 40 = 160\ \mathrm{m}$.
\item (b) Displacement $\Delta x = 120 - 40 = +80\ \mathrm{m}$ (i.e. $80$ m east of the stall).
\item (c) Total time $= 100 + 40 = 140\ \mathrm{s}$. Average speed $= \dfrac{160}{140} \approx 1.14\ \mathrm{m\,s^{-1}}$.
\item (d) Average velocity $= \dfrac{\Delta x}{\Delta t} = \dfrac{80}{140} \approx +0.57\ \mathrm{m\,s^{-1}}$, directed east.
\end{itemize}
Notice the average speed is about twice the magnitude of the average velocity, because of the backtracking.
\end{exoresolu}

\courssub{Instantaneous Velocity}

Average velocity over a whole journey can hide a lot: the bend-skin may average $30\ \mathrm{km/h}$ yet reach $60\ \mathrm{km/h}$ on a clear stretch. To describe the motion \emph{at one instant}, we shrink the time interval.

Imagine computing the average velocity $\dfrac{\Delta x}{\Delta t}$ over the interval $[t,\ t+\Delta t]$, and letting $\Delta t$ become smaller and smaller. The ratio settles down to a definite value: the \cle{instantaneous velocity} at time $t$.

\begin{definition}[Instantaneous velocity]
\[
v = \lim_{\Delta t \to 0} \frac{\Delta x}{\Delta t} = \frac{dx}{dt}.
\]
The instantaneous velocity is the gradient of the position--time graph at time $t$ (we exploit this fully in Section 3). Its magnitude is the instantaneous \cle{speed}.
\end{definition}

\begin{attention}[Speed is not velocity]
Speed is a scalar (how fast); velocity is a vector (how fast \emph{and in which direction}). Two cars moving at $50\ \mathrm{km/h}$ in opposite directions have the \emph{same speed} but \emph{opposite velocities}. In a problem with a chosen positive direction, a velocity of $-8\ \mathrm{m/s}$ means a speed of $8\ \mathrm{m/s}$ in the negative direction. Never drop the sign, and never report a negative \emph{speed}.
\end{attention}

\courssub{Acceleration}

Vehicles rarely keep a constant velocity: they speed up, slow down, stop. \cle{Acceleration} measures how quickly the velocity changes.

\begin{definition}[Acceleration]
The \cle{average acceleration} over a time interval is
\[
a = \frac{\Delta v}{\Delta t} = \frac{v_{\text{final}} - v_{\text{initial}}}{\Delta t},
\]
and the instantaneous acceleration is $a = \dfrac{dv}{dt} = \dfrac{d^2x}{dt^2}$. Acceleration is a \textbf{vector}; its SI unit is the metre per second squared, $\mathrm{m\,s^{-2}}$.
\end{definition}

The sign of $a$ is read along the chosen positive axis:
\begin{itemize}
\item $a > 0$: the velocity is increasing in the positive direction (speeding up if moving positively, slowing down if moving negatively);
\item $a < 0$: the velocity is decreasing in the positive direction.
\end{itemize}

\begin{definition}[Deceleration (retardation)]
When a particle is \emph{slowing down}, its acceleration is opposite to its velocity. We then speak of \cle{deceleration} or \cle{retardation}: a deceleration of $3\ \mathrm{m\,s^{-2}}$ means an acceleration $a = -3\ \mathrm{m\,s^{-2}}$ (with the direction of motion taken as positive).
\end{definition}

\begin{attention}[Negative acceleration is not always slowing down]
A negative acceleration only means the acceleration vector points in the negative direction. If the particle is \emph{already moving} in the negative direction ($v<0$), a negative acceleration makes it go \emph{faster}. Deceleration occurs exactly when $a$ and $v$ have \textbf{opposite signs}.
\end{attention}

\begin{exemplebox}[Reading acceleration]
A car moving in the positive direction increases its velocity from $10\ \mathrm{m/s}$ to $22\ \mathrm{m/s}$ in $4$ s:
\[
a = \frac{22 - 10}{4} = +3\ \mathrm{m\,s^{-2}}.
\]
If instead it brakes from $22\ \mathrm{m/s}$ to $10\ \mathrm{m/s}$ in $4$ s, then $a = \dfrac{10-22}{4} = -3\ \mathrm{m\,s^{-2}}$: a deceleration of $3\ \mathrm{m\,s^{-2}}$.
\end{exemplebox}

\courssub{Uniform and Non-Uniform Acceleration}

Acceleration itself may change with time (a car in town traffic) or stay constant (a freely falling stone, neglecting air resistance).

\begin{definition}[Uniform acceleration]
Acceleration is \cle{uniform} (constant) when the velocity changes by the same amount in every equal time interval: $a$ is a constant. Otherwise the acceleration is \cle{non-uniform} (variable).
\end{definition}

This chapter focuses almost entirely on \textbf{uniform acceleration}, for two reasons: it describes many real situations well (free fall, a vehicle braking firmly), and it leads to the simple, powerful constant-acceleration formulae of Section 2. Variable acceleration requires calculus and is treated later in the course.

\begin{methode}[Choosing signs and keeping them]
\begin{enumerate}
\item Draw a quick sketch of the situation and mark the straight line of motion.
\item Choose an origin $O$ and a positive direction; write this choice down explicitly.
\item Express every given quantity (position, velocity, acceleration) as a signed number along this axis.
\item Carry the signs through every formula; never ``fix'' a sign at the end by guesswork.
\item Interpret the sign of each answer in words (``the particle is $3$ m to the left of $O$'', ``it is moving in the negative direction'').
\end{enumerate}
\end{methode}

\begin{popfigure}
\begin{center}
\begin{tikzpicture}[scale=0.9]
\draw[->, thick, popmuted] (-0.5,0) -- (11.5,0) node[below left] {road};
\foreach \xx/\tt in {0/0, 4/2, 8/4}
  \draw (\xx,0.1) -- (\xx,-0.1) node[below] {\small $t=\tt$ s};
% motorbike at t=0
\begin{scope}[shift={(0,0.35)}]
\draw[fill=popblueL, draw=popblue] (0,0.15) circle (0.14);
\draw[fill=popblueL, draw=popblue] (0.7,0.15) circle (0.14);
\draw[popblue, thick] (0,0.15) -- (0.35,0.5) -- (0.7,0.15) -- (0.25,0.15) -- cycle;
\draw[popblue, thick] (0.35,0.5) -- (0.55,0.62);
\end{scope}
% motorbike at t=2
\begin{scope}[shift={(4,0.35)}]
\draw[fill=popblueL, draw=popblue] (0,0.15) circle (0.14);
\draw[fill=popblueL, draw=popblue] (0.7,0.15) circle (0.14);
\draw[popblue, thick] (0,0.15) -- (0.35,0.5) -- (0.7,0.15) -- (0.25,0.15) -- cycle;
\draw[popblue, thick] (0.35,0.5) -- (0.55,0.62);
\end{scope}
% motorbike at t=4
\begin{scope}[shift={(8,0.35)}]
\draw[fill=popblueL, draw=popblue] (0,0.15) circle (0.14);
\draw[fill=popblueL, draw=popblue] (0.7,0.15) circle (0.14);
\draw[popblue, thick] (0,0.15) -- (0.35,0.5) -- (0.7,0.15) -- (0.25,0.15) -- cycle;
\draw[popblue, thick] (0.35,0.5) -- (0.55,0.62);
\end{scope}
\draw[->, very thick, poporange] (0.9,1.35) -- (2.6,1.35) node[midway, above, poporange] {\small $v_0$};
\draw[->, very thick, poporange] (4.9,1.35) -- (6.2,1.35) node[midway, above, poporange] {\small $v_2$};
\draw[->, very thick, poporange] (8.9,1.35) -- (9.7,1.35) node[midway, above, poporange] {\small $v_4$};
\node[popdark, align=center] at (5.5,-0.9) {\small braking: equal times, shrinking velocity arrows};
\end{tikzpicture}
\end{center}
\legende{A bend-skin braking on a straight road: positions at $t = 0, 2, 4$ s with velocity arrows. The arrows shrink steadily --- the signature of uniform (constant) deceleration.}
\end{popfigure}

\begin{exoresolu}[Velocity and acceleration with signs]
A particle moves along the $x$-axis (positive to the right). At $t = 1$ s its velocity is $+6\ \mathrm{m/s}$; at $t = 5$ s its velocity is $-2\ \mathrm{m/s}$. (a) Find its average acceleration over this interval. (b) Explain in words what happened to the particle.
\tcblower
\textbf{Solution.}
\begin{itemize}
\item (a) $a = \dfrac{\Delta v}{\Delta t} = \dfrac{-2 - 6}{5 - 1} = \dfrac{-8}{4} = -2\ \mathrm{m\,s^{-2}}$.
\item (b) The particle was first moving to the right at $6\ \mathrm{m/s}$. The constant negative acceleration slowed it down, brought it momentarily to rest, then made it move to the left, reaching $2\ \mathrm{m/s}$ leftwards by $t = 5$ s. Note the trap: the acceleration is negative throughout, yet the particle only \emph{decelerates} while $v > 0$; once $v < 0$, the same negative acceleration \emph{speeds it up} in the negative direction.
\end{itemize}
\end{exoresolu}

\begin{exercice}[]
A student jogs $300$ m due north in $150$ s, then $100$ m due south in $50$ s. Taking north as positive, find (a) the distance travelled, (b) the displacement, (c) the average speed, (d) the average velocity. Convert the average speed to $\mathrm{km\,h^{-1}}$.
\end{exercice}

\begin{corrige}[Exercise 1]
(a) Distance $= 300 + 100 = 400\ \mathrm{m}$. (b) Displacement $= 300 - 100 = +200\ \mathrm{m}$ (north). (c) Average speed $= \dfrac{400}{200} = 2\ \mathrm{m\,s^{-1}} = 2 \times 3.6 = 7.2\ \mathrm{km\,h^{-1}}$. (d) Average velocity $= \dfrac{200}{200} = +1\ \mathrm{m\,s^{-1}}$, directed north.
\end{corrige}

\begin{aretenir}[Key points of Section 1]
\begin{itemize}
\item A \cle{particle} is a point mass: a valid model when size is negligible compared with distances.
\item \cle{Displacement} $\Delta x = x_f - x_i$ is a vector; \cle{distance} is a scalar, and distance $\geq |\text{displacement}|$.
\item Average velocity $= \dfrac{\Delta x}{\Delta t}$ (vector); average speed $= \dfrac{\text{distance}}{\text{time}}$ (scalar). Instantaneous velocity $v = \dfrac{dx}{dt}$.
\item Acceleration $a = \dfrac{\Delta v}{\Delta t}$, in $\mathrm{m\,s^{-2}}$; deceleration means $a$ and $v$ have opposite signs.
\item $1\ \mathrm{m\,s^{-1}} = 3.6\ \mathrm{km\,h^{-1}}$: always convert to SI units before calculating.
\item Choose a positive direction once, and keep every sign consistent to the end.
\end{itemize}
\end{aretenir}

\coursec{Section 2: Equations of Motion under Uniform Acceleration (Lesson 93)}

In Section 1 we learned to describe motion in one dimension using \cle{displacement} $s$, \cle{velocity} $v$ and \cle{acceleration} $a$, and we saw that acceleration is the rate of change of velocity. We now specialise to the most important case in the whole of GCE kinematics: motion in a straight line with \textbf{constant acceleration}. This case is important for two reasons. First, it occurs very often in practice: a car pulling away from a junction in Bamenda, a taxi braking at a checkpoint, a stone falling freely near the ground (Section 5). Second, when $a$ is constant, the definitions of Section 1 can be turned into a small set of simple algebraic equations --- the famous \emph{suvat equations} --- which allow us to solve almost any straight-line problem at Advanced Level.

\courssub{The Uniform Acceleration Assumption and the suvat Variables}

\begin{definition}[Uniformly accelerated motion]
A particle moves with \cle{uniform acceleration} (or \emph{constant acceleration}) when its acceleration $a$ keeps the same value --- in magnitude \emph{and} in sign --- throughout the motion. Its velocity then changes by equal amounts in equal intervals of time.
\end{definition}

Throughout this chapter we use the five standard letters, collectively called the \cle{suvat variables}:

\begin{aretenir}[The five suvat variables]
\begin{itemize}
\item $s$ : displacement from the starting point (in metres, m);
\item $u$ : initial velocity, i.e.\ the velocity at time $t=0$ (in $\mathrm{m\,s^{-1}}$);
\item $v$ : final velocity, i.e.\ the velocity at time $t$ (in $\mathrm{m\,s^{-1}}$);
\item $a$ : constant acceleration (in $\mathrm{m\,s^{-2}}$);
\item $t$ : time elapsed (in seconds, s).
\end{itemize}
All five are \emph{signed} quantities: each may be positive or negative depending on the chosen positive direction.
\end{aretenir}

\courssub{First Equation: $v = u + at$}

Start from the definition of constant acceleration: acceleration is the change in velocity divided by the time taken,
\[
a=\frac{v-u}{t}.
\]
Multiplying both sides by $t$ and adding $u$ gives the first equation.

\begin{propriete}[First equation of motion]
For motion with constant acceleration,
\[
\boxed{v=u+at}
\]
This links $u$, $v$, $a$ and $t$; the variable $s$ is absent. The derivation from the definition of acceleration is examinable.
\end{propriete}

Intuitively: the velocity gained in time $t$ is $at$ (acceleration $\times$ time), so the final velocity is the initial velocity plus this gain.

\courssub{Second and Third Equations: the Average Velocity Idea}

Here is the key intuition. If the velocity changes \emph{steadily} from $u$ to $v$, then the average velocity over the journey is simply the mean of the two:
\[
\bar{v}=\frac{u+v}{2}.
\]
But displacement is average velocity multiplied by time, so
\[
s=\frac{u+v}{2}\,t.
\]
This is the second equation. To obtain the third, substitute $v=u+at$ into it:
\[
s=\frac{u+(u+at)}{2}\,t=\frac{2u+at}{2}\,t=ut+\tfrac{1}{2}at^{2}.
\]

\begin{propriete}[Second and third equations of motion]
\[
\boxed{s=\tfrac{1}{2}(u+v)t}\qquad\text{(no }a\text{)}\qquad\qquad \boxed{s=ut+\tfrac{1}{2}at^{2}}\qquad\text{(no }v\text{)}
\]
The term $ut$ is the displacement the particle would have covered at constant velocity $u$; the term $\tfrac12 at^{2}$ is the extra displacement due to the acceleration.
\end{propriete}

\courssub{Fourth Equation: Eliminating the Time}

Sometimes the time is neither given nor required. From $v=u+at$ we get $t=\dfrac{v-u}{a}$ (for $a\neq 0$). Substituting into $s=\tfrac12(u+v)t$:
\[
s=\frac{u+v}{2}\cdot\frac{v-u}{a}=\frac{v^{2}-u^{2}}{2a},
\]
since $(u+v)(v-u)=v^{2}-u^{2}$. Multiplying by $2a$ and rearranging:

\begin{propriete}[Fourth equation of motion]
\[
\boxed{v^{2}=u^{2}+2as}\qquad\text{(no }t\text{)}
\]
This is the equation to use whenever time does not appear in the problem.
\end{propriete}

\courssub{The Fifth Equation and the Summary}

Finally, write $u=v-at$ (rearranging the first equation) and substitute into $s=\tfrac12(u+v)t$:
\[
s=\frac{(v-at)+v}{2}\,t=vt-\tfrac{1}{2}at^{2}.
\]

\begin{propriete}[The five equations of uniformly accelerated motion]
\[
\begin{array}{lll}
v=u+at & \text{missing } s\\[2pt]
s=\tfrac{1}{2}(u+v)t & \text{missing } a\\[2pt]
s=ut+\tfrac{1}{2}at^{2} & \text{missing } v\\[2pt]
v^{2}=u^{2}+2as & \text{missing } t\\[2pt]
s=vt-\tfrac{1}{2}at^{2} & \text{missing } u
\end{array}
\]
Each equation involves exactly four of the five variables; each \emph{omits} a different one. These formulae are valid \textbf{only when $a$ is constant}.
\end{propriete}

\begin{popfigure}
\begin{tikzpicture}[font=\small]
\tikzset{eqbox/.style={draw=popblue, thick, rounded corners=3pt, fill=popblueL, minimum width=4.6cm, minimum height=0.95cm, align=center},
miss/.style={draw=poporange, thick, rounded corners=3pt, fill=white, minimum width=2.6cm, minimum height=0.95cm, align=center}}
\node[eqbox] (e1) at (0,4) {$v=u+at$};
\node[eqbox] (e2) at (0,3) {$s=\tfrac12(u+v)t$};
\node[eqbox] (e3) at (0,2) {$s=ut+\tfrac12 at^{2}$};
\node[eqbox] (e4) at (0,1) {$v^{2}=u^{2}+2as$};
\node[eqbox] (e5) at (0,0) {$s=vt-\tfrac12 at^{2}$};
\node[miss] (m1) at (5.2,4) {no $s$};
\node[miss] (m2) at (5.2,3) {no $a$};
\node[miss] (m3) at (5.2,2) {no $v$};
\node[miss] (m4) at (5.2,1) {no $t$};
\node[miss] (m5) at (5.2,0) {no $u$};
\foreach \i in {1,...,5}{\draw[-, popmuted, thick] (e\i) -- (m\i);}
\node[align=center, text=popdark] at (0,4.85) {\textbf{Equation}};
\node[align=center, text=popdark] at (5.2,4.85) {\textbf{Absent variable}};
\end{tikzpicture}
\legende{The five suvat equations: each one leaves out exactly one variable. To choose an equation, look at which variable is neither given nor wanted.}
\end{popfigure}

\courssub{Strategy: Choosing the Right Equation}

\begin{methode}[The four-step suvat strategy]
\begin{enumerate}
\item \textbf{List.} State the positive direction, then write down the values of the suvat variables that are known, converting all units to metres and seconds.
\item \textbf{Choose.} Identify the unknown you want and the variable that is neither given nor wanted; pick the equation that omits that variable.
\item \textbf{Substitute.} Insert the known values, keeping signs carefully (a deceleration is a \emph{negative} acceleration).
\item \textbf{Solve} and give the answer with correct units and a sensible number of significant figures.
\end{enumerate}
\end{methode}

\begin{attention}[Exam tip: direction first]
Before substituting anything, \textbf{always state the positive direction} and write down the known values with their signs. Most suvat errors at GCE come from sign mistakes made in the first line, not from algebra.
\end{attention}

\begin{exoresolu}[A car accelerating from rest]
A car starts from rest at a junction and accelerates uniformly at $2.5\ \mathrm{m\,s^{-2}}$ for $8\ \mathrm{s}$. Find (a) its velocity after $8\ \mathrm{s}$; (b) the distance covered in this time.
\tcblower
\textbf{Solution.} Take the direction of motion as positive. Known: $u=0$, $a=2.5$, $t=8$.

(a) Want $v$; $s$ is absent, so use $v=u+at$:
\[
v=0+2.5\times 8=20\ \mathrm{m\,s^{-1}}.
\]

(b) Want $s$; use $s=ut+\tfrac12 at^{2}$:
\[
s=0\times 8+\tfrac12\times 2.5\times 8^{2}=1.25\times 64=80\ \mathrm{m}.
\]
\emph{Check} with $v^{2}=u^{2}+2as$: $v^{2}=2\times 2.5\times 80=400$, so $v=20\ \mathrm{m\,s^{-1}}$. \checkmark
\end{exoresolu}

\begin{exoresolu}[A taxi braking to a stop]
A taxi travelling at $25\ \mathrm{m\,s^{-1}}$ brakes uniformly and stops in a distance of $50\ \mathrm{m}$. Find (a) the acceleration; (b) the time taken to stop.
\tcblower
\textbf{Solution.} Positive direction = direction of motion. Known: $u=25$, $v=0$, $s=50$.

(a) Want $a$; $t$ is absent, so use $v^{2}=u^{2}+2as$:
\[
0=25^{2}+2a(50)\;\Longrightarrow\;100a=-625\;\Longrightarrow\;a=-6.25\ \mathrm{m\,s^{-2}}.
\]
The negative sign confirms a deceleration: the taxi slows at $6.25\ \mathrm{m\,s^{-2}}$.

(b) Want $t$; use $v=u+at$:
\[
0=25-6.25t\;\Longrightarrow\;t=\frac{25}{6.25}=4\ \mathrm{s}.
\]
\end{exoresolu}

\begin{attention}[Sign errors with deceleration]
When a particle slows down while moving in the positive direction, its acceleration is \textbf{negative}. Never write $v^{2}=u^{2}-2as$ ``by hand'' to force a deceleration: keep the standard formula $v^{2}=u^{2}+2as$ and let $a<0$ do the work. Mixing the two conventions is the commonest source of wrong braking distances.
\end{attention}

\begin{attention}[Using suvat when $a$ is not constant]
The five equations are valid \emph{only} for constant acceleration. If a problem states that acceleration varies (for example $a$ depends on $t$ or on $v$), suvat does not apply: you must return to calculus, $v=\dfrac{ds}{dt}$ and $a=\dfrac{dv}{dt}$. Always check the ``uniform acceleration'' hypothesis before reaching for a suvat formula.
\end{attention}

\courssub{Two-Phase and Multi-Phase Journeys}

Real journeys rarely have a single acceleration. A typical trip consists of \emph{phases}: accelerate, cruise at constant velocity, then brake. The method is always the same: \textbf{treat each phase separately}, using the final velocity of one phase as the initial velocity of the next.

\begin{exoresolu}[Accelerate, cruise, brake]
A bus accelerates uniformly from rest at $1.5\ \mathrm{m\,s^{-2}}$ for $10\ \mathrm{s}$, travels at constant velocity for $20\ \mathrm{s}$, then brakes uniformly to rest in $5\ \mathrm{s}$. Find the total distance travelled.
\tcblower
\textbf{Solution.} Positive direction = direction of motion.

\textbf{Phase 1} ($u=0$, $a=1.5$, $t=10$):
\[
v_{1}=0+1.5\times 10=15\ \mathrm{m\,s^{-1}},\qquad s_{1}=\tfrac12(0+15)\times 10=75\ \mathrm{m}.
\]

\textbf{Phase 2} (constant velocity $15\ \mathrm{m\,s^{-1}}$, $t=20$, so $a=0$):
\[
s_{2}=15\times 20=300\ \mathrm{m}.
\]

\textbf{Phase 3} ($u=15$, $v=0$, $t=5$): use $s=\tfrac12(u+v)t$:
\[
s_{3}=\tfrac12(15+0)\times 5=37.5\ \mathrm{m}.
\]

\textbf{Total:} $s=75+300+37.5=412.5\ \mathrm{m}$.
\end{exoresolu}

\begin{popfigure}
\begin{tikzpicture}[font=\small, >=stealth]
\draw[->, thick] (0,0) -- (11.5,0) node[below left] {position};
\draw[popblue, very thick] (0,0.4) -- (3,0.4);
\draw[popgreen, very thick] (3,0.4) -- (7,0.4);
\draw[poporange, very thick] (7,0.4) -- (10,0.4);
\foreach \x/\lab in {0/{start}, 3/{$v=15$}, 7/{braking}, 10/{stop}}{
  \draw[popdark] (\x,0.32) -- (\x,0.48);
  \node[below, text=popdark] at (\x,0.28) {\lab};}
\node[above, text=popblue, align=center] at (1.5,0.45) {Phase 1: $a=1.5$, $t=10$ s};
\node[above, text=popgreen, align=center] at (5,0.45) {Phase 2: constant $v$, $20$ s};
\node[above, text=poporange, align=center] at (8.5,0.45) {Phase 3: brake, $5$ s};
\node[below, text=popmuted] at (1.5,-0.35) {$s_1=75$ m};
\node[below, text=popmuted] at (5,-0.35) {$s_2=300$ m};
\node[below, text=popmuted] at (8.5,-0.35) {$s_3=37.5$ m};
\end{tikzpicture}
\legende{A three-phase journey. The final velocity of each phase becomes the initial velocity of the next; distances add up.}
\end{popfigure}

\courssub{Two-Particle Problems (Enrichment for GCE)}

Some GCE questions involve \emph{two} particles moving along the same line, starting at the same time. Write a suvat expression for the displacement of \emph{each} particle as a function of the \emph{same} time $t$, then impose the condition of the problem --- typically a \cle{meeting condition}: the particles meet when their positions are equal.

\begin{exoresolu}[When does the second car catch the first?]
Car $A$ passes a point $O$ at $20\ \mathrm{m\,s^{-1}}$ and continues at constant velocity. At the same instant, car $B$ starts from rest at $O$ and accelerates uniformly at $4\ \mathrm{m\,s^{-2}}$ in the same direction. Find the time at which $B$ catches up with $A$, and the distance from $O$ at that moment.
\tcblower
\textbf{Solution.} Measure displacement from $O$, in the direction of motion, with $t$ the time since $B$ started.

Car $A$ (constant velocity): $s_A=20t$.

Car $B$ ($u=0$, $a=4$): $s_B=\tfrac12(4)t^{2}=2t^{2}$.

Meeting condition $s_A=s_B$:
\[
20t=2t^{2}\;\Longrightarrow\;2t(t-10)=0.
\]
$t=0$ is the start; the catching-up occurs at $t=10\ \mathrm{s}$, at distance $s=20\times 10=200\ \mathrm{m}$ from $O$.
\end{exoresolu}

\begin{exercice}[Braking distance]
A car travelling at $30\ \mathrm{m\,s^{-1}}$ brakes with uniform deceleration $5\ \mathrm{m\,s^{-2}}$. Find (a) the braking distance; (b) the time to stop; (c) the distance covered during the first $2\ \mathrm{s}$ of braking.
\end{exercice}

\begin{corrige}[Exercise 1]
Positive direction = direction of motion: $u=30$, $a=-5$.

(a) $v^{2}=u^{2}+2as$: $0=900-10s$, so $s=90\ \mathrm{m}$.

(b) $v=u+at$: $0=30-5t$, so $t=6\ \mathrm{s}$.

(c) $s=ut+\tfrac12 at^{2}=30(2)-\tfrac12(5)(2^{2})=60-10=50\ \mathrm{m}$.
\end{corrige}

\begin{exercice}[Two phases]
A motorcyclist accelerates uniformly from $5\ \mathrm{m\,s^{-1}}$ to $15\ \mathrm{m\,s^{-1}}$ over $100\ \mathrm{m}$, then rides at constant velocity for $30\ \mathrm{s}$. Find (a) the acceleration and the duration of phase 1; (b) the total distance travelled.
\end{exercice}

\begin{corrige}[Exercise 2]
(a) $v^{2}=u^{2}+2as$: $225=25+200a$, so $a=1\ \mathrm{m\,s^{-2}}$. Then $v=u+at$: $15=5+t$, so $t=10\ \mathrm{s}$.

(b) Phase 2: $s_2=15\times 30=450\ \mathrm{m}$. Total: $100+450=550\ \mathrm{m}$.
\end{corrige}

\begin{aretenir}[Key points of Section 2]
\begin{itemize}
\item The five suvat equations apply \textbf{only} to straight-line motion with \textbf{constant} acceleration.
\item Each equation omits one variable: choose the equation that omits the variable which is neither given nor wanted.
\item Always fix a positive direction first; deceleration means $a<0$.
\item Multi-phase journeys: treat each phase separately; the final velocity of one phase is the initial velocity of the next.
\item Two-particle problems: same clock for both particles; meeting means equal positions.
\end{itemize}
\end{aretenir}

In the next section we give these motions a face: we draw and interpret \emph{displacement--time} and \emph{velocity--time} graphs, and we shall see that the suvat equations reappear as gradients and areas.

\coursec{Section 3: Displacement–Time and Velocity–Time Graphs (Lessons 94, 95, 96)}

In Section 2 we derived the equations of motion for uniform acceleration using algebra and calculus. Those equations are powerful, but they assume constant acceleration. In many real situations — a bus leaving a station in Yaoundé, a ball thrown upwards during a physics practical, a car braking on the highway to Douala — the acceleration changes. Graphs give us a universal language to describe motion, whether acceleration is constant or not. They also provide a visual check on our calculations: the area under a velocity–time graph, for instance, always equals the displacement, no matter how the acceleration varies. In this section we master the two most important graphs in kinematics: displacement–time and velocity–time.

\courssub{Displacement–Time Graphs (x–t graphs)}

A displacement–time graph plots the position $x$ of a particle (measured from a fixed origin along a straight line) against time $t$. The horizontal axis represents time, the vertical axis represents displacement. The graph tells us \emph{where} the particle is at any instant.

\begin{propriete}[Gradient of a Displacement–Time Graph]
The gradient (slope) of the tangent to an $x$–$t$ graph at any point gives the \textbf{instantaneous velocity} $v$ at that time:
\[
v = \frac{dx}{dt} = \text{gradient of } x\text{–}t \text{ graph}.
\]
\begin{itemize}
\item A horizontal line (gradient $=0$) means the particle is \textbf{at rest}.
\item A straight line with constant positive gradient means \textbf{constant velocity} in the positive direction.
\item A straight line with constant negative gradient means \textbf{constant velocity} in the negative direction.
\item A curve means the velocity is changing; the steeper the curve, the greater the speed.
\end{itemize}
\end{propriete}

\begin{popfigure}
\begin{tikzpicture}
\begin{axis}[
    width=0.9\linewidth,
    height=7cm,
    axis lines=middle,
    xlabel={$t$ (s)}, ylabel={$x$ (m)},
    xmin=0, xmax=10,
    ymin=-2, ymax=12,
    xtick={0,2,4,6,8,10},
    ytick={0,2,4,6,8,10},
    tick label style={font=\small},
    label style={font=\small},
    clip=false,
    domain=0:10,
    samples=200
]
% Region 1: positive gradient (0 to 3)
\addplot[popblue, thick, domain=0:3] {2*x};
% Region 2: horizontal (3 to 6)
\addplot[poporange, thick, domain=3:6] {6};
% Region 3: negative gradient (6 to 10)
\addplot[popgreen, thick, domain=6:10] {-x+12};
% Annotations for time points
\draw[popblue, dashed] (3,0) node[below, font=\small] {$t_1$} -- (3,6);
\draw[poporange, dashed] (6,0) node[below, font=\small] {$t_2$} -- (6,6);
\draw[popgreen, dashed] (10,0) node[below, font=\small] {$t_3$} -- (10,2);
% Tangent lines for gradient illustration (corrected slopes)
% Region 1: slope 2. Tangent at t=1 (x=2): line through (1,2) slope 2 -> at t=2, x=4.
\draw[popmuted, dashed] (1,2) -- (2,4) node[above right, font=\tiny, popblue] {gradient $=v>0$};
% Region 2: slope 0.
\draw[popmuted, dashed] (4,6) -- (5,6) node[above right, font=\tiny, poporange] {gradient $=0$};
% Region 3: slope -1. Tangent at t=8 (x=4): line through (8,4) slope -1 -> at t=9, x=3.
\draw[popmuted, dashed] (8,4) -- (9,3) node[below right, font=\tiny, popgreen] {gradient $=v<0$};
% Points
\fill[popblue] (0,0) circle (2pt) node[below left, font=\small] {O};
\fill[popblue] (3,6) circle (2pt);
\fill[poporange] (6,6) circle (2pt);
\fill[popgreen] (10,2) circle (2pt);
\end{axis}
\end{tikzpicture}
\legende{Displacement–time graph with three distinct regions: (1) constant positive velocity, (2) at rest, (3) constant negative velocity (return journey). Tangents show the velocity in each region.}
\end{popfigure}

\begin{exemplebox}[Reading an $x$–$t$ Graph]
The graph above shows a particle that:
\begin{enumerate}
\item moves from the origin with constant velocity $v=2\ \mathrm{m\,s^{-1}}$ for $3\ \mathrm{s}$ (reaching $x=6\ \mathrm{m}$),
\item remains at rest at $x=6\ \mathrm{m}$ for $3\ \mathrm{s}$,
\item then returns towards the origin with constant velocity $v=-1\ \mathrm{m\,s^{-1}}$ for $4\ \mathrm{s}$, ending at $x=2\ \mathrm{m}$.
\end{enumerate}
\textbf{Questions:}
\begin{enumerate}
\item What is the displacement at $t=5\ \mathrm{s}$? \quad \textbf{Answer:} $x=6\ \mathrm{m}$ (read directly from graph).
\item What is the velocity at $t=5\ \mathrm{s}$? \quad \textbf{Answer:} $v=0$ (horizontal segment).
\item What is the displacement between $t=2\ \mathrm{s}$ and $t=8\ \mathrm{s}$? \quad \textbf{Answer:} $\Delta x = x(8)-x(2) = 4 - 4 = 0\ \mathrm{m}$.
\item What is the total distance travelled in the first $10\ \mathrm{s}$? \quad \textbf{Answer:} $6\ \mathrm{m} + 0\ \mathrm{m} + 4\ \mathrm{m} = 10\ \mathrm{m}$.
\end{enumerate}
\end{exemplebox}

\begin{aretenir}[Average Velocity from an $x$–$t$ Graph]
The \textbf{average velocity} between times $t_1$ and $t_2$ is the gradient of the \textbf{chord} joining the points $(t_1, x_1)$ and $(t_2, x_2)$:
\[
v_{\text{avg}} = \frac{x_2 - x_1}{t_2 - t_1} = \frac{\Delta x}{\Delta t}.
\]
This is the slope of the straight line connecting the two points on the curve, \emph{not} the slope of the tangent.
\end{aretenir}

\begin{attention}[Common Mistake: Confusing Position with Velocity]
On an $x$–$t$ graph, the \textbf{height} of the curve gives the position $x$, while the \textbf{slope} gives the velocity $v$. A common error is to read the value on the vertical axis at a given time and call it the velocity. Always remember: \emph{gradient = velocity}.
\end{attention}

\courssub{Velocity–Time Graphs (v–t graphs)}

A velocity–time graph plots velocity $v$ against time $t$. It tells us \emph{how fast} and \emph{in which direction} the particle is moving at each instant.

\begin{propriete}[Gradient and Area of a Velocity–Time Graph]
For a $v$–$t$ graph:
\begin{enumerate}
\item The \textbf{gradient} of the graph at any point equals the \textbf{acceleration} $a$:
\[
a = \frac{dv}{dt} = \text{gradient of } v\text{–}t \text{ graph}.
\]
\item The \textbf{area under the graph} between two times $t_1$ and $t_2$ equals the \textbf{displacement} $\Delta x$ during that interval (areas above the $t$-axis are positive, areas below are negative):
\[
\Delta x = \int_{t_1}^{t_2} v\,dt = \text{signed area under } v\text{–}t \text{ graph}.
\]
\end{enumerate}
\end{propriete}

\begin{aretenir}[Standard Shapes of $v$–$t$ Graphs for Uniform Acceleration]
\begin{itemize}
\item \textbf{Constant velocity:} horizontal line ($a=0$). Area = rectangle $\Rightarrow \Delta x = v \times \Delta t$.
\item \textbf{Uniform acceleration from rest:} straight line through origin with positive slope. Area = triangle $\Rightarrow \Delta x = \frac{1}{2} v_{\text{final}} \times t$.
\item \textbf{Uniform deceleration to rest:} straight line sloping down to the $t$-axis. Area = triangle.
\item \textbf{Combined journey (accelerate – cruise – brake):} trapezoid. Area = area of rectangle + two triangles.
\end{itemize}
\end{aretenir}

\begin{methode}[Computing Displacement from a $v$–$t$ Graph]
To find the displacement between $t_1$ and $t_2$:
\begin{enumerate}
\item Identify the geometric shapes formed by the graph and the $t$-axis (triangles, rectangles, trapezia).
\item Compute the area of each shape using standard formulae:
\begin{itemize}
\item Rectangle: $\text{base} \times \text{height}$.
\item Triangle: $\frac{1}{2} \times \text{base} \times \text{height}$.
\item Trapezium: $\frac{1}{2} \times (\text{sum of parallel sides}) \times \text{height}$.
\end{itemize}
\item Assign a \textbf{positive} sign to areas above the $t$-axis, \textbf{negative} to areas below.
\item Sum the signed areas to get the displacement $\Delta x$.
\item For \textbf{distance travelled}, sum the absolute values of the areas (ignore signs).
\end{enumerate}
\end{methode}

\begin{popfigure}
\begin{tikzpicture}
\begin{axis}[
    width=0.9\linewidth,
    height=7cm,
    axis lines=middle,
    xlabel={$t$ (s)}, ylabel={$v$ ($\mathrm{m\,s^{-1}}$)},
    xmin=0, xmax=30,
    ymin=-1, ymax=16,
    xtick={0,5,10,15,20,25,30},
    ytick={0,5,10,15},
    tick label style={font=\small},
    label style={font=\small},
    clip=false,
    domain=0:30,
    samples=200
]
% Bus journey: accelerate 0-5s, cruise 5-20s, brake 20-30s
\addplot[popblue, thick, domain=0:5] {3*x};          % acceleration: a=3 m/s^2, v=3t
\addplot[poporange, thick, domain=5:20] {15};        % constant velocity 15 m/s
\addplot[popgreen, thick, domain=20:30] {-1.5*x+45}; % deceleration: a=-1.5 m/s^2, v=-1.5t+45
% Shade area under curve
\addplot[popblueL, opacity=0.3, domain=0:5] {3*x} \closedcycle;
\addplot[poporangeL, opacity=0.3, domain=5:20] {15} \closedcycle;
\addplot[popgreenL, opacity=0.3, domain=20:30] {-1.5*x+45} \closedcycle;
% Annotations
\draw[dashed] (5,0) node[below, font=\small] {$5$} -- (5,15);
\draw[dashed] (20,0) node[below, font=\small] {$20$} -- (20,15);
\draw[dashed] (30,0) node[below, font=\small] {$30$} -- (30,0);
% Area labels
\node[popblue, font=\small] at (2.5, 5) {$\Delta x_1$};
\node[poporange, font=\small] at (12.5, 8) {$\Delta x_2$};
\node[popgreen, font=\small] at (25, 5) {$\Delta x_3$};
% Velocity label
\node[left, font=\small] at (0,15) {$15\ \mathrm{m\,s^{-1}}$};
\end{axis}
\end{tikzpicture}
\legende{Velocity–time graph of a bus journey between two stops: uniform acceleration (0–5 s), constant velocity (5–20 s), uniform deceleration (20–30 s). The total shaded area equals the distance between the stops.}
\end{popfigure}

\begin{exemplebox}[Bus Journey – Graphical Analysis]
A bus starts from rest at a bus stop, accelerates uniformly at $3\ \mathrm{m\,s^{-2}}$ for $5\ \mathrm{s}$, travels at constant velocity for $15\ \mathrm{s}$, then decelerates uniformly to rest at the next stop in $10\ \mathrm{s}$.
\begin{enumerate}
\item Sketch the $v$–$t$ graph (as above).
\item Find the maximum velocity reached. \quad $v_{\text{max}} = a t = 3 \times 5 = 15\ \mathrm{m\,s^{-1}}$.
\item Find the distance between the two stops (total area under graph).
\begin{align*}
\Delta x_1 &= \tfrac{1}{2} \times 5 \times 15 = 37.5\ \mathrm{m} \quad \text{(triangle)} \\
\Delta x_2 &= 15 \times 15 = 225\ \mathrm{m} \quad \text{(rectangle)} \\
\Delta x_3 &= \tfrac{1}{2} \times 10 \times 15 = 75\ \mathrm{m} \quad \text{(triangle)} \\
\text{Total distance} &= 37.5 + 225 + 75 = 337.5\ \mathrm{m}.
\end{align*}
\item Find the deceleration during braking. \quad $a = \frac{\Delta v}{\Delta t} = \frac{0-15}{10} = -1.5\ \mathrm{m\,s^{-2}}$.
\end{enumerate}
\end{exemplebox}

\begin{savaistu}[Recovering the suvat Equations Graphically]
The area of the trapezium in a $v$–$t$ graph for uniform acceleration gives $\Delta x = \frac{1}{2}(u+v)t$. Since $v = u+at$, substituting yields $\Delta x = ut + \frac{1}{2}at^2$. The other suvat equations follow similarly. This shows that the algebraic equations are just the geometric areas in disguise!
\end{savaistu}

\courssub{Displacement vs Distance from v–t Graphs}

When the velocity changes sign (the graph crosses the $t$-axis), the particle reverses direction. The \textbf{displacement} is the algebraic sum of the areas (positive above axis, negative below). The \textbf{distance travelled} is the sum of the absolute areas.

\begin{popfigure}
\begin{tikzpicture}
\begin{axis}[
    width=0.9\linewidth,
    height=7cm,
    axis lines=middle,
    xlabel={$t$ (s)}, ylabel={$v$ ($\mathrm{m\,s^{-1}}$)},
    xmin=0, xmax=10,
    ymin=-7, ymax=15,
    xtick={0,2,3,4,6,8,10},
    ytick={-6,-4,-2,0,2,4,6,8,10,12,14},
    tick label style={font=\small},
    label style={font=\small},
    clip=false,
    domain=0:10,
    samples=200
]
% v = 2t - 6 (crosses at t=3)
\addplot[poppurple, thick, domain=0:10] {2*x - 6};
% Shade positive area (t=3 to 10)
\addplot[popgreenL, opacity=0.3, domain=3:10] {2*x - 6} \closedcycle;
% Shade negative area (t=0 to 3)
\addplot[poppinkL, opacity=0.3, domain=0:3] {2*x - 6} \closedcycle;
% Annotations
\draw[dashed] (3,0) node[below, font=\small] {$3$} -- (3,0);
\draw[dashed] (10,0) node[below, font=\small] {$10$} -- (10,14);
% Area labels
\node[popgreen, font=\small, align=center] at (6.5, 4) {Positive area\\$A_+ = 49\ \mathrm{m}$};
\node[poppink, font=\small, align=center] at (1.5, -3) {Negative area\\$A_- = -9\ \mathrm{m}$};
% Velocity intercepts
\node[left, font=\small] at (0,-6) {$-6$};
\node[left, font=\small] at (0,14) {$14$};
\end{axis}
\end{tikzpicture}
\legende{Velocity–time graph crossing the time axis at $t=3\ \mathrm{s}$. The pink region (below axis) contributes negative displacement; the green region (above axis) contributes positive displacement.}
\end{popfigure}

\begin{exemplebox}[Displacement and Distance from a $v$–$t$ Graph Crossing the Axis]
The graph shows $v = 2t - 6$ for $0 \le t \le 10$.
\begin{enumerate}
\item \textbf{Displacement} from $t=0$ to $t=10$:
\[
\Delta x = \int_0^{10} (2t-6)\,dt = \big[t^2 - 6t\big]_0^{10} = 100 - 60 = 40\ \mathrm{m}.
\]
Geometrically: $A_- = -\frac{1}{2}\times 3 \times 6 = -9\ \mathrm{m}$ (triangle below axis), $A_+ = \frac{1}{2}\times 7 \times 14 = 49\ \mathrm{m}$ (triangle above axis). Sum $= 40\ \mathrm{m}$.
\item \textbf{Distance travelled} from $t=0$ to $t=10$:
\[
\text{Distance} = |A_-| + |A_+| = 9 + 49 = 58\ \mathrm{m}.
\]
\end{enumerate}
\end{exemplebox}

\begin{attention}[Common Mistake: Forgetting the Sign of Area Below the Axis]
When a $v$–$t$ graph goes below the $t$-axis, the area is \textbf{negative} for displacement but \textbf{positive} for distance. Many candidates lose marks by giving the distance as the net area (displacement) or by forgetting to take the absolute value of the negative part. Always ask yourself: \emph{Am I asked for displacement or distance?}
\end{attention}

\begin{attention}[Common Mistake: Reading $v$ as Displacement]
Another frequent error: at a given time $t$, reading the value $v$ on the vertical axis and calling it the displacement. Remember: on a $v$–$t$ graph, the \textbf{height} is velocity, the \textbf{area} is displacement.
\end{attention}

\courssub{Real-Life Applications and GCE-Style Questions}

In the GCE examination, you will often be given a worded description of a journey and asked to:
\begin{itemize}
\item sketch the $v$–$t$ graph (with correct shape, labels, and units on axes),
\item calculate the total distance or displacement,
\item find the acceleration at a given stage,
\item find the average speed or average velocity.
\end{itemize}

\begin{aretenir}[Exam Tips for Graph Questions]
\begin{itemize}
\item \textbf{Always label axes} with the quantity \textbf{and} its unit (e.g., $v\ (\mathrm{m\,s^{-1}})$, $t\ (\mathrm{s})$). Marks are awarded for this.
\item Use a sharp pencil and ruler for straight lines; draw curves smoothly.
\item Show key coordinates (intercepts, maximum/minimum points) on the graph.
\item When computing area, show the geometric breakdown (e.g., $\frac{1}{2}\times 4 \times 10$) — this earns method marks even if the final arithmetic is wrong.
\item For average speed: $\frac{\text{total distance}}{\text{total time}}$. For average velocity: $\frac{\text{total displacement}}{\text{total time}}$.
\end{itemize}
\end{aretenir}

\begin{exoresolu}[GCE-Style Question]
A car starts from rest at a traffic light and accelerates uniformly at $2\ \mathrm{m\,s^{-2}}$ for $10\ \mathrm{s}$. It then travels at constant velocity for $20\ \mathrm{s}$ before decelerating uniformly to rest in $5\ \mathrm{s}$ at the next traffic light.
\begin{enumerate}
\item Sketch the velocity–time graph for the whole journey.
\item Calculate the distance between the two traffic lights.
\item Find the deceleration during the final stage.
\item Calculate the average speed of the car for the whole journey.
\end{enumerate}
\tcblower
\textbf{Solution:}
\begin{enumerate}
\item \textbf{Graph:}
\begin{itemize}
\item $0 \le t \le 10$: $v = 2t$ (straight line from $(0,0)$ to $(10,20)$).
\item $10 \le t \le 30$: $v = 20$ (horizontal line).
\item $30 \le t \le 35$: $v = -4(t-30)+20 = -4t+140$ (straight line from $(30,20)$ to $(35,0)$).
\end{itemize}
Axes labelled $v\ (\mathrm{m\,s^{-1}})$ and $t\ (\mathrm{s})$.
\item \textbf{Distance} = total area under graph:
\begin{align*}
A_1 &= \tfrac{1}{2} \times 10 \times 20 = 100\ \mathrm{m} \\
A_2 &= 20 \times 20 = 400\ \mathrm{m} \\
A_3 &= \tfrac{1}{2} \times 5 \times 20 = 50\ \mathrm{m} \\
\text{Total distance} &= 100 + 400 + 50 = 550\ \mathrm{m}.
\end{align*}
\item \textbf{Deceleration} = gradient of last segment:
\[
a = \frac{0 - 20}{35 - 30} = \frac{-20}{5} = -4\ \mathrm{m\,s^{-2}}.
\]
\item \textbf{Average speed} = $\frac{\text{total distance}}{\text{total time}} = \frac{550}{35} = 15.7\ \mathrm{m\,s^{-1}}$ (to 3 s.f.).
\end{enumerate}
\end{exoresolu}

\begin{exercice}[Practice]
A particle moves along a straight line. Its velocity–time graph consists of three straight-line segments:
\begin{itemize}
\item From $t=0$ to $t=4\ \mathrm{s}$: velocity increases uniformly from $0$ to $8\ \mathrm{m\,s^{-1}}$.
\item From $t=4\ \mathrm{s}$ to $t=10\ \mathrm{s}$: velocity decreases uniformly to $-4\ \mathrm{m\,s^{-1}}$.
\item From $t=10\ \mathrm{s}$ to $t=14\ \mathrm{s}$: velocity increases uniformly to $0$.
\end{itemize}
\begin{enumerate}
\item Sketch the $v$–$t$ graph.
\item Find the acceleration during each stage.
\item Calculate the displacement of the particle at $t=14\ \mathrm{s}$ relative to its starting point.
\item Calculate the total distance travelled in the $14\ \mathrm{s}$.
\end{enumerate}
\end{exercice}

\begin{corrige}[Exercise]
\begin{enumerate}
\item Graph: three straight lines connecting $(0,0)$, $(4,8)$, $(10,-4)$, $(14,0)$.
\item Accelerations:
Stage 1: $a_1 = \frac{8-0}{4} = 2\ \mathrm{m\,s^{-2}}$.
Stage 2: $a_2 = \frac{-4-8}{6} = -2\ \mathrm{m\,s^{-2}}$.
Stage 3: $a_3 = \frac{0-(-4)}{4} = 1\ \mathrm{m\,s^{-2}}$.
\item Displacement = signed area:
$A_1 = \frac{1}{2}\times 4 \times 8 = 16\ \mathrm{m}$ (positive).
$A_2$: trapezium from $t=4$ to $10$ with heights $8$ and $-4$. Area = $\frac{1}{2}(8 + (-4)) \times 6 = 12\ \mathrm{m}$ (positive).
$A_3$: triangle below axis from $t=10$ to $14$ with height $-4$. Area = $\frac{1}{2}\times 4 \times (-4) = -8\ \mathrm{m}$.
Total displacement = $16 + 12 - 8 = 20\ \mathrm{m}$.
\item Distance = $|16| + |12| + |-8| = 36\ \mathrm{m}$.
\end{enumerate}
\end{corrige}

\coursec{Section 4: Acceleration--Time Graphs and Connecting the Three Graphs (Lessons 97, 98)}

In Section 3 we learned to read motion from \cle{displacement--time} and \cle{velocity--time} graphs: the gradient of an $x$--$t$ graph gives velocity, and the area under a $v$--$t$ graph gives displacement. We now complete the picture by adding the third member of the family, the \cle{acceleration--time graph}, and we learn to move freely between all three. This is one of the most frequently examined skills in the GCE Advanced Level mechanics paper: an examiner gives you \emph{one} graph and asks you to sketch or interpret the \emph{other two}.

\courssub{The Acceleration--Time Graph}

\begin{definition}[Acceleration--time graph]
An \cle{acceleration--time graph} (or $a$--$t$ graph) is a graph of the acceleration of a particle plotted on the vertical axis against time on the horizontal axis. Since $a = \dfrac{dv}{dt}$, the graph shows how quickly the velocity itself is changing at each instant.
\end{definition}

Because acceleration is the \emph{gradient} of velocity, the $a$--$t$ graph plays for the $v$--$t$ graph exactly the role that the $v$--$t$ graph plays for the $x$--$t$ graph. Everything we learned in Section 3 transfers directly, one step further down the chain:

\begin{center}
\begin{tabularx}{\linewidth}{|l|X|X|}
\hline
\rowcolor{popblueL} \textbf{Pair of graphs} & \textbf{Going down (differentiate)} & \textbf{Going up (integrate)} \\
\hline
$x$--$t$ and $v$--$t$ & gradient of $x$--$t$ $=$ velocity & area under $v$--$t$ $=$ displacement \\
\hline
$v$--$t$ and $a$--$t$ & gradient of $v$--$t$ $=$ acceleration & area under $a$--$t$ $=$ change in velocity \\
\hline
\end{tabularx}
\end{center}

\textbf{Reading the basic shapes.}\par
\begin{itemize}
\item A \textbf{horizontal line above the time axis} means \cle{constant positive acceleration}: velocity increases uniformly.
\item A \textbf{horizontal line below the time axis} means constant negative acceleration (deceleration if the particle is moving forward): velocity decreases uniformly.
\item The \textbf{line $a = 0$} (the graph lies on the time axis) means \emph{zero acceleration}, i.e.\ the velocity is \textbf{constant} --- not necessarily zero.
\item A \textbf{sloping line} means the acceleration itself is changing (non-uniform acceleration); at this level we mostly meet piecewise-constant accelerations.
\end{itemize}

\begin{attention}[Zero acceleration does not mean zero velocity]
The most common error at this point is to read ``$a = 0$'' as ``the particle is at rest''. Wrong. Zero acceleration means the velocity is \emph{not changing}: the particle moves with \textbf{constant velocity} (which may be large!). A bus cruising steadily at $\SI{60}{km.h^{-1}}$ along a straight stretch of the Douala--Yaound\'e road has zero acceleration. Only when the $v$--$t$ graph touches the time axis is the particle instantaneously at rest.
\end{attention}

\courssub{Area Under the Acceleration--Time Graph}

\begin{propriete}[Area under the $a$--$t$ graph]
The (signed) area between the $a$--$t$ graph and the time axis, between times $t_1$ and $t_2$, equals the \textbf{change in velocity} over that interval:
\[
\text{area under } a\text{--}t \text{ graph} \;=\; \Delta v \;=\; v(t_2) - v(t_1).
\]
Areas \emph{below} the time axis count negatively. The proof of this result (it is the Fundamental Theorem of Calculus applied to $a = \frac{dv}{dt}$) is not examinable, but the result itself must be known and used fluently.
\end{propriete}

\textbf{Why this is true (instructive derivation).}\par
For \emph{constant} acceleration $a$ over a time interval $\Delta t$, the defining equation $a = \dfrac{\Delta v}{\Delta t}$ gives $\Delta v = a\,\Delta t$, which is precisely the area of the rectangle of height $a$ and width $\Delta t$ under the graph. If the acceleration varies, split the interval into very small pieces of duration $\delta t$; on each piece the acceleration is nearly constant, so the velocity changes by approximately $a\,\delta t$, the area of a thin strip. Adding all the strips, the total change of velocity equals the total area. Making $\delta t \to 0$ turns this approximation into an exact equality: $\Delta v = \displaystyle\int_{t_1}^{t_2} a\,dt$.

\begin{attention}[Area gives the \emph{change} in velocity, not the velocity]
The area under an $a$--$t$ graph is $\Delta v$, \textbf{not} $v$. To obtain the velocity at time $t$ you must add the accumulated area to the \textbf{initial velocity}:
\[
v(t) = v(0) + (\text{signed area under the } a\text{--}t \text{ graph from } 0 \text{ to } t).
\]
Forgetting $v(0)$ is the second most common mistake in this topic. A particle can have zero area under its $a$--$t$ graph yet be moving fast, simply because it started fast.
\end{attention}

\courssub{Reconstructing $v$--$t$ from $a$--$t$: Accumulating Areas}

\begin{methode}[From the $a$--$t$ graph to the $v$--$t$ graph]
\begin{enumerate}
\item Write down the \textbf{initial velocity} $v(0)$ (given in the question).
\item Work \textbf{phase by phase}. On each phase where $a$ is constant, compute the signed area $a \times \Delta t$ and add it to the running velocity: $v_{\text{end}} = v_{\text{start}} + a\,\Delta t$.
\item Plot the velocity values at each phase boundary, then join them: constant acceleration means the $v$--$t$ graph is a \textbf{straight line} on each phase (rising if $a>0$, falling if $a<0$, horizontal if $a=0$).
\item Check: the gradient of each segment of your $v$--$t$ sketch must equal the value of $a$ on that phase.
\end{enumerate}
\end{methode}

\begin{exemplebox}[Accumulating areas]
A particle starts with velocity $v(0) = \SI{2}{m.s^{-1}}$. Its acceleration is $\SI{3}{m.s^{-2}}$ for $0 \le t \le 4$, then $0$ for $4 \le t \le 7$, then $\SI{-2}{m.s^{-2}}$ for $7 \le t \le 12$.

Phase 1: $\Delta v = 3 \times 4 = 12$, so $v(4) = 2 + 12 = \SI{14}{m.s^{-1}}$.

Phase 2: $\Delta v = 0$, so $v(7) = \SI{14}{m.s^{-1}}$ (constant velocity).

Phase 3: $\Delta v = (-2)\times 5 = -10$, so $v(12) = 14 - 10 = \SI{4}{m.s^{-1}}$.

The $v$--$t$ graph is therefore: straight line from $(0,2)$ to $(4,14)$, horizontal from $(4,14)$ to $(7,14)$, straight line down from $(7,14)$ to $(12,4)$. Note that the velocity stays positive throughout: the particle never stops, it only slows down.
\end{exemplebox}

\courssub{From $v$--$t$ to $x$--$t$: Accumulating Areas Again}

Exactly the same philosophy, one level up: the \textbf{area under the $v$--$t$ graph is the change in displacement}, so
\[
x(t) = x(0) + (\text{signed area under the } v\text{--}t \text{ graph from } 0 \text{ to } t).
\]
There is one new subtlety concerning \emph{shape}. On a phase of constant acceleration, the $v$--$t$ graph is a straight line, so the displacement accumulates like the area under a sloping line: the $x$--$t$ graph is a \textbf{curve (parabola)}, not a straight line. Concretely:
\begin{itemize}
\item velocity increasing $\Rightarrow$ $x$--$t$ curve bends \textbf{upwards} (concave up, getting steeper);
\item velocity decreasing (but positive) $\Rightarrow$ $x$--$t$ curve bends \textbf{downwards} (concave down, flattening);
\item constant velocity $\Rightarrow$ $x$--$t$ graph is a straight line whose gradient is that velocity.
\end{itemize}

\courssub{From $x$--$t$ to $v$--$t$ to $a$--$t$: Taking Gradients}

Going the other way down the chain, we \emph{differentiate graphically}: at each instant, the velocity is the gradient of the $x$--$t$ graph, and the acceleration is the gradient of the $v$--$t$ graph. For the piecewise-linear graphs met at this level:

\begin{itemize}
\item Each \textbf{straight segment} of the $x$--$t$ graph has constant gradient, so it produces a \textbf{horizontal segment} on the $v$--$t$ graph.
\item A \textbf{horizontal segment} of the $v$--$t$ graph has zero gradient, so it produces $a = 0$ on the $a$--$t$ graph.
\item A \textbf{corner} (sudden change of gradient) in the $x$--$t$ graph produces a \textbf{jump} in the $v$--$t$ graph; a corner in the $v$--$t$ graph produces a jump in the $a$--$t$ graph.
\item A \textbf{curved} part of the $x$--$t$ graph (parabola) produces a \textbf{sloping straight line} on the $v$--$t$ graph and a \textbf{non-zero horizontal line} on the $a$--$t$ graph.
\end{itemize}

\begin{methode}[The ``gradient going down, area going up'' rule]
To connect the three graphs, remember one sentence:
\begin{center}
\emph{Going down the chain ($x \to v \to a$), take \textbf{gradients}.\\
Going up the chain ($a \to v \to x$), accumulate \textbf{areas} (and add the initial value).}
\end{center}
Concretely:
\begin{enumerate}
\item \textbf{Down:} value of $v$ at time $t$ $=$ gradient of $x$--$t$ at $t$; value of $a$ at time $t$ $=$ gradient of $v$--$t$ at $t$.
\item \textbf{Up:} $v(t) = v(0) + $ area under $a$--$t$; $x(t) = x(0) + $ area under $v$--$t$.
\item Always mark the phase-change instants on \textbf{all three} graphs before sketching any curve.
\end{enumerate}
\end{methode}

\courssub{Complete Worked Chain: A Lift Journey}

\begin{exoresolu}[The lift journey]
A lift (elevator) in a Douala office block starts from rest at the ground floor. Its acceleration is modelled in three phases: $a = \SI{1.2}{m.s^{-2}}$ for $0 \le t \le 5$ (speeding up), $a = 0$ for $5 \le t \le 11$ (steady motion), $a = \SI{-1.5}{m.s^{-2}}$ for $11 \le t \le 15$ (slowing down to a stop). Take $x(0) = 0$. (a) Find the velocity at the end of each phase. (b) Sketch the $a$--$t$, $v$--$t$ and $x$--$t$ graphs. (c) Find the total height risen.
\tcblower
\textbf{(a) Velocities by accumulating areas.} The lift starts from rest: $v(0) = 0$.

Phase 1 ($0$ to $\SI{5}{s}$): $\Delta v = 1.2 \times 5 = 6$, so $v(5) = 0 + 6 = \SI{6}{m.s^{-1}}$.

Phase 2 ($5$ to $\SI{11}{s}$): $\Delta v = 0 \times 6 = 0$, so $v(11) = \SI{6}{m.s^{-1}}$.

Phase 3 ($11$ to $\SI{15}{s}$): $\Delta v = (-1.5) \times 4 = -6$, so $v(15) = 6 - 6 = 0$. The lift stops exactly at $t = \SI{15}{s}$, as expected.

\textbf{(b) The three graphs.}
\begin{itemize}
\item $a$--$t$: horizontal segments at $1.2$, then $0$, then $-1.5$.
\item $v$--$t$: straight line from $(0,0)$ to $(5,6)$; horizontal from $(5,6)$ to $(11,6)$; straight line from $(11,6)$ down to $(15,0)$. The gradient of each piece ($1.2$, $0$, $-1.5$) matches the acceleration.
\item $x$--$t$: parabola bending upwards from $(0,0)$ to $(5,15)$; straight line of gradient $6$ from $(5,15)$ to $(11,51)$; parabola bending downwards, flattening to a horizontal tangent, from $(11,51)$ to $(15,63)$.
\end{itemize}

\textbf{(c) Height risen = area under the $v$--$t$ graph.}

Phase 1 (triangle): $\tfrac12 \times 5 \times 6 = \SI{15}{m}$, so $x(5) = \SI{15}{m}$.

Phase 2 (rectangle): $6 \times 6 = \SI{36}{m}$, so $x(11) = 15 + 36 = \SI{51}{m}$.

Phase 3 (triangle): $\tfrac12 \times 4 \times 6 = \SI{12}{m}$, so $x(15) = 51 + 12 = \SI{63}{m}$.

The lift rises a total of $\SI{63}{m}$. \textbf{Consistency check:} the total signed area under the $a$--$t$ graph is $1.2\times5 + 0 - 1.5\times4 = 0$, matching the fact that the velocity returns to its initial value $0$.
\end{exoresolu}

\begin{popfigure}
\begin{center}
\begin{tikzpicture}
\begin{axis}[
  name=pa,
  width=11cm, height=3.6cm,
  axis lines=left,
  xlabel={$t$ (s)}, ylabel={$a$ (m\,s$^{-2}$)},
  xmin=0, xmax=16.5, ymin=-2.4, ymax=2.4,
  xtick={0,5,11,15}, ytick={-1.5,0,1.2},
  yticklabels={$-1.5$,$0$,$1.2$},
  clip=false, label style={font=\small}, tick label style={font=\small}
]
\fill[popblueL, opacity=0.7] (axis cs:0,0) rectangle (axis cs:5,1.2);
\fill[poporange, opacity=0.35] (axis cs:11,0) rectangle (axis cs:15,-1.5);
\draw[very thick, popblue] (axis cs:0,1.2) -- (axis cs:5,1.2);
\draw[very thick, popblue] (axis cs:5,0) -- (axis cs:11,0);
\draw[very thick, popblue] (axis cs:11,-1.5) -- (axis cs:15,-1.5);
\draw[dashed, popmuted] (axis cs:5,-2.4) -- (axis cs:5,2.4);
\draw[dashed, popmuted] (axis cs:11,-2.4) -- (axis cs:11,2.4);
\node[font=\small, popdark] at (axis cs:2.5,0.6) {$+6$};
\node[font=\small, popdark] at (axis cs:13,-0.75) {$-6$};
\end{axis}
\begin{axis}[
  name=pv, at={($(pa.south)-(0,2.6cm)$)}, anchor=north,
  width=11cm, height=3.6cm,
  axis lines=left,
  xlabel={$t$ (s)}, ylabel={$v$ (m\,s$^{-1}$)},
  xmin=0, xmax=16.5, ymin=0, ymax=8,
  xtick={0,5,11,15}, ytick={0,6},
  clip=false, label style={font=\small}, tick label style={font=\small}
]
\draw[very thick, popgreen] (axis cs:0,0) -- (axis cs:5,6) -- (axis cs:11,6) -- (axis cs:15,0);
\draw[dashed, popmuted] (axis cs:5,0) -- (axis cs:5,8);
\draw[dashed, popmuted] (axis cs:11,0) -- (axis cs:11,8);
\node[font=\small, popdark] at (axis cs:2.5,2.2) {gradient $1.2$};
\node[font=\small, popdark] at (axis cs:8,6.6) {gradient $0$};
\node[font=\small, popdark] at (axis cs:13.4,2.2) {gradient $-1.5$};
\end{axis}
\begin{axis}[
  name=px, at={($(pv.south)-(0,2.6cm)$)}, anchor=north,
  width=11cm, height=4.2cm,
  axis lines=left,
  xlabel={$t$ (s)}, ylabel={$x$ (m)},
  xmin=0, xmax=16.5, ymin=0, ymax=72,
  xtick={0,5,11,15}, ytick={0,15,51,63},
  clip=false, label style={font=\small}, tick label style={font=\small}
]
\draw[very thick, poppurple] (axis cs:0,0) parabola[bend at start] (axis cs:5,15);
\draw[very thick, poppurple] (axis cs:5,15) -- (axis cs:11,51);
\draw[very thick, poppurple] (axis cs:11,51) parabola[bend at end] (axis cs:15,63);
\draw[dashed, popmuted] (axis cs:5,0) -- (axis cs:5,72);
\draw[dashed, popmuted] (axis cs:11,0) -- (axis cs:11,72);
\draw[dashed, gray] (axis cs:0,15) -- (axis cs:5,15);
\draw[dashed, gray] (axis cs:0,51) -- (axis cs:11,51);
\draw[dashed, gray] (axis cs:0,63) -- (axis cs:15,63);
\end{axis}
\end{tikzpicture}
\end{center}
\legende{The lift journey on three aligned graphs. The dashed vertical lines mark the phase changes at $t = 5$ and $t = \SI{11}{s}$. Going down: gradients. Going up: areas (the shaded areas on the $a$--$t$ graph, $+6$ and $-6$, are the velocity changes).}
\end{popfigure}

\courssub{Consistency Checks Between the Three Graphs}

Before handing in any sketch, run through this checklist. Each row is a direct consequence of ``gradient going down, area going up''.

\begin{center}
\begin{tabularx}{\linewidth}{|l|X|X|}
\hline
\rowcolor{popblueL} \textbf{Feature on a graph} & \textbf{Consequence on the graph immediately below} & \textbf{Consequence on the graph two steps below} \\
\hline
$x$--$t$ has a maximum or minimum (horizontal tangent) & $v$--$t$ crosses zero (sign change of $v$) & $a$--$t$ \dots (no fixed rule) \\
\hline
$v$--$t$ has a maximum or minimum & $a$--$t$ crosses zero (sign change of $a$) & $x$--$t$ has an inflection point (change of concavity) \\
\hline
$v$--$t$ is horizontal & $a$--$t$ lies on the time axis ($a = 0$) & $x$--$t$ is a straight line \\
\hline
$a$--$t$ is positive & $v$--$t$ is increasing & $x$--$t$ bends upwards (concave up) \\
\hline
$a$--$t$ is negative & $v$--$t$ is decreasing & $x$--$t$ bends downwards (concave down) \\
\hline
$v$--$t$ crosses the time axis (changes sign) & particle reverses direction & $x$--$t$ has a turning point (maximum or minimum) \\
\hline
\end{tabularx}
\end{center}

\begin{attention}[Deceleration is not ``negative velocity'']
A negative acceleration does \emph{not} mean the particle moves backwards. In the lift example, $a = \SI{-1.5}{m.s^{-2}}$ during the last phase, yet the lift keeps moving \emph{upwards} (velocity still positive) --- it simply slows down. Negative acceleration only means the velocity is \emph{decreasing}. The particle reverses direction only when the \emph{velocity} changes sign.
\end{attention}

\begin{exemplebox}[Quick consistency test]
A $v$--$t$ graph reaches a maximum of $\SI{10}{m.s^{-1}}$ at $t = \SI{4}{s}$ and crosses the time axis at $t = \SI{9}{s}$. What must the other graphs show? \textbf{Answer:} at $t = 4$ the $a$--$t$ graph crosses zero (from positive to negative, since the velocity stops increasing and starts decreasing); at $t = 9$ the particle reverses direction, so the $x$--$t$ graph has a turning point (a maximum) there.
\end{exemplebox}

\courssub{Enrichment: The Link with Calculus (Beyond the Syllabus)}

\begin{savaistu}[A preview of Further Mathematics]
Everything in this section is secretly calculus. If $x(t)$ is the displacement, then
\[
v = \frac{dx}{dt} = \dot{x}, \qquad a = \frac{dv}{dt} = \frac{d^2x}{dt^2} = \ddot{x},
\]
and conversely $v(t) = v(0) + \displaystyle\int_0^t a\,ds$ and $x(t) = x(0) + \displaystyle\int_0^t v\,ds$. ``Gradient going down'' \emph{is} differentiation; ``area going up'' \emph{is} integration. For example, constant acceleration $a$ integrates to $v = u + at$ and again to $x = ut + \tfrac12 at^2$ --- the very equations of Lesson 93. You will meet this formally in Further Mathematics, where the same ideas handle \emph{non-uniform} acceleration such as $a = 3t^2 - 2t$.
\end{savaistu}

\begin{exercice}[Reconstructing the graphs]
A motorbike taxi (\emph{bendskin}) in Bamenda starts from rest. Its acceleration is $\SI{2}{m.s^{-2}}$ for $0 \le t \le 6$, then $\SI{-1}{m.s^{-2}}$ for $6 \le t \le 10$, then $0$ for $10 \le t \le 16$.
\begin{enumerate}
\item Find the velocity at $t = 6$, $t = 10$ and $t = 16$.
\item Sketch the three graphs $a$--$t$, $v$--$t$ and $x$--$t$, aligned vertically.
\item Find the total distance travelled in the $\SI{16}{s}$.
\end{enumerate}
\end{exercice}

\begin{corrige}[Exercise 1]
\textbf{1.} $v(0) = 0$. Phase 1: $\Delta v = 2 \times 6 = 12$, so $v(6) = \SI{12}{m.s^{-1}}$. Phase 2: $\Delta v = (-1)\times 4 = -4$, so $v(10) = 12 - 4 = \SI{8}{m.s^{-1}}$. Phase 3: $\Delta v = 0$, so $v(16) = \SI{8}{m.s^{-1}}$.

\textbf{2.} $a$--$t$: horizontal segments at $2$, $-1$, $0$. $v$--$t$: line from $(0,0)$ to $(6,12)$, line down to $(10,8)$, horizontal to $(16,8)$. $x$--$t$: parabola bending up to $(6,36)$, parabola bending down to $(10,76)$, straight line of gradient $8$ to $(16,124)$.

\textbf{3.} Distance $=$ area under $v$--$t$: triangle $\tfrac12(6)(12) = 36$; trapezium $\tfrac12(12+8)(4) = 40$; rectangle $8 \times 6 = 48$. Total $= 36 + 40 + 48 = \SI{124}{m}$.
\end{corrige}

\begin{aretenir}[Key points of Section 4]
\begin{itemize}
\item The $a$--$t$ graph shows how velocity changes; a horizontal line means constant acceleration, and $a = 0$ means \textbf{constant velocity} (not rest).
\item \textbf{Area under the $a$--$t$ graph $=$ change in velocity} $\Delta v$; always add the initial velocity: $v(t) = v(0) + \Delta v$.
\item \textbf{Gradient going down} the chain $x \to v \to a$; \textbf{area going up} the chain $a \to v \to x$, never forgetting the initial value.
\item Constant acceleration phases give: horizontal $a$--$t$ segments, straight $v$--$t$ segments, parabolic $x$--$t$ curves.
\item Consistency: a turning point on one graph corresponds to a zero crossing on the graph below it.
\item Exam tip: mark every phase-change instant on \textbf{all three} graphs first, then sketch one phase at a time, checking gradients against areas as you go.
\end{itemize}
\end{aretenir}

\coursec{Section 5: Free Fall and Motion under Gravity (Lessons 99, 100)}

In Section 4 we saw how the three graphs of motion — displacement–time, velocity–time and acceleration–time — are linked: the gradient of the $x$--$t$ graph gives velocity, the gradient of the $v$--$t$ graph gives acceleration, and areas under the $v$--$t$ and $a$--$t$ graphs give displacement and change of velocity. We now apply this whole machinery to the most important natural motion of all: the motion of a body moving under \cle{gravity} alone. Every stone dropped from a bridge over the River Wouri, every mango falling from a tree in Bamenda, every ball tossed into the air on a playground in Buea obeys the same simple laws — and these laws are a favourite of the GCE examiners.

\courssub{The Concept of Free Fall (Lesson 99)}

\textbf{An observation to start with.}\par
Drop a coin and a crumpled piece of paper from the same height at the same instant. They hit the ground together. Now drop the coin and an \emph{uncrumpled} sheet of paper: the paper flutters down much later. What changed? Not the mass of the paper — it is the same paper. What changed is the effect of \emph{air resistance}. This simple classroom experiment reveals a deep truth: when air resistance is removed or neglected, \textbf{all bodies fall with the same acceleration}, whatever their mass.

This was Galileo Galilei's great insight in the early seventeenth century, contradicting the ancient belief (due to Aristotle) that heavier bodies fall faster. Legend says Galileo dropped different masses from the Leaning Tower of Pisa; whether or not the story is literally true, his conclusion has been confirmed by countless experiments, most dramatically when astronaut David Scott dropped a hammer and a feather on the airless Moon in 1971 — they struck the ground simultaneously.

\begin{definition}[Free fall]
A body is in \cle{free fall} when the \emph{only} force acting on it is its weight, so that its only acceleration is the acceleration due to gravity $\vec{g}$, directed \textbf{vertically downward}. In free fall, air resistance and all other forces (such as buoyancy or tension) are neglected. Near the Earth's surface, the magnitude of $\vec{g}$ is approximately constant.
\end{definition}

\begin{propriete}[Acceleration due to gravity]
Near the surface of the Earth, every freely falling body has a constant downward acceleration of magnitude
\[
g \approx 9.8\ \mathrm{m\,s^{-2}} \qquad (\text{often taken as } 10\ \mathrm{m\,s^{-2}} \text{ in GCE exercises}).
\]
This acceleration is the \emph{same for all bodies}, whatever their mass, shape or material. Its direction is \textbf{always vertically downward}, even when the body is moving upward.
\end{propriete}

\begin{attention}[A heavier object does NOT fall faster]
A very common misconception is that a stone falls faster than a piece of chalk because it is heavier. In free fall, \textbf{mass does not appear anywhere} in the kinematics: the acceleration is $g$ for every body. The only reason everyday objects seem to fall at different rates is \emph{air resistance}, which we neglect in this chapter. In your answers, never write ``the heavier object lands first''.
\end{attention}

\begin{attention}[Exam tip — which value of $g$?]
GCE questions frequently state ``take $g = 10\ \mathrm{m\,s^{-2}}$''. If the question says so and you use $9.8$, you will obtain different numerical answers and \textbf{lose accuracy marks}. If nothing is said, $g = 9.8\ \mathrm{m\,s^{-2}}$ is the safe choice. Always read the question to the end before starting.
\end{attention}

\courssub{Applying the \emph{suvat} Equations to Vertical Motion}

Free fall is motion with \emph{uniform acceleration}, so all the equations of Section 2 apply unchanged. They can be derived from calculus: $a = \frac{dv}{dt} = -g$ (with up positive) $\Rightarrow v = u - gt$; $v = \frac{dx}{dt} \Rightarrow x = ut - \frac{1}{2}gt^2$; eliminating $t$ gives $v^2 = u^2 - 2gx$. The standard set is:
\[
v = u + at, \qquad s = ut + \tfrac{1}{2}at^2, \qquad v^2 = u^2 + 2as, \qquad s = \tfrac{1}{2}(u+v)t.
\]
The only new ingredient is that the acceleration is known in advance: it is $g$, downward. Because the motion is vertical, we must \emph{choose a positive direction} and stick to it.

\begin{methode}[Vertical \emph{suvat} checklist]
\begin{enumerate}
\item \textbf{Draw a diagram.} Sketch the motion, indicate the chosen positive direction (up or down), and mark the initial velocity $\vec{u}$ and the acceleration $\vec{g}$.
\item \textbf{Fix the positive direction.} Decide whether ``up'' or ``down'' is positive, and write it clearly on your script, e.g. ``take upward as positive''.
\item \textbf{Set the acceleration.} If up is positive, $a = -g$; if down is positive, $a = +g$. This is the single most common source of error.
\item \textbf{List the known quantities} $u, v, s, a, t$ with their correct signs, and identify the unknown.
\item \textbf{Choose the \emph{suvat} equation} containing the three knowns and the unknown, exactly as in Section 2.
\item \textbf{Solve and interpret} the sign of the answer (a negative velocity with ``up positive'' means the body is moving downward).
\end{enumerate}
\end{methode}

\begin{attention}[Sign of $g$ — the classic trap]
If you choose \textbf{upward as positive}, then the acceleration is $a = -9.8\ \mathrm{m\,s^{-2}}$, because gravity points \emph{down}. Writing $a = +9.8$ while taking up as positive is the most frequent mistake in this topic and it destroys the whole question. Remember: $g = 9.8\ \mathrm{m\,s^{-2}}$ is a \emph{magnitude} (always positive); the \emph{acceleration} is $\pm g$ depending on your chosen direction.
\end{attention}

\courssub{Body Dropped from Rest}

Suppose a stone is \emph{released} (not thrown) from a height $h$ above the ground. Take \textbf{downward as positive}. Then
\[
u = 0, \qquad a = +g, \qquad s = h.
\]
From $v^2 = u^2 + 2as$ we get $v^2 = 2gh$, and from $s = ut + \tfrac12 a t^2$ we get $h = \tfrac12 g t^2$. Hence:

\begin{propriete}[Body dropped from rest — proof examinable]
For a body dropped from rest at height $h$:
\[
\boxed{\,t = \sqrt{\dfrac{2h}{g}}\,} \quad \text{(time of fall)}, \qquad
\boxed{\,v = \sqrt{2gh}\,} \quad \text{(velocity on impact)}.
\]
Neither formula contains the mass of the body — Galileo's observation in algebraic form.
\end{propriete}

\begin{exemplebox}[Stone dropped from a bridge]
A stone is dropped from a bridge and hits the water \SI{3}{s} later. Taking $g = 10\ \mathrm{m\,s^{-2}}$ and down as positive: $h = \tfrac12 g t^2 = \tfrac12(10)(3^2) = \SI{45}{m}$, and the impact velocity is $v = gt = 10 \times 3 = \SI{30}{m.s^{-1}}$.
\end{exemplebox}

\courssub{Body Projected Vertically Upwards}

Now throw a ball straight up with initial speed $u$. Take \textbf{upward as positive}, so $a = -g$. The ball slows down on the way up, stops \emph{instantaneously} at the top, then falls back. The displacement $s$ is measured from the launch point (positive upward).

\textbf{Time to maximum height.} At the top, $v = 0$. Using $v = u + at$:
\[
0 = u - gt \quad\Longrightarrow\quad \boxed{\,t_{\text{up}} = \dfrac{u}{g}\,}.
\]

\textbf{Maximum height.} Using $v^2 = u^2 + 2as$ with $v = 0$, $a = -g$:
\[
0 = u^2 - 2gH \quad\Longrightarrow\quad \boxed{\,H = \dfrac{u^2}{2g}\,}.
\]

\begin{propriete}[At the top: $v = 0$ but $a \neq 0$]
At the highest point of the motion, the velocity is instantaneously zero, but the acceleration is still $g$ directed downward ($a = -g$ if up is positive). The body is \emph{not} in equilibrium at the top — gravity never switches off.
\end{propriete}

\begin{attention}[``Acceleration is zero at the top'' — wrong!]
Many candidates write that at maximum height ``$v = 0$ so $a = 0$''. This is false. If the acceleration were zero at the top, the ball would stay there forever! The velocity changes sign at the top precisely \emph{because} the acceleration is still $-g$ throughout the whole flight.
\end{attention}

\textbf{Total time of flight (back to launch level).} The displacement when the ball returns to its launch point is $s = 0$. Using $s = ut + \tfrac12 at^2$ with $a = -g$:
\[
0 = ut - \tfrac12 gt^2 = t(u - \tfrac12 gt) \quad\Longrightarrow\quad t = 0 \text{ (launch) or } t = \frac{2u}{g}.
\]
Hence the total time of flight is $\boxed{\,T = \dfrac{2u}{g}\,}$.

\textbf{Symmetry of the motion.} Consider the ball passing a given height on the way up and on the way down. From $v^2 = u^2 - 2gs$, the value of $v^2$ at a given height $s$ is fixed, so the \emph{speed} is the same at the same height — only the direction differs. It follows that the time to go up to the highest point equals the time to come back down from it, and the ball returns to its launch point with speed $u$ (moving downward).

\begin{popfigure}
\begin{tikzpicture}[scale=0.9]
% balcony
\fill[popblueL] (0,0) rectangle (2.2,0.25);
\draw[thick] (0,0) rectangle (2.2,0.25);
\draw[thick] (0,0) -- (0,-0.6); \draw[thick] (2.2,0) -- (2.2,-0.6);
\node[align=center, font=\small] at (1.1,-0.85) {balcony, Buea};
% ground
\draw[thick] (-0.5,-3.2) -- (3.2,-3.2);
\node[font=\small] at (2.9,-3.45) {ground};
% vertical path
\draw[dashed, popmuted] (1.1,0.25) -- (1.1,4.6);
% positions at equal time intervals
\foreach \y/\lab in {0.25/{}, 1.45/{}, 2.35/{}, 2.95/{}, 3.25/{}}{
  \fill[poporange] (1.1,\y) circle (0.09);
}
\node[font=\small, right] at (1.25,3.25) {top: $v=0$, $a=-g$};
% velocity arrows up
\draw[-latex, very thick, popblue] (0.55,0.25) -- (0.55,1.55) node[midway, left, font=\small] {$u$};
\draw[-latex, thick, popblue] (0.55,1.7) -- (0.55,2.5);
\draw[-latex, thick, popblue] (0.55,2.6) -- (0.55,3.0);
% velocity arrows down
\draw[-latex, thick, poppink] (1.65,3.0) -- (1.65,2.5);
\draw[-latex, thick, poppink] (1.65,2.35) -- (1.65,1.55);
\draw[-latex, very thick, poppink] (1.65,1.4) -- (1.65,0.3) node[midway, right, font=\small] {$v$};
% g arrow
\draw[-latex, very thick, popgreen] (2.6,2.6) -- (2.6,1.6) node[midway, right, font=\small] {$\vec{g}$};
% height label
\draw[latex-latex, popdark] (3.0,0.25) -- (3.0,3.25) node[midway, right, font=\small] {$H=\dfrac{u^2}{2g}$};
\end{tikzpicture}
\legende{A ball thrown vertically upward from a balcony in Buea, shown at equal time intervals. The spacing decreases on the way up (deceleration) and increases on the way down (acceleration); the acceleration $\vec{g}$ is downward throughout.}
\end{popfigure}

\begin{exoresolu}[Ball thrown upward from a balcony]
A ball is thrown vertically upward with speed $\SI{14}{m.s^{-1}}$ from a balcony $\SI{15}{m}$ above the ground. Take $g = 10\ \mathrm{m\,s^{-2}}$ and upward as positive. Find: (a) the maximum height above the ground; (b) the total time until the ball hits the ground; (c) the speed of impact.
\tcblower
\textbf{(a) Maximum height above launch point:}
\[
H = \frac{u^2}{2g} = \frac{14^2}{2 \times 10} = \frac{196}{20} = \SI{9.8}{m}.
\]
Maximum height above the ground: $15 + 9.8 = \SI{24.8}{m}$.

\textbf{(b) Total time of flight.} Use $s = ut + \tfrac12 a t^2$ with $s = -15\ \mathrm{m}$ (the ball finishes $15$ m \emph{below} the launch point, and up is positive), $u = 14$, $a = -10$:
\[
-15 = 14t - 5t^2 \quad\Longrightarrow\quad 5t^2 - 14t - 15 = 0.
\]
By the quadratic formula:
\[
t = \frac{14 \pm \sqrt{196 + 300}}{10} = \frac{14 \pm \sqrt{496}}{10} \approx \frac{14 \pm 22.27}{10}.
\]
Taking the positive root: $t \approx \SI{3.63}{s}$. (The negative root $t \approx -0.83$ s is rejected: time cannot be negative.)

\textbf{(c) Impact velocity.} $v = u + at = 14 - 10(3.63) \approx \SI{-22.3}{m.s^{-1}}$. The negative sign means the ball is moving \emph{downward}; the impact speed is about $\SI{22.3}{m.s^{-1}}$.
\end{exoresolu}

\courssub{Graphical Representation of Motion under Gravity (Lesson 100)}

Everything we learned in Sections 3 and 4 about graphs now pays off. For a ball projected vertically upward with initial speed $u$, taking \textbf{upward as positive}:
\[
a(t) = -g, \qquad v(t) = u - gt, \qquad x(t) = ut - \tfrac12 g t^2.
\]

\begin{itemize}
\item The \textbf{acceleration--time graph} is a horizontal straight line at $-g$: the acceleration is constant and downward for the \emph{whole} flight, including at the top.
\item The \textbf{velocity--time graph} is a straight line of gradient $-g$, starting at $v = u$ and crossing the time axis at $t = u/g$ (the instant the ball is at the top). The area between the graph and the time axis (counted negative below the axis) gives the displacement.
\item The \textbf{displacement--time graph} is a \emph{parabola} opening downward, with maximum $H = u^2/(2g)$ at $t = u/g$. Its gradient at any point is the velocity: positive on the way up, zero at the top, negative on the way down.
\end{itemize}

\begin{popfigure}
\begin{center}
\begin{tikzpicture}
\begin{axis}[
  name=xt, width=7.2cm, height=4.6cm,
  xlabel={$t$ (s)}, ylabel={$x$ (m)},
  xmin=0, xmax=3, ymin=-2, ymax=12,
  xtick={0,1,2,3}, ytick={0,5,10},
  grid=both, grid style={popline!40},
  title={$x$--$t$: parabola}, title style={font=\small},
  label style={font=\small}, ticklabel style={font=\small},
]
\addplot[very thick, popblue, domain=0:2.8, samples=60] {10*x - 5*x^2};
\addplot[mark=*, poporange, mark size=1.5pt] coordinates {(1,5)};
\node[font=\small, popdark] at (axis cs:1.6,6.5) {top $(1,\,5)$};
\end{axis}
\begin{axis}[
  name=vt, at={(xt.right of origin)}, anchor=left of origin, xshift=1.1cm,
  width=7.2cm, height=4.6cm,
  xlabel={$t$ (s)}, ylabel={$v$ (m\,s$^{-1}$)},
  xmin=0, xmax=3, ymin=-15, ymax=12,
  xtick={0,1,2,3}, ytick={-10,0,10},
  grid=both, grid style={popline!40},
  title={$v$--$t$: line of gradient $-g$}, title style={font=\small},
  label style={font=\small}, ticklabel style={font=\small},
]
\addplot[very thick, poppink, domain=0:2.8, samples=10] {10 - 10*x};
\draw[dashed, popmuted] (axis cs:0,0) -- (axis cs:3,0);
\end{axis}
\begin{axis}[
  name=at, at={(vt.right of origin)}, anchor=left of origin, xshift=1.1cm,
  width=7.2cm, height=4.6cm,
  xlabel={$t$ (s)}, ylabel={$a$ (m\,s$^{-2}$)},
  xmin=0, xmax=3, ymin=-15, ymax=12,
  xtick={0,1,2,3}, ytick={-10,0,10},
  grid=both, grid style={popline!40},
  title={$a$--$t$: constant $-g$}, title style={font=\small},
  label style={font=\small}, ticklabel style={font=\small},
]
\addplot[very thick, popgreen, domain=0:2.8, samples=5] {-10};
\draw[dashed, popmuted] (axis cs:0,0) -- (axis cs:3,0);
\end{axis}
\end{tikzpicture}
\end{center}
\legende{The three graphs for a ball projected upward at $\SI{10}{m.s^{-1}}$ with $g = 10\ \mathrm{m\,s^{-2}}$, upward positive. The $v$--$t$ line has gradient $-10$ and crosses the axis at the top of the flight; the $a$--$t$ line stays at $-g$ even at the highest point. The $x$--$t$ parabola crosses the time axis at $t=2$ s (return to launch level).}
\end{popfigure}

Note how the three graphs are perfectly consistent with Section 4: the gradient of the parabola at $t = 0$ is $u$ (the initial slope of the $x$--$t$ curve equals the initial velocity), the gradient of the $v$--$t$ line equals the constant value of the $a$--$t$ graph, and the area under the $v$--$t$ graph from $0$ to $2u/g$ is zero — the ball is back where it started.

\courssub{Dropped, Thrown Down, Thrown Up: a Comparison}

A classic GCE question compares three bodies released from the \emph{same} height $h$: one \emph{dropped}, one \emph{thrown downward} with speed $u_0$, one \emph{thrown upward} with speed $u_0$. Which lands first? Which lands fastest?

\begin{exoresolu}[Three balls from the same height]
From a point $\SI{20}{m}$ above the ground, ball A is \emph{dropped}, ball B is \emph{thrown vertically downward} at $\SI{5}{m.s^{-1}}$, and ball C is \emph{thrown vertically upward} at $\SI{5}{m.s^{-1}}$. Take $g = 10\ \mathrm{m\,s^{-2}}$. Find the time of flight and impact speed of each ball.
\tcblower
Take \textbf{downward as positive}, so $a = +10$ and $s = +20$ for all three balls.

\textbf{Ball A (dropped):} $u = 0$.
\[
t_A = \sqrt{\frac{2h}{g}} = \sqrt{\frac{2 \times 20}{10}} = \SI{2}{s}, \qquad
v_A = \sqrt{2gh} = \sqrt{400} = \SI{20}{m.s^{-1}}.
\]

\textbf{Ball B (thrown down):} $u = +5$. From $s = ut + \tfrac12 a t^2$:
\[
20 = 5t + 5t^2 \;\Longrightarrow\; t^2 + t - 4 = 0 \;\Longrightarrow\; t_B = \frac{-1 + \sqrt{17}}{2} \approx \SI{1.56}{s}.
\]
\[
v_B^2 = u^2 + 2gh = 25 + 400 = 425 \;\Longrightarrow\; v_B \approx \SI{20.6}{m.s^{-1}}.
\]

\textbf{Ball C (thrown up):} $u = -5$ (upward is negative here). The same equation gives
\[
20 = -5t + 5t^2 \;\Longrightarrow\; t^2 - t - 4 = 0 \;\Longrightarrow\; t_C = \frac{1 + \sqrt{17}}{2} \approx \SI{2.56}{s}.
\]
\[
v_C^2 = u^2 + 2gh = 25 + 400 = 425 \;\Longrightarrow\; v_C \approx \SI{20.6}{m.s^{-1}}.
\]

\textbf{Comparison.} B lands first ($1.56$ s), then A ($2$ s), then C ($2.56$ s). But B and C hit the ground with the \emph{same speed} ($20.6$ m\,s$^{-1}$): by symmetry, C passes its launch point moving downward at exactly $\SI{5}{m.s^{-1}}$ — the same state as B at launch — so from that instant on, B and C have identical motion. Note also $t_C - t_B = 1\ \mathrm{s} = \dfrac{2u_0}{g}$, exactly the time C spends going up and coming back to the launch level.
\end{exoresolu}

\begin{aretenir}[Key points of Section 5]
\begin{itemize}
\item Free fall: the only acceleration is $g \approx \SI{9.8}{m.s^{-2}}$, always directed \emph{vertically downward}; it is the same for all bodies regardless of mass.
\item Fix a positive direction first: up positive $\Rightarrow a = -g$; down positive $\Rightarrow a = +g$.
\item Dropped from rest: $t = \sqrt{2h/g}$, $v = \sqrt{2gh}$.
\item Projected upward: $t_{\text{up}} = u/g$, $H = u^2/(2g)$; at the top $v = 0$ but $a = -g \neq 0$; total flight time back to launch level is $2u/g$.
\item Up--down symmetry: same speed at the same height; time up = time down.
\item Graphs (up positive): $x$--$t$ is a downward parabola, $v$--$t$ a line of gradient $-g$, $a$--$t$ a horizontal line at $-g$.
\item Check the question: use $g = 10\ \mathrm{m\,s^{-2}}$ if told to.
\end{itemize}
\end{aretenir}

\begin{exercice}[Practice — free fall]
Take $g = 10\ \mathrm{m\,s^{-2}}$ unless stated otherwise.
\begin{enumerate}
\item A mango falls from a branch $\SI{5}{m}$ above the ground. Find the time of fall and the speed on impact.
\item A stone is thrown vertically upward at $\SI{20}{m.s^{-1}}$ from ground level. Find (a) the maximum height, (b) the total time of flight, (c) its velocity after $\SI{3}{s}$, and interpret the sign.
\item A ball is dropped from a tower and takes $\SI{4}{s}$ to reach the ground. Find the height of the tower and the impact speed.
\item From the top of a $\SI{45}{m}$ cliff, a stone is thrown vertically \emph{downward} at $\SI{10}{m.s^{-1}}$. Find the time to reach the foot of the cliff and the impact speed.
\item A ball is projected vertically upward from the ground with speed $u$. Show that the times at which it is at height $h$ (with $h < u^2/2g$) are the roots of $gt^2 - 2ut + 2h = 0$, and deduce that the sum of these two times is the total time of flight.
\end{enumerate}
\end{exercice}

\begin{corrige}[Exercise 5]
\textbf{1.} $t = \sqrt{2h/g} = \sqrt{10/10} = \SI{1}{s}$; $v = \sqrt{2gh} = \sqrt{100} = \SI{10}{m.s^{-1}}$.

\textbf{2.} Up positive, $a = -10$. (a) $H = u^2/2g = 400/20 = \SI{20}{m}$. (b) $T = 2u/g = \SI{4}{s}$. (c) $v = u - gt = 20 - 30 = \SI{-10}{m.s^{-1}}$: the stone is moving \emph{downward} at $\SI{10}{m.s^{-1}}$ (it is on its way back down).

\textbf{3.} $h = \tfrac12 g t^2 = \tfrac12(10)(16) = \SI{80}{m}$; $v = gt = \SI{40}{m.s^{-1}}$.

\textbf{4.} Down positive: $45 = 10t + 5t^2 \Rightarrow t^2 + 2t - 9 = 0 \Rightarrow t = \dfrac{-2+\sqrt{40}}{2} \approx \SI{2.16}{s}$. $v^2 = 100 + 2(10)(45) = 1000 \Rightarrow v \approx \SI{31.6}{m.s^{-1}}$.

\textbf{5.} With up positive, $h = ut - \tfrac12 g t^2$, i.e. $gt^2 - 2ut + 2h = 0$ after multiplying by 2 and rearranging. The two roots $t_1, t_2$ are the times of passing height $h$ (up and down). By Viète's formula, $t_1 + t_2 = \dfrac{2u}{g}$, which is exactly the total time of flight — a beautiful confirmation of the up--down symmetry.
\end{corrige}

\begin{savaistu}[Did you know?]
The value of $g$ is not exactly the same everywhere on Earth: it is slightly larger at the poles (about $\SI{9.83}{m.s^{-2}}$) and slightly smaller at the equator (about $\SI{9.78}{m.s^{-2}}$), because the Earth is not a perfect sphere and because of its rotation. Mount Cameroon, rising about $\SI{4100}{m}$ above sea level, experiences a value of $g$ marginally smaller than at the coast of Limbe below it — the difference is tiny, but modern instruments called \emph{gravimeters} can detect it, and geophysicists use such variations to prospect for oil and minerals.
\end{savaistu}

\coursec{Section 6: Real-Life Applications and Synthesis (Lessons 94 \& 101, consolidation)}

In Section 5 we studied free fall and motion under gravity: we saw that every object near the Earth's surface, when air resistance is neglected, accelerates downward at $g \approx 9.81\ \mathrm{m\,s^{-2}}$, and we learned to read this motion on displacement--time and velocity--time graphs. We now close the chapter by putting \emph{all} of our kinematics to work on real situations: stopping a car on the Douala--Yaound\'e highway, estimating the height of a building with a stopwatch, predicting where two moving objects meet, and reading a full journey from a verbal description. This is exactly the style of the long, multi-part questions set at the GCE Advanced Level.

\courssub{Road Safety: Thinking Distance, Braking Distance, Stopping Distance}

Imagine you are driving at a steady speed $v$ when a goat suddenly crosses the road. Two things happen before the car stops. First, there is a \emph{reaction time} $t_r$ (typically $0.5$ to $1\ \mathrm{s}$) during which you see the danger, decide to brake, and press the pedal. During this time the car is still moving at full speed $v$. Second, once the brakes bite, the car decelerates uniformly with acceleration $a$ (negative) until it stops.

\begin{definition}[Thinking, braking and stopping distances]
The \cle{thinking distance} is the distance travelled at constant speed during the reaction time: $d_{\text{think}} = v\,t_r$. The \cle{braking distance} is the distance travelled while decelerating uniformly from speed $v$ to rest. The \cle{stopping distance} is the sum:
\[
d_{\text{stop}} = d_{\text{think}} + d_{\text{brake}}.
\]
\end{definition}

\begin{propriete}[Stopping-distance formula]
For a vehicle moving at speed $v$, with reaction time $t_r$ and braking deceleration of magnitude $|a|$,
\[
d_{\text{stop}} = v\,t_r + \frac{v^{2}}{2|a|}.
\]
In particular the braking distance grows with the \textbf{square} of the speed: doubling the speed multiplies the braking distance by four.
\end{propriete}

\textbf{Derivation (examinable).} During the reaction time the velocity is constant, so $d_{\text{think}} = v t_r$ immediately. During braking the acceleration is constant, so we may use $v_f^{2} = v_i^{2} + 2as$ with final velocity $v_f = 0$, initial velocity $v_i = v$, and acceleration $a = -|a|$:
\[
0 = v^{2} - 2|a|\,s \quad\Longrightarrow\quad s = \frac{v^{2}}{2|a|}. \qquad \blacksquare
\]

\begin{popfigure}
\begin{center}
\begin{tikzpicture}[scale=0.9]
\draw[thick, popline] (0,0) -- (12.6,0);
\draw[fill=popblueL, draw=popblue] (0,0.15) rectangle (1.1,0.75);
\node[popink] at (0.55,0.45) {\small car};
\draw[->, thick, popblue] (1.1,0.45) -- (2.1,0.45) node[midway, above, popink]{\small $v$};
\draw[fill=poporange!40, draw=poporange] (11.6,0.1) rectangle (12.4,0.9);
\node[popink, align=center] at (12.0,0.5) {\small goat};
\draw[<->, thick, popgreen] (0,-0.5) -- (4,-0.5);
\node[popgreen, below] at (2,-0.5) {\small thinking $= v t_r$};
\draw[<->, thick, poppurple] (4,-0.5) -- (11.6,-0.5);
\node[poppurple, below] at (7.8,-0.5) {\small braking $= \dfrac{v^{2}}{2|a|}$};
\draw[<->, thick, popdark] (0,-1.4) -- (11.6,-1.4);
\node[popdark, below] at (5.8,-1.4) {\small stopping distance};
\draw[dashed, popmuted] (4,0) -- (4,-1.5);
\draw[dashed, popmuted] (11.6,0) -- (11.6,-1.5);
\end{tikzpicture}
\end{center}
\legende{Stopping distance = thinking distance + braking distance.}
\end{popfigure}

\begin{exemplebox}[Stopping distances on the Douala--Yaound\'e highway]
Take a reaction time $t_r = 0.7\ \mathrm{s}$ and a braking deceleration $|a| = 6\ \mathrm{m\,s^{-2}}$ (good tyres, dry road). Convert speeds: $50\ \mathrm{km/h} = 13.9\ \mathrm{m\,s^{-1}}$, $80\ \mathrm{km/h} = 22.2\ \mathrm{m\,s^{-1}}$, $110\ \mathrm{km/h} = 30.6\ \mathrm{m\,s^{-1}}$.
\begin{center}
\begin{tabularx}{\linewidth}{|l|X|X|X|}
\hline
\rowcolor{popblueL} Speed & Thinking $v t_r$ & Braking $v^2/(2|a|)$ & Stopping \\
\hline
$50\ \mathrm{km/h}$ & $9.7\ \mathrm{m}$ & $16.1\ \mathrm{m}$ & $25.8\ \mathrm{m}$ \\
\hline
$80\ \mathrm{km/h}$ & $15.6\ \mathrm{m}$ & $41.1\ \mathrm{m}$ & $56.7\ \mathrm{m}$ \\
\hline
$110\ \mathrm{km/h}$ & $21.4\ \mathrm{m}$ & $78.0\ \mathrm{m}$ & $99.4\ \mathrm{m}$ \\
\hline
\end{tabularx}
\end{center}
Going from $50$ to $100\ \mathrm{km/h}$ (doubling the speed) would multiply the braking part by $4$; here, from $50$ to $110\ \mathrm{km/h}$ the braking distance is multiplied by about $4.8$. Speed kills through the \emph{square} in the formula.
\end{exemplebox}

\begin{popfigure}
\begin{center}
\begin{tikzpicture}
\begin{axis}[
    ybar, bar width=0.7cm,
    width=10cm, height=6cm,
    symbolic x coords={50, 80, 110},
    xtick=data,
    xlabel={Speed (km/h)}, ylabel={Braking distance (m)},
    ymin=0, ymax=90,
    nodes near coords,
    every node near coord/.append style={font=\small, color=popink},
    ymajorgrids, grid style={popline!40},
    axis line style={popline},
    tick label style={color=popink},
    label style={color=popink},
]
\addplot[fill=popblue, draw=popdark] coordinates {(50,16.1) (80,41.1) (110,78.0)};
\end{axis}
\end{tikzpicture}
\end{center}
\legende{Braking distance at 50, 80 and $110\ \mathrm{km/h}$ for $|a| = 6\ \mathrm{m\,s^{-2}}$.}
\end{popfigure}

\begin{attention}[Forgetting the thinking distance]
A very common GCE error is to compute only $v^{2}/(2|a|)$ and call it the stopping distance. Unless the question says the driver reacts instantly, you \textbf{must} add $v\,t_r$. Read the question twice and underline ``reaction time''.
\end{attention}

\courssub{Falling Objects in Daily Life: Measuring Heights with a Stopwatch}

Drop a stone into a well and time its fall. Since the stone starts from rest and falls freely, $h = \tfrac{1}{2} g t^{2}$, so measuring $t$ gives the depth $h$. This is a genuine field technique.

\begin{exoresolu}[Depth of a well]
A stone dropped from rest into a well is heard to splash after $2.4\ \mathrm{s}$. Estimate the depth of the well, taking $g = 9.81\ \mathrm{m\,s^{-2}}$ and neglecting the time sound takes to come back up.
\tcblower
\textbf{Solution.} The stone starts from rest ($u=0$) with acceleration $g$:
\[
h = \tfrac{1}{2} g t^{2} = \tfrac{1}{2}(9.81)(2.4)^{2} = \tfrac{1}{2}(9.81)(5.76) \approx 28.3\ \mathrm{m}.
\]
The well is about $28\ \mathrm{m}$ deep. (Including the sound travel time would reduce the answer slightly.)
\end{exoresolu}

The same idea estimates the height of a storey building: an object taking $1.0\ \mathrm{s}$ to fall covers $\tfrac{1}{2}(9.81)(1)^2 \approx 4.9\ \mathrm{m}$, roughly one and a half storeys.

\courssub{Vertical Motion in Sport and Construction}

A ball kicked vertically upward, or a spanner slipped from scaffolding, is motion under uniform acceleration $-g$ (taking upward as positive). The full toolkit of Section 5 applies: time to the top is $u/g$, maximum height is $u^{2}/(2g)$, and the speed of return to the launch level equals the launch speed.

\begin{exoresolu}[A tool dropped from scaffolding]
A mason drops a hammer from scaffolding $15\ \mathrm{m}$ above the ground. A colleague stands below. (a) How long does the hammer take to reach the ground? (b) With what speed does it arrive? (c) The colleague needs $0.8\ \mathrm{s}$ to step aside after a warning shout; is a shout at the moment of release sufficient? Take $g = 9.81\ \mathrm{m\,s^{-2}}$.
\tcblower
\textbf{Solution.} Take downward positive, $u=0$, $a = g$, $s = 15\ \mathrm{m}$.
(a) $s = \tfrac{1}{2} g t^{2} \Rightarrow t = \sqrt{2s/g} = \sqrt{30/9.81} \approx 1.75\ \mathrm{s}$.
(b) $v = gt = 9.81 \times 1.75 \approx 17.2\ \mathrm{m\,s^{-1}}$ (about $62\ \mathrm{km/h}$!).
(c) $1.75\ \mathrm{s} > 0.8\ \mathrm{s}$, so the colleague has just under a second to spare --- but a shout even $0.5\ \mathrm{s}$ late would leave no safety margin. This is why hard hats and exclusion zones matter on construction sites in Douala and Yaound\'e.
\end{exoresolu}

\courssub{Two-Particle Problems: The Meeting Condition}

Many GCE questions involve two objects moving together: one dropped from a height, another projected upward from the ground. The whole secret is organisation.

\begin{methode}[Solving meeting problems]
\begin{enumerate}
\item Choose \textbf{one} origin and \textbf{one} positive direction for both particles.
\item Start \textbf{one} clock: $t=0$ at the same instant for both particles (if one starts later, shift its time, e.g.\ replace $t$ by $t-2$).
\item Write the position of each particle: $x_1(t)$ and $x_2(t)$, using $x = x_0 + ut + \tfrac{1}{2}at^{2}$.
\item The meeting condition is $x_1(t) = x_2(t)$. Solve for $t$, then substitute back to find where they meet.
\item Check that the answer is physically sensible (positive time, position within the allowed region).
\end{enumerate}
\end{methode}

\begin{attention}[Different origins or different clocks]
The classic trap is to measure the dropped stone's position from the top and the thrown ball's position from the ground, then equate them --- this compares distances measured from different origins and gives nonsense. Similarly, using $t=0$ at the drop for one particle and at the throw for the other corrupts everything. One origin, one clock.
\end{attention}

\begin{exoresolu}[Drop and throw]
A stone is dropped from rest from the top of a tower $45\ \mathrm{m}$ high. At the same instant, a ball is thrown vertically upward from the foot of the tower with speed $30\ \mathrm{m\,s^{-1}}$. Taking $g = 10\ \mathrm{m\,s^{-2}}$, find when and where they meet.
\tcblower
\textbf{Solution.} Origin at the ground, upward positive, $t=0$ at the instant of release.
Stone: $x_0 = 45$, $u = 0$, $a = -10$:
\[
x_1(t) = 45 - 5t^{2}.
\]
Ball: $x_0 = 0$, $u = 30$, $a = -10$:
\[
x_2(t) = 30t - 5t^{2}.
\]
Meeting: $x_1 = x_2 \Rightarrow 45 - 5t^{2} = 30t - 5t^{2} \Rightarrow 45 = 30t \Rightarrow t = 1.5\ \mathrm{s}$.
Position: $x_2(1.5) = 30(1.5) - 5(1.5)^{2} = 45 - 11.25 = 33.75\ \mathrm{m}$ above the ground.
Check: the ball is still rising at $t=1.5\ \mathrm{s}$ (its top is at $t = u/g = 3\ \mathrm{s}$), and the stone has fallen only $5(1.5)^2 = 11.25\ \mathrm{m}$: consistent.
\end{exoresolu}

\begin{popfigure}
\begin{center}
\begin{tikzpicture}[scale=0.85]
\draw[thick, popline] (0,0) -- (0,4.5);
\node[popink, left] at (0,4.5) {\small $45$ m};
\node[popink, left] at (0,0) {\small ground};
\draw[fill=popblue, draw=popdark] (0,4.5) circle (0.09);
\draw[->, thick, popblue] (0.15,4.4) .. controls (0.7,3.6) and (0.7,2.6) .. (0.5,1.6);
\node[popblue, right, align=center] at (0.75,3.4) {\small stone dropped\\ \small $x_1 = 45 - 5t^2$};
\draw[fill=poporange, draw=popdark] (2.2,0) circle (0.09);
\draw[->, thick, poporange] (2.2,0.15) .. controls (2.3,1.2) and (2.0,2.2) .. (1.35,3.3);
\node[poporange, right, align=center] at (2.35,1.6) {\small ball thrown\\ \small $x_2 = 30t - 5t^2$};
\draw[fill=poppurple, draw=popdark] (1.05,3.38) circle (0.12);
\node[poppurple, right] at (1.25,3.38) {\small meet: $t = 1.5$ s, $x = 33.75$ m};
\draw[dashed, popmuted] (1.05,3.38) -- (0,3.38);
\node[popink, left] at (0,3.38) {\small $33.75$ m};
\end{tikzpicture}
\end{center}
\legende{The drop-and-throw meeting problem: one origin (ground), one clock.}
\end{popfigure}

\courssub{Graphical Synthesis: From a Verbal Description to All Three Graphs}

Long GCE questions often describe a journey in words and ask for the displacement--time, velocity--time and acceleration--time graphs, plus numerical answers. The strategy: cut the journey into phases of constant acceleration; compute the velocity at the end of each phase; remember that the \emph{area} under the $v$--$t$ graph gives displacement, and the \emph{gradient} of the $v$--$t$ graph gives acceleration.

\begin{exoresolu}[A complete journey]
A bus starts from rest at a stop, accelerates uniformly at $1.5\ \mathrm{m\,s^{-2}}$ for $8\ \mathrm{s}$, travels at constant velocity for $20\ \mathrm{s}$, then brakes uniformly to rest in $6\ \mathrm{s}$. Find (a) the maximum velocity, (b) the total distance travelled, (c) the deceleration during braking, and describe the three graphs.
\tcblower
\textbf{Solution.}
(a) $v_{\max} = u + at = 0 + 1.5 \times 8 = 12\ \mathrm{m\,s^{-1}}$.
(b) Use areas under the $v$--$t$ graph (triangle + rectangle + triangle):
\[
s = \tfrac{1}{2}(8)(12) + (20)(12) + \tfrac{1}{2}(6)(12) = 48 + 240 + 36 = 324\ \mathrm{m}.
\]
(c) $a = \dfrac{0 - 12}{6} = -2\ \mathrm{m\,s^{-2}}$, a deceleration of $2\ \mathrm{m\,s^{-2}}$.
Graphs: the $v$--$t$ graph is a trapezium rising from $0$ to $12$ over $[0,8]$, flat over $[8,28]$, falling to $0$ over $[28,34]$. The $a$--$t$ graph is $+1.5$, then $0$, then $-2\ \mathrm{m\,s^{-2}}$ on the same intervals. The $x$--$t$ graph is a curve bending upward (parabola) for $8\ \mathrm{s}$, then a straight line of slope $12$, then a curve flattening to a horizontal tangent at $t = 34\ \mathrm{s}$.
\end{exoresolu}

\courssub{Chapter Synthesis and GCE Strategy}

\begin{aretenir}[Chapter summary map]
\begin{itemize}
\item \textbf{Language:} displacement $s$, velocity $v = \dot{s}$, acceleration $a = \dot{v}$; scalars vs vectors.
\item \textbf{Uniform acceleration:} $v = u + at$,\; $s = ut + \tfrac12 at^{2}$,\; $s = \tfrac12(u+v)t$,\; $v^{2} = u^{2} + 2as$.
\item \textbf{Graphs:} gradient of $x$--$t$ = velocity; gradient of $v$--$t$ = acceleration; area under $v$--$t$ = displacement; area under $a$--$t$ = change of velocity.
\item \textbf{Gravity:} $a = -g \approx -9.81\ \mathrm{m\,s^{-2}}$ (upward positive); time to top $u/g$; max height $u^{2}/(2g)$.
\item \textbf{Applications:} stopping distance $v t_r + v^{2}/(2|a|)$; meeting problems via $x_1(t) = x_2(t)$ with one origin and one clock.
\end{itemize}
\end{aretenir}

\textbf{Common GCE question patterns.} (i) A two-phase journey with graph sketching and area calculations. (ii) A vertical projection with ``greatest height'' and ``time of flight'' parts. (iii) Two particles (drop-and-throw, or overtaking) with a meeting condition. (iv) A stopping-distance question with reaction time. In all of these, later parts reuse earlier results.

\begin{attention}[Exam tip: keep intermediate values exact]
In long questions, part (b) almost always uses the answer of part (a). Keep intermediate values \textbf{exact} (fractions, surds, or unrounded decimals in your calculator memory) and only round at the final line. Rounding $v_{\max} = 12$ early is safe, but rounding $\sqrt{30/9.81}$ to $1.7$ before squaring it again can cost accuracy marks.
\end{attention}

\begin{savaistu}[Why real falling objects do not accelerate forever]
Our model ignores air resistance. In reality, air drag grows with speed until it balances weight; the object then falls at a constant \cle{terminal velocity}. A skydiver in free fall reaches roughly $190$--$200\ \mathrm{km/h}$ before opening the parachute; a raindrop only a few metres per second --- which is why rain does not hurt us. At the GCE level you only need the concept: drag explains why the uniform-acceleration model is an approximation, excellent for dense objects over short falls.
\end{savaistu}

\begin{exercice}[Consolidation]
\begin{enumerate}
\item A car travels at $90\ \mathrm{km/h}$ on a wet road where the braking deceleration is only $4\ \mathrm{m\,s^{-2}}$; the driver's reaction time is $0.9\ \mathrm{s}$. Find the stopping distance.
\item A mango falls from a tree and takes $1.1\ \mathrm{s}$ to reach the ground. How high was it? With what speed does it land? ($g = 9.81\ \mathrm{m\,s^{-2}}$)
\item A ball is thrown vertically upward at $20\ \mathrm{m\,s^{-1}}$ from the ground at the same instant as a stone is dropped from a $60\ \mathrm{m}$ tower beside it. Taking $g = 10\ \mathrm{m\,s^{-2}}$, find when and where they meet.
\item A lorry accelerates from rest at $0.8\ \mathrm{m\,s^{-2}}$ for $15\ \mathrm{s}$, keeps constant speed for $40\ \mathrm{s}$, then decelerates uniformly to rest over $100\ \mathrm{m}$. Find the total distance and the total time.
\end{enumerate}
\end{exercice}

\begin{corrige}[Exercise 1]
\textbf{1.} $v = 90/3.6 = 25\ \mathrm{m\,s^{-1}}$. Thinking: $25 \times 0.9 = 22.5\ \mathrm{m}$. Braking: $25^{2}/(2 \times 4) = 78.1\ \mathrm{m}$. Stopping $\approx 100.6\ \mathrm{m}$.
\textbf{2.} $h = \tfrac12(9.81)(1.1)^{2} \approx 5.9\ \mathrm{m}$; $v = 9.81 \times 1.1 \approx 10.8\ \mathrm{m\,s^{-1}}$.
\textbf{3.} $x_1 = 60 - 5t^{2}$, $x_2 = 20t - 5t^{2}$; $x_1 = x_2 \Rightarrow t = 3\ \mathrm{s}$, $x = 20(3) - 5(9) = 15\ \mathrm{m}$ above the ground.
\textbf{4.} $v = 0.8 \times 15 = 12\ \mathrm{m\,s^{-1}}$. Distances: $\tfrac12(15)(12) = 90$; $40 \times 12 = 480$; braking $100\ \mathrm{m}$ given. Total $= 670\ \mathrm{m}$. Braking time: $100 = \tfrac12(12+0)t \Rightarrow t = 16.7\ \mathrm{s}$. Total time $\approx 71.7\ \mathrm{s}$.
\end{corrige}

\newpage

\coursec{Integration activity}

\coursec{Real-Life Applications and Synthesis}

\courssub{Aims of the Activity}

This integration activity brings together everything you have learnt in this chapter: the \cle{suvat equations}, the interpretation and construction of \cle{displacement--time}, \cle{velocity--time} and \cle{acceleration--time} graphs, and \cle{free fall} under gravity. Working in groups, you will act as road-safety engineers studying a real emergency stop in Bamenda, and you will communicate your findings on a poster, exactly as an engineer presents a technical report.

By the end of the activity you should be able to:
\begin{itemize}
\item convert speeds between $\mathrm{km\,h^{-1}}$ and $\mathrm{m\,s^{-1}}$ fluently;
\item decompose a stopping distance into \cle{thinking distance} and \cle{braking distance};
\item use $v^2 = u^2 + 2as$ and the \emph{area under a velocity--time graph} to find distances, and check that both methods agree;
\item sketch coherent $v$--$t$, $x$--$t$ and $a$--$t$ graphs for the same motion;
\item model a falling object with $s = \tfrac{1}{2}gt^2$ and combine it with horizontal motion at constant velocity;
\item compare stopping distances at different speeds and write a justified safety recommendation.
\end{itemize}

\courssub{Situation}

\textbf{The Bend-Skin Safety Report.}\par
In Bamenda, \emph{bend-skins} (motorbike taxis) are one of the most common means of transport between Commercial Avenue and Nkwen. The Students' Road Safety Club of your school has been asked by the local council to prepare a short technical report on emergency braking.

A bend-skin is travelling at a steady $60\ \mathrm{km\,h^{-1}}$ along a straight, level, tarred road when a goat suddenly crosses $40\ \mathrm{m}$ ahead. The rider's \cle{reaction time} (the time between seeing the danger and the brakes starting to act) is $0.8\ \mathrm{s}$. Once applied, the brakes produce a constant deceleration of $4\ \mathrm{m\,s^{-2}}$ until the bike stops. Take $g = 10\ \mathrm{m\,s^{-2}}$.

During the braking, the rider's phone slips from his pocket at a height of $1.2\ \mathrm{m}$ above the ground. The phone's \emph{vertical} motion may be treated as free fall from rest vertically, while the bike continues horizontally.

Your group must produce a poster containing: the numerical results, the three motion graphs ($v$--$t$, $x$--$t$, $a$--$t$), and a one-paragraph safety recommendation comparing stopping distances at $40$, $60$ and $80\ \mathrm{km\,h^{-1}}$.

\courssub{Tasks}

\begin{enumerate}
\item \textbf{Unit conversion.} Express $60\ \mathrm{km\,h^{-1}}$ in $\mathrm{m\,s^{-1}}$. Show the conversion factor you use.
\item \textbf{Thinking distance.} Compute the distance travelled during the reaction time (before braking). Why is the speed constant during this phase?
\item \textbf{Braking distance and time.} Using a suitable suvat equation, compute the braking distance and the braking time. Hence find the \cle{total stopping distance}. Does the bike stop before reaching the goat $40\ \mathrm{m}$ ahead?
\item \textbf{Velocity--time graph.} Draw the velocity--time graph of the whole emergency stop (reaction phase + braking phase). Compute the total stopping distance as the \emph{area under the graph} and verify that it agrees with Task 3.
\item \textbf{The other two graphs.} On the same time axis, sketch the displacement--time graph and the acceleration--time graph. Indicate clearly the gradient of the $x$--$t$ curve at $t = 0$ and at the stopping instant, and the value of $a$ in each phase.
\item \textbf{The falling phone.} Compute the time the phone takes to fall $1.2\ \mathrm{m}$ vertically. Assuming the phone slips at the instant braking begins, compute the horizontal distance travelled by the \emph{bike} during the fall. (Bonus: what horizontal distance does the \emph{phone} cover if it keeps the horizontal velocity it had when it slipped?)
\item \textbf{Speed comparison.} Repeat the stopping-distance calculation for $40\ \mathrm{km\,h^{-1}}$ and $80\ \mathrm{km\,h^{-1}}$ (same reaction time, same deceleration). Present the three results in a table.
\item \textbf{Recommendation.} In at most six lines, write a safety recommendation for bend-skin riders, using your table to justify a speed limit near schools and markets.
\end{enumerate}

\courssub{Model Answer}

\textbf{Task 1. Unit conversion.}
Since $1\ \mathrm{km\,h^{-1}} = \dfrac{1000}{3600}\ \mathrm{m\,s^{-1}} = \dfrac{1}{3.6}\ \mathrm{m\,s^{-1}}$,
\[ u = \frac{60}{3.6} = \frac{50}{3} \approx 16.7\ \mathrm{m\,s^{-1}}. \]

\textbf{Task 2. Thinking distance.}
During the reaction time the brakes are not yet applied, so the velocity is constant (no horizontal force acts on the bike):
\[ s_{\text{think}} = u\,t_r = \frac{50}{3}\times 0.8 = \frac{40}{3} \approx 13.3\ \mathrm{m}. \]

\textbf{Task 3. Braking distance and time.}
Braking phase: $u = \tfrac{50}{3}\ \mathrm{m\,s^{-1}}$, $v = 0$, $a = -4\ \mathrm{m\,s^{-2}}$.
From $v^2 = u^2 + 2as$:
\[ 0 = \left(\tfrac{50}{3}\right)^2 + 2(-4)s_{\text{brake}} \;\Longrightarrow\; s_{\text{brake}} = \frac{2500/9}{8} = \frac{625}{18} \approx 34.7\ \mathrm{m}. \]
Braking time from $v = u + at$:
\[ t_b = \frac{u}{|a|} = \frac{50/3}{4} = \frac{25}{6} \approx 4.17\ \mathrm{s}. \]
Total stopping distance:
\[ \boxed{\,s_{\text{stop}} = \tfrac{40}{3} + \tfrac{625}{18} = \tfrac{865}{18} \approx 48.1\ \mathrm{m}\,} \]
Since $48.1 > 40$, \textbf{the bike hits the goat}: it cannot stop within $40\ \mathrm{m}$.

\textbf{Task 4. Velocity--time graph and area check.}
The $v$--$t$ graph consists of a horizontal segment at $v = 50/3\ \mathrm{m\,s^{-1}}$ from $t=0$ to $t=0.8\ \mathrm{s}$, followed by a straight line of gradient $-4\ \mathrm{m\,s^{-2}}$ reaching $v=0$ at $t = 0.8 + 25/6 = 149/30 \approx 4.97\ \mathrm{s}$.

\begin{popfigure}
\begin{tikzpicture}[scale=0.7]
\begin{axis}[
    axis lines=middle,
    xlabel={$t\ (\mathrm{s})$}, ylabel={$v\ (\mathrm{m\,s^{-1}})$},
    xmin=0, xmax=5.5, ymin=0, ymax=18,
    xtick={0,0.8,4.97}, xticklabels={$0$,$0.8$,$4.97$},
    ytick={0,16.67}, yticklabels={$0$,$16.7$},
    width=10cm, height=6cm,
    clip=false
]
\addplot[popblue, thick] coordinates {(0,16.67) (0.8,16.67)};
\addplot[popblue, thick] coordinates {(0.8,16.67) (4.97,0)};
\addplot[poporange, dashed] coordinates {(0.8,0) (0.8,16.67)};
\addplot[poporange, dashed] coordinates {(4.97,0) (4.97,16.67)};
\node[popblue, anchor=south west] at (axis cs:0.2,17) {$v = 50/3$};
\node[popblue, anchor=south] at (axis cs:2.5,8) {$a = -4$};
\end{axis}
\end{tikzpicture}
\legende{Velocity--time graph for the emergency stop.}
\end{popfigure}

Area under the graph = rectangle + triangle:
\[ A = \underbrace{\left(\frac{50}{3}\times 0.8\right)}_{\text{rectangle}} + \underbrace{\frac{1}{2}\left(\frac{50}{3}\right)\left(\frac{25}{6}\right)}_{\text{triangle}} = \frac{40}{3} + \frac{625}{18} = \frac{865}{18} \approx 48.1\ \mathrm{m}. \quad \checkmark \]

\textbf{Task 5. Displacement--time and acceleration--time graphs.}

\begin{popfigure}
\begin{tikzpicture}[scale=0.7]
\begin{axis}[
    axis lines=middle,
    xlabel={$t\ (\mathrm{s})$}, ylabel={$x\ (\mathrm{m})$},
    xmin=0, xmax=5.5, ymin=0, ymax=52,
    xtick={0,0.8,4.97}, xticklabels={$0$,$0.8$,$4.97$},
    ytick={0,13.33,48.06}, yticklabels={$0$,$13.3$,$48.1$},
    width=10cm, height=6cm,
    clip=false,
    samples=100
]
% Phase 1: constant velocity
\addplot[popblue, thick, domain=0:0.8] {50/3*x};
% Phase 2: constant deceleration
\addplot[popblue, thick, domain=0.8:4.97] {40/3 + 50/3*(x-0.8) - 2*(x-0.8)^2};
% Tangent at t=0
\addplot[poporange, dashed, domain=0:0.8] {50/3*x};
% Tangent at t=4.97 (horizontal)
\addplot[poporange, dashed, domain=4.5:5.2] {48.06};
\node[popblue, anchor=south west] at (axis cs:0.2,5) {gradient $= 16.7$};
\node[popblue, anchor=south] at (axis cs:4.5,50) {gradient $= 0$};
\end{axis}
\end{tikzpicture}
\legende{Displacement--time graph. Gradient at $t=0$ is $u = 16.7\ \mathrm{m\,s^{-1}}$; gradient at stop is $0$.}
\end{popfigure}

\begin{popfigure}
\begin{tikzpicture}[scale=0.7]
\begin{axis}[
    axis lines=middle,
    xlabel={$t\ (\mathrm{s})$}, ylabel={$a\ (\mathrm{m\,s^{-2}})$},
    xmin=0, xmax=5.5, ymin=-5, ymax=1,
    xtick={0,0.8,4.97}, xticklabels={$0$,$0.8$,$4.97$},
    ytick={-4,0}, yticklabels={$-4$,$0$},
    width=10cm, height=4cm,
    clip=false
]
\addplot[popgreen, thick] coordinates {(0,0) (0.8,0)};
\addplot[popgreen, thick] coordinates {(0.8,-4) (4.97,-4)};
\addplot[popgreen, thick] coordinates {(4.97,0) (5.5,0)};
% open/closed circles
\fill[popgreen] (axis cs:0,0) circle (2pt);
\fill[popgreen] (axis cs:0.8,0) circle (2pt);
\draw[popgreen, fill=white] (axis cs:0.8,-4) circle (2pt);
\fill[popgreen] (axis cs:4.97,-4) circle (2pt);
\draw[popgreen, fill=white] (axis cs:4.97,0) circle (2pt);
\node[popgreen, anchor=north] at (axis cs:0.4,0.3) {$a=0$};
\node[popgreen, anchor=north] at (axis cs:2.9,-3.7) {$a=-4$};
\node[popgreen, anchor=south] at (axis cs:5.2,0.3) {$a=0$};
\end{axis}
\end{tikzpicture}
\legende{Acceleration--time graph. $a=0$ during reaction, $a=-4\ \mathrm{m\,s^{-2}}$ during braking, $a=0$ after stop.}
\end{popfigure}

\textbf{Task 6. The falling phone.}
Vertical free fall from rest: $h = \tfrac{1}{2}gt^2$.
\[ t_{\text{fall}} = \sqrt{\frac{2h}{g}} = \sqrt{\frac{2\times 1.2}{10}} = \sqrt{0.24} \approx 0.49\ \mathrm{s}. \]
During this time the bike is still braking (since $0.49 < 4.17$). Distance travelled by bike in $0.49\ \mathrm{s}$ after brakes applied:
\[ s_{\text{bike}} = u t - \tfrac{1}{2} a t^2 = \frac{50}{3}(0.49) - 2(0.24) \approx 8.17 - 0.48 = 7.69\ \mathrm{m}. \]
(Bonus: the phone retains the horizontal velocity $u = 50/3\ \mathrm{m\,s^{-1}}$ at the moment it slips, so its horizontal distance is $s_{\text{phone}} = u\,t_{\text{fall}} = \frac{50}{3}\times 0.49 \approx 8.17\ \mathrm{m}$.)

\textbf{Task 7. Speed comparison table.}

\begin{center}
\begin{tabular}{|l|c|c|c|}
\hline
\rowcolor{popblueL} Speed & Thinking distance & Braking distance & Total stopping distance \\ \hline
$40\ \mathrm{km\,h^{-1}}$ ($11.1\ \mathrm{m\,s^{-1}}$) & $8.9\ \mathrm{m}$ & $15.4\ \mathrm{m}$ & $24.3\ \mathrm{m}$ \\ \hline
$60\ \mathrm{km\,h^{-1}}$ ($16.7\ \mathrm{m\,s^{-1}}$) & $13.3\ \mathrm{m}$ & $34.7\ \mathrm{m}$ & $48.1\ \mathrm{m}$ \\ \hline
$80\ \mathrm{km\,h^{-1}}$ ($22.2\ \mathrm{m\,s^{-1}}$) & $17.8\ \mathrm{m}$ & $61.7\ \mathrm{m}$ & $79.5\ \mathrm{m}$ \\ \hline
\end{tabular}
\end{center}

\textbf{Task 8. Safety recommendation.}
\emph{``Doubling the speed from $40$ to $80\ \mathrm{km\,h^{-1}}$ more than \emph{triples} the stopping distance ($24.3\ \mathrm{m} \to 79.5\ \mathrm{m}$), because braking distance grows like $u^2$. At $60\ \mathrm{km\,h^{-1}}$ a bend-skin needs about $48\ \mathrm{m}$ to stop — more than the length of a football pitch's half. We recommend a $40\ \mathrm{km\,h^{-1}}$ limit near schools and markets, and that riders keep at least $50\ \mathrm{m}$ of clear road ahead.''}

\begin{attention}[Marking Points]
\begin{itemize}
\item Do not forget to convert $\mathrm{km\,h^{-1}}$ to $\mathrm{m\,s^{-1}}$ \emph{before} using suvat equations.
\item Do not mix the reaction phase (constant velocity, $a=0$) with the braking phase (uniform deceleration, $a=-4\ \mathrm{m\,s^{-2}}$).
\item The phone's fall time depends \emph{only} on the vertical height and $g$, not on the bike's speed.
\item When sketching graphs, ensure the time axis is consistent across $v$--$t$, $x$--$t$ and $a$--$t$ plots.
\item The gradient of the $x$--$t$ graph gives velocity; the gradient of the $v$--$t$ graph gives acceleration; the area under the $v$--$t$ graph gives displacement.
\end{itemize}
\end{attention}

\newpage

\coursec{Methods and worked examples}

\coursec{Kinematics of particles moving in a straight line}

\courssub{Concept of Motion in One Dimension}

\begin{definition}[Displacement, velocity and acceleration]
A particle moves along a straight line. Choose an origin $O$ and a positive direction.
\begin{itemize}
\item \textbf{Displacement} $s$ (or $x$) is the signed distance from $O$; unit: metre (\si{m}).
\item \textbf{Velocity} $v$ is the rate of change of displacement: $v = \dfrac{ds}{dt}$; unit: \si{m.s^{-1}}.
\item \textbf{Acceleration} $a$ is the rate of change of velocity: $a = \dfrac{dv}{dt} = \dfrac{d^2s}{dt^2}$; unit: \si{m.s^{-2}}.
\item \textbf{Speed} is the magnitude of velocity: $|v|$.
\item \textbf{Retardation} (deceleration) means acceleration opposite to the velocity; its magnitude is positive.
\end{itemize}
\end{definition}

\begin{propriete}[Uniform acceleration]
If acceleration $a$ is constant, the motion is \emph{uniformly accelerated}. The five quantities $s$ (displacement from start), $u$ (initial velocity), $v$ (final velocity), $a$, $t$ (time elapsed) are linked by the \cle{suvat equations}:
\begin{align}
v &= u + at \tag{1}\\
s &= ut + \tfrac{1}{2}at^2 \tag{2}\\
v^2 &= u^2 + 2as \tag{3}\\
s &= \frac{u+v}{2}\,t \tag{4}\\
s &= vt - \tfrac{1}{2}at^2 \tag{5}
\end{align}
\textbf{Derivation (examinable):} From $a = \frac{dv}{dt} \Rightarrow v = u + at$. Then $v = \frac{ds}{dt} \Rightarrow s = \int_0^t (u+at)\,dt = ut + \frac{1}{2}at^2$. Eliminate $t$ from (1) and (2) to get (3). Equation (4) follows from average velocity $\frac{u+v}{2}$ for constant $a$. Equation (5) is obtained by substituting $u = v-at$ into (2).
\end{propriete}

\begin{aretenir}[Sign convention]
Always fix a positive direction at the start. All vector quantities ($s, u, v, a$) carry signs relative to that direction. A negative velocity means motion in the negative direction; a negative acceleration opposite to the velocity means retardation.
\end{aretenir}

\begin{attention}[Common pitfalls]
\begin{itemize}
\item Forgetting to convert $\si{km.h^{-1}}$ to $\si{m.s^{-1}}$ (multiply by $\frac{5}{18}$).
\item Using $g = \SI{9.8}{m.s^{-2}}$ unless the question specifies $g=\SI{10}{m.s^{-2}}$.
\item Confusing \emph{distance} (scalar, sum of path lengths) with \emph{displacement} (vector, final position minus initial position).
\item Taking the wrong root when solving a quadratic for time; always reject negative or unphysical times.
\end{itemize}
\end{attention}

\courssub{Equations of Motion under Uniform Acceleration}

\begin{methode}[Solving a uniform-acceleration problem with the suvat equations]
\begin{enumerate}
\item \textbf{Draw a diagram} and fix a positive direction (usually the direction of motion). Write down the five quantities: $s$ (displacement), $u$ (initial velocity), $v$ (final velocity), $a$ (acceleration), $t$ (time). Give each known quantity its correct sign.
\item \textbf{Identify the three known and the one unknown} you need. Note which quantity is \emph{not} mentioned.
\item \textbf{Select the equation that omits the missing quantity}:
\begin{center}
\begin{tabular}{|c|c|c|}
\hline
\rowcolor{popblueL} Equation & Contains & Omits \\
\hline
$v = u + at$ & $v,u,a,t$ & $s$ \\
\hline
$s = ut + \tfrac{1}{2}at^{2}$ & $s,u,a,t$ & $v$ \\
\hline
$v^{2} = u^{2} + 2as$ & $v,u,a,s$ & $t$ \\
\hline
$s = \dfrac{(u+v)}{2}\,t$ & $s,u,v,t$ & $a$ \\
\hline
$s = vt - \tfrac{1}{2}at^{2}$ & $s,v,a,t$ & $u$ \\
\hline
\end{tabular}
\end{center}
\item \textbf{Substitute with units}, solve, and check the sign and the magnitude of the answer against common sense.
\item \textbf{Interpret}: a negative velocity means motion in the negative direction; a negative acceleration opposite to motion means a \cle{retardation} (deceleration).
\end{enumerate}
\textbf{Reflexes:} ``starts from rest'' $\Rightarrow u=0$; ``comes to rest'' $\Rightarrow v=0$; convert $\si{km.h^{-1}}$ to $\si{m.s^{-1}}$ by multiplying by $\tfrac{1000}{3600}=\tfrac{5}{18}$.
\end{methode}

\begin{exoresolu}[Braking of a taxi in Douala]
A taxi travelling at $\SI{72}{km.h^{-1}}$ along a straight road in Douala sees an obstacle $\SI{50}{m}$ ahead. The driver applies the brakes, producing a uniform retardation. The taxi just stops at the obstacle. Find (a) the retardation, (b) the time taken to stop.
\tcblower
\textbf{Solution.} Take the direction of motion as positive.

\textbf{Step 1: convert units.}
\[u = 72\times\frac{5}{18} = \SI{20}{m.s^{-1}}, \qquad v = 0, \qquad s = \SI{50}{m}.\]

\textbf{(a)} Time is not given and not required, so use $v^{2}=u^{2}+2as$:
\[0 = 20^{2} + 2a(50) \;\Longrightarrow\; 100a = -400 \;\Longrightarrow\; a = \SI{-4}{m.s^{-2}}.\]
The retardation is therefore $\boxed{\SI{4}{m.s^{-2}}}$ (the negative sign indicates acceleration opposing the motion).

\textbf{(b)} Now use $v = u + at$:
\[0 = 20 - 4t \;\Longrightarrow\; t = \boxed{\SI{5}{s}}.\]

\textbf{Check:} average velocity $=\dfrac{u+v}{2}=\SI{10}{m.s^{-1}}$, so $s = 10\times 5 = \SI{50}{m}$. \checkmark
\end{exoresolu}

\courssub{Displacement–Time and Velocity–Time Graphs}

\begin{methode}[Interpreting a displacement–time ($s$–$t$) graph]
\begin{enumerate}
\item The \textbf{gradient} of the $s$–$t$ graph at any instant gives the \cle{velocity}: $v=\dfrac{ds}{dt}$.
\item A \textbf{horizontal section} means the particle is at rest; a \textbf{straight inclined section} means constant velocity; a \textbf{curve} means the velocity is changing (acceleration).
\item A \textbf{negative gradient} means motion back towards the origin (velocity in the negative direction).
\item The \textbf{total distance travelled} is found by adding the magnitudes of all changes of position; the \textbf{displacement} is the final position minus the initial position (signs included).
\item To \emph{sketch} an $s$–$t$ graph from a description: mark the key times, join them with straight segments (constant velocity) or smooth curves (acceleration), making the gradient steeper where the speed is greater.
\end{enumerate}
\end{methode}

\begin{exoresolu}[A walk to the market]
A trader leaves home and walks along a straight road. Her displacement $s$ (in metres) from home at time $t$ (in seconds) is recorded: she walks at constant velocity and reaches the market $\SI{400}{m}$ away after $\SI{200}{s}$; she stays there for $\SI{100}{s}$; she then walks back home at constant speed, arriving at $t=\SI{500}{s}$.
(a) Sketch the displacement–time graph.
(b) Find her velocity on each stage of the journey.
(c) Find her total distance travelled and her average speed for the whole journey.
\tcblower
\textbf{Solution.}

\textbf{(a)} The graph has three straight sections: rising from $(0,0)$ to $(200,400)$, horizontal from $(200,400)$ to $(300,400)$, and falling from $(300,400)$ to $(500,0)$.

\begin{popfigure}
\begin{tikzpicture}
\begin{axis}[
  width=9cm, height=6cm,
  axis lines=left,
  xlabel={$t$ (s)}, ylabel={$s$ (m)},
  xmin=0, xmax=520, ymin=0, ymax=460,
  xtick={0,100,200,300,400,500},
  ytick={0,100,200,300,400},
  grid=both, grid style={popline, very thin},
  label style={font=\small},
]
\addplot[very thick, popblue] coordinates {(0,0) (200,400) (300,400) (500,0)};
\addplot[only marks, mark=*, poporange] coordinates {(200,400) (300,400)};
\end{axis}
\end{tikzpicture}
\legende{Displacement–time graph for the trader's journey.}
\end{popfigure}

\textbf{(b)} Velocity = gradient of the $s$–$t$ graph.
\begin{align*}
\text{Outward: } & v_1 = \frac{400-0}{200-0} = \SI{2}{m.s^{-1}}.\\
\text{At market: } & v_2 = \frac{400-400}{300-200} = \SI{0}{m.s^{-1}} \quad (\text{at rest}).\\
\text{Return: } & v_3 = \frac{0-400}{500-300} = \SI{-2}{m.s^{-1}}.
\end{align*}
The negative sign shows she is moving back towards home.

\textbf{(c)} Total distance $= 400 + 0 + 400 = \SI{800}{m}$. Total time $=\SI{500}{s}$.
\[\text{Average speed} = \frac{800}{500} = \boxed{\SI{1.6}{m.s^{-1}}}.\]
Note: her average \emph{velocity} is $\dfrac{0-0}{500} = \SI{0}{m.s^{-1}}$, since the total displacement is zero.
\end{exoresolu}

\begin{methode}[Extracting information from a velocity–time ($v$–$t$) graph]
\begin{enumerate}
\item The \textbf{gradient} of the $v$–$t$ graph gives the \cle{acceleration}: $a=\dfrac{dv}{dt}$.
\item The \textbf{area between the graph and the time axis} gives the \cle{displacement}. Areas \emph{below} the axis count as negative.
\item Split the area into triangles, rectangles and trapezia: triangle $=\tfrac12 bh$, rectangle $=bh$, trapezium $=\tfrac12(a+b)h$.
\item \textbf{Total distance} = sum of the \emph{magnitudes} of the areas; \textbf{displacement} = algebraic sum.
\item To \emph{sketch} a $v$–$t$ graph from a description: constant acceleration $\Rightarrow$ straight sloping line; constant velocity $\Rightarrow$ horizontal line; rest $\Rightarrow$ line on the time axis.
\end{enumerate}
\end{methode}

\begin{exoresolu}[Journey of a bus between two stops]
A bus starts from rest at a stop and accelerates uniformly for $\SI{10}{s}$ until it reaches $\SI{15}{m.s^{-1}}$. It maintains this velocity for $\SI{30}{s}$, then decelerates uniformly, coming to rest at the next stop after a further $\SI{15}{s}$.
(a) Sketch the velocity–time graph.
(b) Find the acceleration during each stage.
(c) Find the distance between the two stops.
\tcblower
\textbf{Solution.}

\textbf{(a)} Three straight sections: rising from $(0,0)$ to $(10,15)$, horizontal from $(10,15)$ to $(40,15)$, falling from $(40,15)$ to $(55,0)$.

\begin{popfigure}
\begin{tikzpicture}
\begin{axis}[
  width=9.5cm, height=6cm,
  axis lines=left,
  xlabel={$t$ (s)}, ylabel={$v$ (m\,s$^{-1}$)},
  xmin=0, xmax=58, ymin=0, ymax=18,
  xtick={0,10,20,30,40,50},
  ytick={0,5,10,15},
  grid=both, grid style={popline, very thin},
  label style={font=\small},
]
\addplot[very thick, popgreen] coordinates {(0,0) (10,15) (40,15) (55,0)};
\addplot[popgreenL, draw=none] coordinates {(0,0) (10,15) (40,15) (55,0)} \closedcycle;
\end{axis}
\end{tikzpicture}
\legende{Velocity–time graph of the bus; the shaded area gives the distance.}
\end{popfigure}

\textbf{(b)} Acceleration = gradient.
\begin{align*}
\text{Stage 1: } & a_1 = \frac{15-0}{10-0} = \SI{1.5}{m.s^{-2}}.\\
\text{Stage 2: } & a_2 = 0 \quad (\text{constant velocity}).\\
\text{Stage 3: } & a_3 = \frac{0-15}{55-40} = \SI{-1}{m.s^{-2}} \quad (\text{retardation } \SI{1}{m.s^{-2}}).
\end{align*}

\textbf{(c)} Distance = total area under the graph = triangle $+$ rectangle $+$ triangle:
\[s = \tfrac12(10)(15) + (30)(15) + \tfrac12(15)(15) = 75 + 450 + 112.5 = \boxed{\SI{637.5}{m}}.\]
\textbf{Check (trapezium formula):} the whole area is a trapezium of parallel sides $30$ and $55$, height $15$: $s=\tfrac12(30+55)(15)=\SI{637.5}{m}$. \checkmark
\end{exoresolu}

\courssub{Acceleration–Time Graphs and Connecting the Three Graphs}

\begin{methode}[Moving between $a$–$t$, $v$–$t$ and $s$–$t$ graphs]
\begin{enumerate}
\item \textbf{From $a$–$t$ to $v$–$t$:} the change in velocity over an interval equals the \emph{area} under the $a$–$t$ graph: $\Delta v = \int a\,dt$. Add the initial velocity $u$.
\item \textbf{From $v$–$t$ to $s$–$t$:} the change in displacement equals the \emph{area} under the $v$–$t$ graph: $\Delta s = \int v\,dt$.
\item \textbf{From $s$–$t$ to $v$–$t$ to $a$–$t$:} differentiate: $v=\dfrac{ds}{dt}$, $a=\dfrac{dv}{dt}=\dfrac{d^{2}s}{dt^{2}}$.
\item \textbf{Key correspondences to memorise:}
\begin{center}
\begin{tabular}{|c|c|c|}
\hline
\rowcolor{popblueL} $s$–$t$ graph & $v$–$t$ graph & $a$–$t$ graph \\
\hline
Straight, sloping & Horizontal (constant $v$) & On the axis ($a=0$) \\
\hline
Curve, steepening & Straight, sloping & Horizontal ($a\neq 0$) \\
\hline
Curve, changing curvature & Curve & Sloping line \\
\hline
\end{tabular}
\end{center}
\item Always check \textbf{consistency}: where the $v$–$t$ graph crosses the axis, the $s$–$t$ graph has a turning point; where the $a$–$t$ graph crosses the axis, the $v$–$t$ graph has a turning point.
\end{enumerate}
\end{methode}

\begin{exoresolu}[From an acceleration–time graph to velocity and displacement]
A particle starts from rest at $O$ and moves in a straight line. Its acceleration is $\SI{4}{m.s^{-2}}$ for the first $\SI{5}{s}$, then zero for the next $\SI{5}{s}$, then $\SI{-2}{m.s^{-2}}$ for the final $\SI{10}{s}$.
(a) Sketch the acceleration–time and velocity–time graphs.
(b) Find the velocity at $t=5$, $t=10$ and $t=20$.
(c) Find the total displacement after $\SI{20}{s}$.
\tcblower
\textbf{Solution.}

\textbf{(a)} The $a$–$t$ graph consists of three horizontal segments at heights $4$, $0$, $-2$.

\begin{popfigure}
\begin{tikzpicture}
\begin{axis}[
  width=9.5cm, height=5.5cm,
  axis lines=left,
  xlabel={$t$ (s)}, ylabel={$a$ (m\,s$^{-2}$)},
  xmin=0, xmax=21, ymin=-3, ymax=5,
  xtick={0,5,10,15,20},
  ytick={-2,0,2,4},
  grid=both, grid style={popline, very thin},
  label style={font=\small},
]
\addplot[very thick, poppurple] coordinates {(0,4) (5,4)};
\addplot[very thick, poppurple] coordinates {(5,0) (10,0)};
\addplot[very thick, poppurple] coordinates {(10,-2) (20,-2)};
\end{axis}
\end{tikzpicture}
\legende{Acceleration–time graph of the particle.}
\end{popfigure}

\textbf{(b)} Change in velocity = area under the $a$–$t$ graph, with $u=0$.
\begin{align*}
v(5) &= 0 + (4)(5) = \SI{20}{m.s^{-1}}.\\
v(10) &= 20 + (0)(5) = \SI{20}{m.s^{-1}}.\\
v(20) &= 20 + (-2)(10) = \SI{0}{m.s^{-1}}.
\end{align*}
The particle ends at rest — a useful check.

The $v$–$t$ graph is therefore: rising line from $(0,0)$ to $(5,20)$, horizontal from $(5,20)$ to $(10,20)$, falling from $(10,20)$ to $(20,0)$.

\begin{popfigure}
\begin{tikzpicture}
\begin{axis}[
  width=9.5cm, height=5.5cm,
  axis lines=left,
  xlabel={$t$ (s)}, ylabel={$v$ (m\,s$^{-1}$)},
  xmin=0, xmax=21, ymin=0, ymax=24,
  xtick={0,5,10,15,20},
  ytick={0,10,20},
  grid=both, grid style={popline, very thin},
  label style={font=\small},
]
\addplot[very thick, popblue] coordinates {(0,0) (5,20) (10,20) (20,0)};
\end{axis}
\end{tikzpicture}
\legende{Velocity–time graph obtained from the areas under the $a$–$t$ graph.}
\end{popfigure}

\textbf{(c)} Displacement = area under the $v$–$t$ graph:
\[s = \tfrac12(5)(20) + (5)(20) + \tfrac12(10)(20) = 50 + 100 + 100 = \boxed{\SI{250}{m}}.\]
\end{exoresolu}

\courssub{Free Fall and Motion under Gravity}

\begin{definition}[Free fall]
A body moving under the influence of gravity alone (air resistance neglected) is in \cle{free fall}. Near the Earth's surface the acceleration has constant magnitude $g = \SI{9.8}{m.s^{-2}}$ (use $g=\SI{10}{m.s^{-2}}$ only if the question instructs so) directed vertically downwards.
\end{definition}

\begin{methode}[Solving free-fall problems]
\begin{enumerate}
\item Choose a \textbf{positive direction} (usually upwards) and \textbf{keep it for the whole solution}. With upwards positive, $a = -g = \SI{-9.8}{m.s^{-2}}$.
\item List $s, u, v, a, t$ with signs. Remember: at the \textbf{greatest height}, $v=0$; when the body \textbf{returns to the launch level}, $s=0$.
\item Apply the suvat equations exactly as for any uniform acceleration — gravity \emph{is} a uniform acceleration.
\item For a body \textbf{dropped} from rest: $u=0$; for a body \textbf{projected downwards}, $u>0$ downwards.
\item \textbf{Time of flight} for a body projected upwards from ground level: $T=\dfrac{2u}{g}$; \textbf{greatest height}: $H=\dfrac{u^{2}}{2g}$. (Derive these in the exam rather than quoting blindly.)
\item Ignore air resistance unless the question states otherwise — this is the standard GCE modelling assumption.
\end{enumerate}
\end{methode}

\begin{exoresolu}[Stone thrown upwards from a cliff]
A stone is thrown vertically upwards with velocity $\SI{14}{m.s^{-1}}$ from the top of a cliff $\SI{30}{m}$ above the sea. Taking $g=\SI{9.8}{m.s^{-2}}$ and neglecting air resistance, find
(a) the greatest height above the cliff top reached by the stone,
(b) the time at which the stone passes the cliff top on its way down,
(c) the total time until the stone hits the sea,
(d) the velocity with which it hits the sea.
\tcblower
\textbf{Solution.} Take \textbf{upwards as positive}, origin at the cliff top. Then $u = \SI{14}{m.s^{-1}}$, $a = \SI{-9.8}{m.s^{-2}}$.

\textbf{(a)} At greatest height, $v=0$. Using $v^{2}=u^{2}+2as$:
\[0 = 14^{2} + 2(-9.8)s \;\Longrightarrow\; s = \frac{196}{19.6} = \boxed{\SI{10}{m}}.\]

\textbf{(b)} Passing the cliff top again means $s=0$. Using $s=ut+\tfrac12 at^{2}$:
\[0 = 14t - 4.9t^{2} = t(14-4.9t).\]
The solutions are $t=0$ (the throw) and $t=\dfrac{14}{4.9}=\dfrac{20}{7}\approx\boxed{\SI{2.86}{s}}$.

\textbf{(c)} Hitting the sea means $s=\SI{-30}{m}$ (below the origin):
\[-30 = 14t - 4.9t^{2} \;\Longrightarrow\; 4.9t^{2} - 14t - 30 = 0.\]
By the quadratic formula:
\[t = \frac{14 \pm \sqrt{14^{2}+4(4.9)(30)}}{2(4.9)} = \frac{14 \pm \sqrt{196+588}}{9.8} = \frac{14 \pm 28}{9.8}.\]
Taking the positive root: $t = \dfrac{42}{9.8} = \dfrac{30}{7} \approx \boxed{\SI{4.29}{s}}$. (The negative root $t=-\tfrac{12}{9.8}$ is rejected.)

\textbf{(d)} Using $v = u + at$:
\[v = 14 - 9.8\left(\tfrac{30}{7}\right) = 14 - 42 = \SI{-28}{m.s^{-1}}.\]
The stone hits the sea with speed $\boxed{\SI{28}{m.s^{-1}}}$ directed downwards.

\textbf{Check:} $v^{2}=u^{2}+2as = 196 + 2(-9.8)(-30) = 196+588=784$, so $v=\pm 28$. \checkmark\ The negative root is chosen since the stone is moving downwards.
\end{exoresolu}

\courssub{Two-Particle Problems (Meeting and Overtaking)}

\begin{methode}[Solving problems with two moving particles]
\begin{enumerate}
\item Write the \textbf{displacement of each particle as a function of time}, measured from the \emph{same} origin in the \emph{same} positive direction: $s_A(t)$ and $s_B(t)$.
\item If one particle starts later, use $t - t_0$ for its time of motion.
\item \textbf{Meeting / overtaking condition:} set $s_A(t) = s_B(t)$ and solve for $t$. Reject negative or otherwise impossible roots.
\item Substitute back to find the \textbf{position} of the meeting and the \textbf{velocities} at that instant if required.
\item Alternatively, read the meeting point off a pair of $s$–$t$ graphs drawn on the \emph{same} axes: the intersection of the two curves.
\end{enumerate}
\end{methode}

\begin{exoresolu}[The police car and the speeding lorry]
A lorry travels along a straight highway at a constant velocity of $\SI{20}{m.s^{-1}}$. As it passes a police checkpoint, a police car starts from rest and accelerates uniformly at $\SI{2.5}{m.s^{-2}}$ in pursuit. Find
(a) the time at which the police car catches the lorry,
(b) the distance from the checkpoint at which this happens,
(c) the velocity of the police car at that instant.
\tcblower
\textbf{Solution.} Take the checkpoint as origin, direction of travel positive, and $t=0$ when the lorry passes.

\textbf{Displacements:}
\[s_L = 20t \qquad \text{(lorry, constant velocity)},\]
\[s_P = 0\cdot t + \tfrac12(2.5)t^{2} = 1.25t^{2} \qquad \text{(police car, from rest)}.\]

\textbf{(a)} Catching up: $s_P = s_L$:
\[1.25t^{2} = 20t \;\Longrightarrow\; t(1.25t - 20) = 0.\]
$t=0$ corresponds to the moment the lorry passes the checkpoint. The overtaking occurs at
\[t = \frac{20}{1.25} = \boxed{\SI{16}{s}}.\]

\textbf{(b)} Distance: $s_L = 20\times 16 = \boxed{\SI{320}{m}}$ from the checkpoint.

\textbf{(c)} Police car's velocity: $v = u + at = 0 + 2.5\times 16 = \boxed{\SI{40}{m.s^{-1}}}$.

\textbf{Remark:} on the same $s$–$t$ axes, the lorry's graph is the straight line $s=20t$ and the police car's graph is the parabola $s=1.25t^{2}$; they intersect at $t=0$ and $t=16$, confirming the answer graphically.
\end{exoresolu}

\courssub{Graphical Treatment of Motion under Gravity}

\begin{methode}[Sketching and using graphs for vertical motion]
\begin{enumerate}
\item For a body projected \textbf{upwards} (upwards positive):
\begin{itemize}
\item the $v$–$t$ graph is a \textbf{straight line of gradient $-g$}, starting at $u$ and crossing the axis at the time of greatest height;
\item the $s$–$t$ graph is a \textbf{parabola} opening downwards, with its vertex at the greatest height;
\item the $a$–$t$ graph is a \textbf{horizontal line at $-g$} for the entire flight (including at the top!).
\end{itemize}
\item The \textbf{area} under the $v$–$t$ graph up to the crossing time gives the greatest height; the signed area up to any later time gives the displacement.
\item The \textbf{speed–time} graph differs from the velocity–time graph: speed is never negative, so the part below the axis is \emph{reflected} upwards (a ``V'' shape).
\item Use graph areas to \textbf{verify} algebraic answers — a powerful checking reflex in the examination.
\end{enumerate}
\end{methode}

\begin{exoresolu}[Ball dropped from a water tower]
A ball is released from rest from the top of a water tower in Bamenda, $\SI{45}{m}$ above the ground. Taking $g=\SI{10}{m.s^{-2}}$,
(a) sketch the velocity–time and displacement–time graphs (taking \emph{downwards} as positive) for the whole fall;
(b) use the $v$–$t$ graph to find the time of fall and the velocity of impact;
(c) verify your answers using the equations of motion.
\tcblower
\textbf{Solution.} Downwards positive: $u=0$, $a=g=\SI{10}{m.s^{-2}}$.

\textbf{(a)} The $v$–$t$ graph is a straight line through the origin of gradient $10$: $v=10t$. The $s$–$t$ graph is the parabola $s=5t^{2}$, ending at $s=45$.

\begin{popfigure}
\begin{tikzpicture}
\begin{axis}[
  width=7.5cm, height=5.5cm,
  axis lines=left,
  xlabel={$t$ (s)}, ylabel={$v$ (m\,s$^{-1}$)},
  xmin=0, xmax=3.4, ymin=0, ymax=34,
  xtick={0,1,2,3},
  ytick={0,10,20,30},
  grid=both, grid style={popline, very thin},
  label style={font=\small},
]
\addplot[very thick, poporange, domain=0:3] {10*x};
\addplot[popblueL, draw=none, domain=0:3] {10*x} \closedcycle;
\end{axis}
\end{tikzpicture}
\legende{Velocity–time graph of the falling ball; the shaded area equals $\SI{45}{m}$.}
\end{popfigure}

\textbf{(b)} The fall ends when the area under the $v$–$t$ graph equals $\SI{45}{m}$. If $T$ is the time of fall, the area is a triangle:
\[\tfrac12 \times T \times 10T = 45 \;\Longrightarrow\; 5T^{2} = 45 \;\Longrightarrow\; T^{2}=9 \;\Longrightarrow\; T=\boxed{\SI{3}{s}}.\]
Impact velocity: $v = 10T = \boxed{\SI{30}{m.s^{-1}}}$.

\textbf{(c)} Verification with equations. Using $s=ut+\tfrac12 at^{2}$:
\[45 = 0 + \tfrac12(10)t^{2} \;\Longrightarrow\; t^{2}=9 \;\Longrightarrow\; t=\SI{3}{s}.\ \checkmark\]
Using $v^{2}=u^{2}+2as$: $v^{2}=0+2(10)(45)=900$, so $v=\SI{30}{m.s^{-1}}$. \checkmark

The graphical and algebraic methods agree, as they must: the suvat equations are exactly the ``area and gradient'' statements for a straight-line $v$–$t$ graph.
\end{exoresolu}

\courssub{Real-Life Applications and Synthesis}

\begin{experience}[Reaction time and stopping distance]
A driver travelling at $\SI{30}{m.s^{-1}}$ ($\approx \SI{108}{km.h^{-1}}$) sees a hazard. The typical human reaction time is $\SI{0.7}{s}$. During this time the car travels at constant speed. The driver then brakes with a uniform retardation of $\SI{6}{m.s^{-2}}$.
\begin{enumerate}
\item Find the \textbf{thinking distance} (distance during reaction time).
\item Find the \textbf{braking distance}.
\item Find the \textbf{total stopping distance}.
\item Sketch the $v$–$t$ graph for the whole event.
\end{enumerate}
\tcblower
\textbf{Solution.}
\begin{enumerate}
\item Thinking distance $= u \times t_{\text{react}} = 30 \times 0.7 = \SI{21}{m}$.
\item Braking: $u=30$, $v=0$, $a=-6$. Use $v^2=u^2+2as$: $0 = 900 - 12s \Rightarrow s = \SI{75}{m}$.
\item Total stopping distance $= 21 + 75 = \boxed{\SI{96}{m}}$.
\item $v$–$t$ graph: horizontal line at $v=30$ from $t=0$ to $t=0.7$, then straight line down to $v=0$ at $t = 0.7 + \frac{30}{6} = 5.7\,\text{s}$. The area under the graph is $21 + 75 = 96\,\text{m}$.
\end{enumerate}
\end{experience}

\begin{savaistu}[Did you know?]
The \emph{Highway Code} in Cameroon (and many countries) quotes stopping distances based on a reaction time of $\SI{0.7}{s}$ and a deceleration of about $\SI{6}{m.s^{-2}}$ on dry roads. At $\SI{50}{km.h^{-1}}$ ($\approx \SI{13.9}{m.s^{-1}}$) the total stopping distance is roughly $\SI{23}{m}$; at $\SI{100}{km.h^{-1}}$ it quadruples to about $\SI{96}{m}$ because braking distance is proportional to $u^2$. This is why speed limits are critical for safety.
\end{savaistu}

\begin{exercice}[Consolidation exercises]
\begin{enumerate}
\item A particle moves with uniform acceleration. It passes point $A$ with velocity $\SI{10}{m.s^{-1}}$ and point $B$ with velocity $\SI{20}{m.s^{-1}}$. The distance $AB$ is $\SI{45}{m}$. Find the acceleration and the time taken to travel from $A$ to $B$.
\item A ball is thrown vertically upwards from ground level with speed $\SI{24.5}{m.s^{-1}}$. Taking $g=\SI{9.8}{m.s^{-2}}$, find (a) the greatest height, (b) the total time of flight, (c) the velocity when it is $\SI{14.7}{m}$ above the ground on its way down.
\item Two particles $P$ and $Q$ start simultaneously from the same point $O$ on a straight line. $P$ moves with constant velocity $\SI{5}{m.s^{-1}}$. $Q$ starts from rest with constant acceleration $\SI{2}{m.s^{-2}}$. Find (a) the time when $Q$ overtakes $P$, (b) the distance from $O$ at that moment, (c) the velocity of $Q$ then.
\item A car accelerates uniformly from rest to $\SI{20}{m.s^{-1}}$ in $\SI{10}{s}$, travels at constant speed for $\SI{20}{s}$, then decelerates uniformly to rest in $\SI{5}{s}$. Sketch the $v$–$t$ graph and find the total distance travelled.
\item A stone is dropped from a bridge $\SI{78.4}{m}$ above a river. At the same instant, another stone is thrown vertically downwards from the bridge with speed $\SI{19.6}{m.s^{-1}}$. Taking $g=\SI{9.8}{m.s^{-2}}$, find the time interval between the two stones hitting the water.
\end{enumerate}
\end{exercice}

\begin{corrige}[Exercise 1]
$a = \SI{3.33}{m.s^{-2}}$ (or $\frac{10}{3}$), $t = \SI{3}{s}$.
\end{corrige}
\begin{corrige}[Exercise 2]
(a) $H = \frac{u^2}{2g} = \frac{24.5^2}{19.6} = \SI{30.625}{m}$. (b) $T = \frac{2u}{g} = \frac{49}{9.8} = \SI{5}{s}$. (c) Downwards positive from max height: $s = 30.625 - 14.7 = 15.925\,\text{m}$, $u=0$, $a=9.8$, $v^2 = 2\times9.8\times15.925 = 312.13$, $v \approx \SI{17.67}{m.s^{-1}}$ downwards.
\end{corrige}
\begin{corrige}[Exercise 3]
(a) $5t = t^2 \Rightarrow t=5\,\text{s}$ (excluding $t=0$). (b) $s = 5\times5 = \SI{25}{m}$. (c) $v = 2\times5 = \SI{10}{m.s^{-1}}$.
\end{corrige}
\begin{corrige}[Exercise 4]
Area = $\frac12(10)(20) + (20)(20) + \frac12(5)(20) = 100 + 400 + 50 = \SI{550}{m}$.
\end{corrige}
\begin{corrige}[Exercise 5]
Dropped stone: $78.4 = \frac12(9.8)t_1^2 \Rightarrow t_1 = 4\,\text{s}$. Thrown stone: $78.4 = 19.6t_2 + 4.9t_2^2 \Rightarrow t_2^2 + 4t_2 - 16 = 0 \Rightarrow t_2 = -2 + 2\sqrt{5} \approx 2.47\,\text{s}$. Interval $\approx 4 - 2.47 = \SI{1.53}{s}$.
\end{corrige}

\newpage

\coursec{Exercises}

\courssub{I check my knowledge}

\begin{exercice}[MCQ: Uniform acceleration]
A car starts from rest and accelerates uniformly at \SI{2}{m.s^{-2}} for \SI{10}{s}. What is its final velocity?
\begin{enumerate}
\item \SI{5}{m.s^{-1}}
\item \SI{10}{m.s^{-1}}
\item \SI{20}{m.s^{-1}}
\item \SI{40}{m.s^{-1}}
\end{enumerate}
\end{exercice}

\begin{exercice}[True/False with justification]
State whether each statement is true or false. Justify your answer.
\begin{enumerate}
\item The gradient of a displacement–time graph gives the acceleration.
\item The area under a velocity–time graph gives the displacement.
\item For a particle in free fall, the acceleration–time graph is a horizontal line.
\item If the velocity–time graph is a straight line with negative slope, the particle is decelerating.
\end{enumerate}
\end{exercice}

\begin{exercice}[Definitions and direct calculations]
\begin{enumerate}
\item Define \emph{displacement}, \emph{velocity} and \emph{acceleration} for motion in a straight line.
\item A particle moves with constant velocity \SI{15}{m.s^{-1}} for \SI{8}{s}. Calculate its displacement.
\item A body accelerates uniformly from \SI{5}{m.s^{-1}} to \SI{25}{m.s^{-1}} in \SI{4}{s}. Find its acceleration and the distance covered.
\end{enumerate}
\end{exercice}

\begin{exercice}[Graph interpretation]
The velocity–time graph below shows the motion of a particle for \SI{12}{s}.
\begin{center}
\begin{tikzpicture}[scale=0.7]
\draw[->] (0,0) -- (13,0) node[right] {$t$ (s)};
\draw[->] (0,0) -- (0,9) node[above] {$v$ (m/s)};
\draw[thick,popblue] (0,0) -- (4,8) -- (8,8) -- (12,0);
\foreach \x in {4,8,12} \draw (\x,0.1) -- (\x,-0.1) node[below] {\x};
\foreach \y in {2,4,6,8} \draw (0.1,\y) -- (-0.1,\y) node[left] {\y};
\end{tikzpicture}
\end{center}
\begin{enumerate}
\item Describe the motion in each of the three phases.
\item Calculate the acceleration during the first \SI{4}{s}.
\item Find the total displacement after \SI{12}{s}.
\end{enumerate}
\end{exercice}

\courssub{I practise}

\begin{exercice}[Car journey with three phases]
A car accelerates uniformly from rest to \SI{25}{m.s^{-1}} in \SI{10}{s}, cruises at this speed for \SI{20}{s}, then brakes uniformly to rest in \SI{5}{s}.
\begin{enumerate}
\item Sketch the velocity–time graph for the whole journey.
\item Calculate the total distance travelled.
\item Find the average velocity for the journey.
\end{enumerate}
\end{exercice}

\begin{exercice}[Free fall from a building]
A stone is dropped from the top of a storey building and hits the ground after \SI{3}{s}. Taking $g = \SI{10}{m.s^{-2}}$,
\begin{enumerate}
\item find the height of the building,
\item find the impact velocity.
\end{enumerate}
\end{exercice}

\begin{exercice}[Vertical projection]
A ball is thrown vertically upward at \SI{20}{m.s^{-1}} from ground level. Taking $g = \SI{10}{m.s^{-2}}$,
\begin{enumerate}
\item find the maximum height reached,
\item find the total time of flight,
\item find the velocity at $t = \SI{3}{s}$.
\end{enumerate}
\end{exercice}

\begin{exercice}[Connecting graphs]
A particle starts from rest at $x=0$. Its acceleration–time graph consists of two rectangles: $a = \SI{2}{m.s^{-2}}$ for \SI{4}{s}, then $a = \SI{-3}{m.s^{-2}}$ until the particle comes to rest.
\begin{enumerate}
\item Sketch the velocity–time and displacement–time graphs.
\item Find the total displacement when the particle stops.
\end{enumerate}
\end{exercice}

\courssub{I prepare for the GCE}

\begin{exercice}[GCE style: Trapezium velocity–time graph] \hfill (12 marks)
A particle moves along a straight line. Its velocity $v$ (in m/s) at time $t$ (in s) is given by the trapezium-shaped graph below:
\begin{center}
\begin{tikzpicture}[scale=0.7]
\draw[->] (0,0) -- (11,0) node[right] {$t$};
\draw[->] (0,0) -- (0,7) node[above] {$v$};
\draw[thick,popblue] (0,0) -- (2,6) -- (6,6) -- (10,0);
\foreach \x in {2,6,10} \draw (\x,0.1) -- (\x,-0.1) node[below] {\x};
\foreach \y in {2,4,6} \draw (0.1,\y) -- (-0.1,\y) node[left] {\y};
\end{tikzpicture}
\end{center}
\begin{enumerate}
\item Find the acceleration during each of the three phases. \hfill (3 marks)
\item Calculate the total distance travelled in \SI{10}{s}. \hfill (4 marks)
\item Determine the average velocity for the whole journey. \hfill (2 marks)
\item Sketch the corresponding acceleration–time graph. \hfill (3 marks)
\end{enumerate}
\end{exercice}

\begin{exercice}[GCE style: Two objects meeting] \hfill (14 marks)
A coconut falls from the top of a \SI{20}{m} palm tree at the same instant a stone is thrown vertically upward from the ground with an initial velocity of \SI{15}{m.s^{-1}}. Take $g = \SI{10}{m.s^{-2}}$.
\begin{enumerate}
\item Write down the equations of motion for the coconut and the stone, choosing the upward direction as positive. \hfill (4 marks)
\item Find the time at which they cross each other. \hfill (4 marks)
\item Find the height above the ground at which they cross. \hfill (3 marks)
\item Determine the velocities of both objects at the crossing instant. \hfill (3 marks)
\end{enumerate}
\end{exercice}

\begin{exercice}[GCE style: Overtaking and stopping distance] \hfill (14 marks)
Two taxis leave the same park. Taxi A moves at a constant velocity of \SI{15}{m.s^{-1}}. Taxi B starts from rest \SI{5}{s} later with a uniform acceleration of \SI{2}{m.s^{-2}}.
\begin{enumerate}
\item Find the time (measured from the departure of taxi A) when taxi B overtakes taxi A. \hfill (5 marks)
\item Find the distance from the park where the overtaking occurs. \hfill (3 marks)
\end{enumerate}
A driver is travelling at \SI{90}{km.h^{-1}} and sees an obstacle \SI{60}{m} ahead. His reaction time is \SI{0.7}{s} and the maximum braking deceleration is \SI{5}{m.s^{-2}}.
\begin{enumerate}
\setcounter{enumi}{2}
\item Calculate the distance travelled during the reaction time. \hfill (2 marks)
\item Calculate the braking distance. \hfill (2 marks)
\item Does the car stop before hitting the obstacle? Justify. \hfill (2 marks)
\end{enumerate}
\end{exercice}

\newpage

\coursec{Solutions to the exercises}

\begin{corrige}[Exercise 1 — MCQ: Uniform acceleration]
The car starts from rest, so $u = 0$, with $a = \SI{2}{m.s^{-2}}$ and $t = \SI{10}{s}$. Using the first equation of motion:
\[
v = u + at = 0 + 2 \times 10 = \SI{20}{m.s^{-1}}.
\]
\textbf{Correct answer: (c)} $\SI{20}{m.s^{-1}}$.
\end{corrige}

\begin{corrige}[Exercise 2 — True/False with justification]
\begin{enumerate}
\item \textbf{False.} The gradient of a displacement–time graph gives the \emph{velocity}, since $v = \dfrac{dx}{dt}$. Acceleration is the gradient of the \emph{velocity}–time graph.
\item \textbf{True.} Since $v = \dfrac{dx}{dt}$, we have $x = \displaystyle\int v\,dt$, which is exactly the (signed) area under the velocity–time graph.
\item \textbf{True.} In free fall (neglecting air resistance) the acceleration is constant and equal to $g$, so the acceleration–time graph is a horizontal line $a = -g$ (taking upward as positive).
\item \textbf{False in general.} A straight velocity–time line with negative slope means the \emph{acceleration} is negative. The particle decelerates (speed decreases) only if velocity and acceleration have opposite signs. If the velocity is itself negative, a negative acceleration makes the particle \emph{speed up} in the negative direction.
\end{enumerate}
\end{corrige}

\begin{corrige}[Exercise 3 — Definitions and direct calculations]
\begin{enumerate}
\item \textbf{Displacement} is the change in position of the particle, measured from a fixed origin on the line; it is a vector quantity (in 1-D, a signed scalar), unit: metre (m). \textbf{Velocity} is the rate of change of displacement with respect to time, $v = \dfrac{dx}{dt}$; unit: $\mathrm{m\,s^{-1}}$. \textbf{Acceleration} is the rate of change of velocity with respect to time, $a = \dfrac{dv}{dt}$; unit: $\mathrm{m\,s^{-2}}$.
\item Constant velocity: $x = vt = 15 \times 8 = \SI{120}{m}$.
\item Acceleration: $a = \dfrac{v-u}{t} = \dfrac{25-5}{4} = \SI{5}{m.s^{-2}}$.\\
Distance: $s = \dfrac{u+v}{2}\,t = \dfrac{5+25}{2}\times 4 = \SI{60}{m}$. (Check: $s = ut + \tfrac12 at^2 = 20 + 40 = \SI{60}{m}$.)
\end{enumerate}
\end{corrige}

\begin{corrige}[Exercise 4 — Graph interpretation]
\begin{enumerate}
\item \textbf{Phase 1} ($0 \le t \le 4$): velocity increases uniformly from $0$ to $\SI{8}{m.s^{-1}}$ — uniform acceleration. \textbf{Phase 2} ($4 \le t \le 8$): velocity constant at $\SI{8}{m.s^{-1}}$ — zero acceleration. \textbf{Phase 3} ($8 \le t \le 12$): velocity decreases uniformly from $\SI{8}{m.s^{-1}}$ to $0$ — uniform deceleration (retardation).
\item $a = \dfrac{\Delta v}{\Delta t} = \dfrac{8-0}{4-0} = \SI{2}{m.s^{-2}}$.
\item The displacement is the area of the trapezium under the graph (bases $12\ \mathrm{s}$ and $4\ \mathrm{s}$, height $8\ \mathrm{m\,s^{-1}}$):
\[
s = \tfrac12(12+4)\times 8 = \SI{64}{m}.
\]
(Equivalently: triangle $\tfrac12(4)(8)=16$, rectangle $4\times 8 = 32$, triangle $\tfrac12(4)(8)=16$; total $16+32+16=\SI{64}{m}$.)
\end{enumerate}
\end{corrige}

\begin{corrige}[Exercise 5 — Car journey with three phases]
\begin{enumerate}
\item The graph consists of: a straight line from $(0,0)$ to $(10,25)$; a horizontal segment from $(10,25)$ to $(30,25)$; a straight line from $(30,25)$ to $(35,0)$.
\begin{center}
\begin{tikzpicture}[scale=0.22]
\draw[->] (0,0) -- (40,0) node[right] {$t$ (s)};
\draw[->] (0,0) -- (0,30) node[above] {$v$ (m/s)};
\draw[thick,popblue] (0,0) -- (10,25) -- (30,25) -- (35,0);
\foreach \x in {10,30,35} \draw (\x,0.8) -- (\x,-0.8) node[below] {\x};
\draw (0.4,25) -- (-0.4,25) node[left] {25};
\draw[dashed,gray] (10,0) -- (10,25);
\draw[dashed,gray] (30,0) -- (30,25);
\end{tikzpicture}
\end{center}
\item Total distance = area under the graph (trapezium with parallel sides $35\ \mathrm{s}$ and $20\ \mathrm{s}$, height $25\ \mathrm{m\,s^{-1}}$):
\[
s = \tfrac12(35+20)\times 25 = \tfrac12 \times 55 \times 25 = \SI{687.5}{m}.
\]
(Phase by phase: $\tfrac12(10)(25)=125$; $20\times 25 = 500$; $\tfrac12(5)(25)=62.5$; total $687.5$ m.)
\item Average velocity $= \dfrac{\text{total displacement}}{\text{total time}} = \dfrac{687.5}{35} \approx \SI{19.6}{m.s^{-1}}$.
\end{enumerate}
\end{corrige}

\begin{corrige}[Exercise 6 — Free fall from a building]
Take the downward direction as positive. Then $u = 0$ (dropped), $a = g = \SI{10}{m.s^{-2}}$, $t = \SI{3}{s}$.
\begin{enumerate}
\item Height of the building:
\[
h = ut + \tfrac12 gt^2 = 0 + \tfrac12 \times 10 \times 3^2 = \SI{45}{m}.
\]
\item Impact velocity:
\[
v = u + gt = 0 + 10 \times 3 = \SI{30}{m.s^{-1}} \quad (\text{downwards}).
\]
\end{enumerate}
\end{corrige}

\begin{corrige}[Exercise 7 — Vertical projection]
Take upward as positive: $u = \SI{20}{m.s^{-1}}$, $a = -g = \SI{-10}{m.s^{-2}}$.
\begin{enumerate}
\item At maximum height, $v = 0$. Using $v^2 = u^2 + 2as$:
\[
0 = 20^2 - 2(10)h \implies h = \dfrac{400}{20} = \SI{20}{m}.
\]
\item Time to reach the top: $0 = 20 - 10t \implies t = \SI{2}{s}$. By symmetry, the descent takes the same time, so the total time of flight is $T = 2 \times 2 = \SI{4}{s}$. (Alternatively, solve $s = 20t - 5t^2 = 0$, giving $t = 0$ or $t = 4$.)
\item At $t = \SI{3}{s}$: $v = u - gt = 20 - 10(3) = \SI{-10}{m.s^{-1}}$. The negative sign means the ball is moving \emph{downwards} at $\SI{10}{m.s^{-1}}$ (it is on its way down).
\end{enumerate}
\end{corrige}

\begin{corrige}[Exercise 8 — Connecting graphs]
\begin{enumerate}
\item \textbf{Phase 1} ($0 \le t \le 4$): $a = \SI{2}{m.s^{-2}}$, so $v = 2t$ (straight line through the origin); at $t=4$, $v_{\max} = \SI{8}{m.s^{-1}}$. \textbf{Phase 2}: $a = \SI{-3}{m.s^{-2}}$, so $v = 8 - 3(t-4)$; the particle stops when $v=0$, i.e. $t - 4 = \dfrac{8}{3}$, so $t = \dfrac{20}{3} \approx \SI{6.67}{s}$.\\
The velocity–time graph is a triangle with vertices $(0,0)$, $(4,8)$, $(20/3, 0)$. The displacement–time graph is a parabola opening upwards for $0\le t\le 4$ ($x = t^2$), then a parabola opening downwards from $t=4$ to $t = 20/3$, with horizontal tangent at the end (since $v=0$).
\begin{center}
\begin{tikzpicture}[scale=0.55]
\draw[->] (0,0) -- (8.5,0) node[right] {$t$ (s)};
\draw[->] (0,0) -- (0,9.5) node[above] {$v$ (m/s)};
\draw[thick,popblue] (0,0) -- (4,8) -- (6.67,0);
\draw (4,0.15) -- (4,-0.15) node[below] {4};
\draw (6.67,0.15) -- (6.67,-0.15) node[below] {$\frac{20}{3}$};
\draw (0.15,8) -- (-0.15,8) node[left] {8};
\draw[dashed,gray] (4,0) -- (4,8);
\end{tikzpicture}
\end{center}
\item Total displacement = area of the triangle:
\[
s = \tfrac12 \times \dfrac{20}{3} \times 8 = \dfrac{80}{3} \approx \SI{26.7}{m}.
\]
(Phase by phase: $s_1 = \tfrac12(4)(8) = \SI{16}{m}$; $s_2 = \tfrac12\left(\tfrac83\right)(8) = \tfrac{32}{3}\ \mathrm{m}$; total $16 + \tfrac{32}{3} = \tfrac{80}{3}\ \mathrm{m}$.)
\end{enumerate}
\end{corrige}

\begin{corrige}[Exercise 9 — GCE style: Trapezium velocity–time graph]
\textbf{(a) Accelerations} (gradient of each segment):
\begin{itemize}
\item Phase 1 ($0\le t\le 2$): $a = \dfrac{6-0}{2-0} = \SI{3}{m.s^{-2}}$.
\item Phase 2 ($2\le t\le 6$): $a = \dfrac{6-6}{6-2} = \SI{0}{m.s^{-2}}$.
\item Phase 3 ($6\le t\le 10$): $a = \dfrac{0-6}{10-6} = \SI{-1.5}{m.s^{-2}}$ (retardation of $\SI{1.5}{m.s^{-2}}$).
\end{itemize}
\textbf{(b) Total distance} = area of the trapezium (parallel sides $10\ \mathrm{s}$ and $4\ \mathrm{s}$, height $6\ \mathrm{m\,s^{-1}}$):
\[
s = \tfrac12(10+4)\times 6 = \SI{42}{m}.
\]
\textbf{(c) Average velocity} $= \dfrac{s}{t} = \dfrac{42}{10} = \SI{4.2}{m.s^{-1}}$.\\[2pt]
\textbf{(d) Acceleration–time graph}: three horizontal segments: $a=3$ on $[0,2]$, $a=0$ on $[2,6]$, $a=-1.5$ on $[6,10]$.
\begin{center}
\begin{tikzpicture}[scale=0.55]
\draw[->] (0,-2) -- (11,-2) node[right] {$t$ (s)};
\draw[->] (0,-3) -- (0,4) node[above] {$a$ (m/s$^2$)};
% Phase 1: a=3
\draw[thick,popblue] (0,3) -- (2,3);
% Phase 2: a=0
\draw[thick,popblue] (2,0) -- (6,0);
% Phase 3: a=-1.5
\draw[thick,popblue] (6,-1.5) -- (10,-1.5);
% Discontinuities
\draw[dashed,gray] (2,3) -- (2,0);
\draw[dashed,gray] (6,0) -- (6,-1.5);
% Tick marks and labels
\foreach \x/\lbl in {2/2, 6/6, 10/10} {
    \draw (\x,-1.85) -- (\x,-2.15) node[below] {\lbl};
}
\draw (0.15,3) -- (-0.15,3) node[left] {3};
\draw (0.15,0) -- (-0.15,0) node[left] {0};
\draw (0.15,-1.5) -- (-0.15,-1.5) node[left] {$-1.5$};
\end{tikzpicture}
\end{center}
\textbf{Indicative mark scheme (12 marks):} (a) M1 for gradient method, A1 for $3$ and $0$, A1 for $-1.5$ (3 marks). (b) M1 area method, A1 split into three correct parts, M1 summing, A1 for $\SI{42}{m}$ (4 marks). (c) M1 formula, A1 for $\SI{4.2}{m.s^{-1}}$ (2 marks). (d) B1 horizontal segments, B1 correct values, B1 correct time intervals with discontinuities shown (3 marks). \emph{Examiner's expectations:} units must be quoted; a negative acceleration must be interpreted as deceleration only because velocity is positive; the area must be identified with displacement.
\end{corrige}

\begin{corrige}[Exercise 10 — GCE style: Two objects meeting]
Take upward as positive, origin at the ground, $g = \SI{10}{m.s^{-2}}$.
\begin{enumerate}
\item \textbf{Coconut} (dropped from rest at $x = 20$): $u = 0$, $a = -10$:
\[
x_C = 20 - 5t^2, \qquad v_C = -10t.
\]
\textbf{Stone} (projected from the ground at $\SI{15}{m.s^{-1}}$): $u = 15$, $a = -10$:
\[
x_S = 15t - 5t^2, \qquad v_S = 15 - 10t.
\]
\item They cross when $x_C = x_S$:
\[
20 - 5t^2 = 15t - 5t^2 \implies 20 = 15t \implies t = \dfrac{4}{3}\ \mathrm{s} \approx \SI{1.33}{s}.
\]
\item Height of crossing: $x_C = 20 - 5\left(\dfrac{4}{3}\right)^2 = 20 - \dfrac{80}{9} = \dfrac{100}{9} \approx \SI{11.1}{m}$ above the ground.
\item Velocities at $t = \dfrac{4}{3}$ s:
\[
v_C = -10\times\tfrac{4}{3} = -\tfrac{40}{3} \approx \SI{-13.3}{m.s^{-1}} \ (\text{coconut moving down}),
\]
\[
v_S = 15 - 10\times\tfrac{4}{3} = \tfrac{5}{3} \approx \SI{1.67}{m.s^{-1}} \ (\text{stone still moving up}).
\]
\end{enumerate}
\textbf{Indicative mark scheme (14 marks):} (a) B1 sign convention stated, B1 coconut equation, B1 stone equation, B1 both velocity equations (4 marks). (b) M1 equating positions, A1 simplification (the $t^2$ terms cancel), M1 solving, A1 $t = 4/3$ s (4 marks). (c) M1 substitution, A1 exact value, A1 decimal with unit (3 marks). (d) M1 use of $v = u + at$, A1 $v_C$, A1 $v_S$ with directions interpreted (3 marks). \emph{Examiner's expectations:} a clearly stated sign convention; the physical meaning of the signs of the velocities; exact fractions accepted.
\end{corrige}

\begin{corrige}[Exercise 11 — GCE style: Overtaking and stopping distance]
\textbf{(a) Overtaking time.} Let $t$ be the time measured from the departure of taxi A. Taxi A (constant velocity): $x_A = 15t$. Taxi B starts at $t = 5$ from rest with $a = \SI{2}{m.s^{-2}}$, so for $t \ge 5$:
\[
x_B = \tfrac12 (2)(t-5)^2 = (t-5)^2.
\]
B overtakes A when $x_B = x_A$:
\[
(t-5)^2 = 15t \implies t^2 - 10t + 25 = 15t \implies t^2 - 25t + 25 = 0.
\]
\[
t = \dfrac{25 \pm \sqrt{625 - 100}}{2} = \dfrac{25 \pm \sqrt{525}}{2} = \dfrac{25 \pm 5\sqrt{21}}{2}.
\]
Numerically $t \approx \dfrac{25 \pm 22.91}{2}$, giving $t \approx 1.04$ s (rejected, since B has not started yet, $t < 5$) or
\[
t \approx \SI{23.96}{s} \approx \SI{24.0}{s} \quad \text{after A's departure}.
\]
\textbf{(b) Distance from the park:}
\[
x_A = 15 \times 23.96 \approx \SI{359}{m}.
\]
\textbf{(c) Reaction distance.} Convert the speed: $90\ \mathrm{km\,h^{-1}} = \dfrac{90}{3.6} = \SI{25}{m.s^{-1}}$.
\[
d_{\text{reaction}} = vt = 25 \times 0.7 = \SI{17.5}{m}.
\]
\textbf{(d) Braking distance.} Using $v^2 = u^2 + 2as$ with final velocity $0$ and $a = \SI{-5}{m.s^{-2}}$:
\[
0 = 25^2 - 2(5)d \implies d = \dfrac{625}{10} = \SI{62.5}{m}.
\]
\textbf{(e) Total stopping distance:}
\[
d_{\text{total}} = 17.5 + 62.5 = \SI{80}{m} > \SI{60}{m}.
\]
\textbf{No}, the car does \emph{not} stop in time: it needs $\SI{80}{m}$ but the obstacle is only $\SI{60}{m}$ ahead, so it hits the obstacle (with a residual speed $v = \sqrt{625 - 10(60-17.5)} = \sqrt{200} \approx \SI{14.1}{m.s^{-1}}$ — optional comment).
\\[2pt]
\textbf{Indicative mark scheme (14 marks):} (a) B1 expression for $x_A$, B1 expression for $x_B$ with the time shift $(t-5)$, M1 equating, M1 solving the quadratic, A1 selecting the valid root $t \approx 24.0$ s with rejection of the other root justified (5 marks). (b) M1 substitution, A1 value, A1 unit (3 marks). (c) M1 conversion of km/h to m/s, A1 $\SI{17.5}{m}$ (2 marks). (d) M1 use of $v^2 = u^2 + 2as$, A1 $\SI{62.5}{m}$ (2 marks). (e) M1 adding both distances, A1 correct conclusion with comparison (2 marks). \emph{Examiner's expectations:} the time delay of taxi B must appear as $(t-5)$; the extraneous root must be rejected with a reason; the final answer to (e) must be an explicit justified statement, not just a number.
\end{corrige}

\newpage

\coursec{Summary sheet}

\begin{popfigure}
\begin{tikzpicture}[
  mindmap, grow cyclic, every node/.style={concept},
  root concept/.append style={concept color=popblue, font=\bfseries\small, text=white},
  level 1/.append style={level distance=3.6cm, sibling angle=60, font=\scriptsize},
  level 1 concept/.append style={minimum size=2.2cm},
  every node/.append style={scale=0.85}
]
\node[root concept]{Kinematics in a Straight Line}
  child[concept color=poporange]{ node {Position, displacement, velocity, acceleration} }
  child[concept color=popgreen]{ node {Uniform acceleration: the suvat equations} }
  child[concept color=poppurple]{ node {$x$--$t$ graph: gradient $=$ velocity} }
  child[concept color=poppink]{ node {$v$--$t$ graph: gradient $=$ $a$, area $=$ $s$} }
  child[concept color=popgold]{ node {$a$--$t$ graph: area $=$ change of velocity} }
  child[concept color=popblue!70!black]{ node {Free fall: $a=-g$, $g=9.81\ \mathrm{m\,s^{-2}}$} }
  child[concept color=popgreen!60!black]{ node {Applications: braking, lifts, falling objects} };
\end{tikzpicture}
\legende{Mind map of the chapter: kinematics of particles moving in a straight line.}
\end{popfigure}

\begin{aretenir}[Summary Sheet --- Kinematics of Particles Moving in a Straight Line]

\textbf{Key definitions}
\begin{itemize}
\item \cle{Position} $x$: location of a particle relative to a fixed origin $O$ on the line; \cle{displacement} $s$: change of position, a \emph{vector} quantity (its sign gives the direction).
\item \cle{Distance}: total length of the path travelled, a scalar (always positive or zero).
\item \cle{Velocity}: $v=\dfrac{ds}{dt}$, the rate of change of displacement; \cle{speed} $=|v|$, the magnitude of velocity.
\item \cle{Acceleration}: $a=\dfrac{dv}{dt}=\dfrac{d^{2}s}{dt^{2}}$; \cle{deceleration} (retardation) means the acceleration is in the direction opposite to the motion.
\item \cle{Uniform acceleration}: the acceleration $a$ remains constant throughout the motion.
\item \cle{Free fall}: motion under gravity alone (air resistance neglected), with $a=g\approx \SI{9.81}{m.s^{-2}}$ directed vertically downwards (take $g=\SI{10}{m.s^{-2}}$ only when the question instructs it).
\end{itemize}

\textbf{Formulas to know (uniform acceleration, ``suvat'')}
\[
v=u+at,\qquad s=ut+\tfrac12 at^{2},\qquad s=vt-\tfrac12 at^{2},
\]
\[
v^{2}=u^{2}+2as,\qquad s=\tfrac12(u+v)t,
\]
where $u$ = initial velocity, $v$ = final velocity, $a$ = acceleration, $s$ = displacement, $t$ = time. Each equation involves exactly four of the five quantities. These are valid \emph{only} when $a$ is constant.

\textbf{Graphical facts}
\begin{itemize}
\item $x$--$t$ graph: \textbf{gradient} $=$ velocity; a horizontal part means the particle is at rest; a curve bending upwards means it is speeding up.
\item $v$--$t$ graph: \textbf{gradient} $=$ acceleration; \textbf{signed area} between the graph and the time axis $=$ displacement; area below the time axis counts \emph{negative}.
\item $a$--$t$ graph: \textbf{area} between the graph and the time axis $=$ change in velocity $v-u$; for uniform acceleration the graph is a horizontal line.
\item The three graphs are linked by calculus: the gradient of the $x$--$t$ graph gives the $v$--$t$ graph; the gradient of the $v$--$t$ graph gives the $a$--$t$ graph; going the other way uses areas (integration).
\end{itemize}

\textbf{Free fall essentials}
\begin{itemize}
\item Choose a positive direction (usually upwards) and keep it throughout: then $a=-g$ for the whole motion, going up \emph{and} coming down.
\item At the highest point, $v=0$ but $a=-g$: the acceleration is never zero during free fall.
\item For a particle projected vertically upwards with speed $u$: time to the highest point $t_{\text{up}}=\dfrac{u}{g}$, maximum height $H=\dfrac{u^{2}}{2g}$; the time up equals the time down to the same level, and the particle returns with speed $u$.
\item An object dropped from rest at height $h$: $h=\tfrac12 gt^{2}$, impact speed $v=\sqrt{2gh}$.
\end{itemize}

\textbf{Methods}
\begin{enumerate}
\item Draw a diagram, mark the origin and the positive direction.
\item List the known suvat quantities and the one required; pick the equation containing exactly those four.
\item Substitute with \emph{signs}, solve, then interpret the result (a negative time, or a negative speed for a car still moving forwards, may have to be rejected).
\item For multi-stage motion, treat each stage separately: the final velocity of one stage is the initial velocity of the next.
\item For graph questions: compute gradients and areas stage by stage; use the total signed area for displacement, and the sum of the absolute areas for distance.
\end{enumerate}

\textbf{Common mistakes}
\begin{itemize}
\item Using the suvat equations when the acceleration is not constant.
\item Mixing signs: forgetting that $a=-g$ when ``up'' is positive, or giving a falling object $a=+g$ while taking upwards as positive.
\item Confusing distance with displacement on a $v$--$t$ graph (ignoring the area below the axis).
\item Thinking that $v=0$ at maximum height means $a=0$ --- false, $a=-g$ throughout the flight.
\item Reading the gradient of a $v$--$t$ graph as velocity instead of acceleration.
\item Forgetting units, or mixing km/h with m/s (divide by $3.6$ to convert km/h to m/s).
\end{itemize}

\textbf{GCE tips}
\begin{itemize}
\item Always state the equation you use before substituting --- method marks are awarded for this.
\item Use $g=9.81\ \mathrm{m\,s^{-2}}$ unless the paper says $g=10\ \mathrm{m\,s^{-2}}$; give final answers to 3 significant figures.
\item Sketch the $v$--$t$ graph even when not asked: it often reveals the quickest solution (the area method).
\item Check your answer physically: a braking car must have $a<0$ and $v\ge 0$ until it stops.
\item In ``two particles'' problems (one thrown up, one dropped from height $h$), write $s_1+s_2=h$ at the meeting instant and solve for $t$.
\end{itemize}
\end{aretenir}

\findecours

\end{document}
